% =====================================================================
%  CAPITULO 2 - CONCEITOS INICIAIS
%  Leis de Ohm, Potencia, Fontes e Leis de Kirchhoff (versao revisada)
% =====================================================================
\section{Conceitos Iniciais: Leis de Ohm, Potência, Fontes e Kirchhoff}

\begin{conceito}[Antes de começar]
\textbf{Corrente elétrica:} $I = \dfrac{Q}{t}$, medida em ampères
($\SI{1}{\ampere}=\SI{1}{\coulomb\per\second}$). Como $Q = n\,e$, o número de
elétrons é $n = Q/e$, com $e=\SI{1.6e-19}{\coulomb}$.\par
\textbf{1\textsuperscript{a} lei de Ohm:} $U = R\,I$.\par
\textbf{2\textsuperscript{a} lei de Ohm (resistividade):}
$R = \rho\,\dfrac{L}{A}$, onde $\rho$ é a resistividade do material, $L$ o
comprimento e $A$ a área da seção.\par
\textbf{Potência elétrica:}
$P = U\,I = R\,I^2 = \dfrac{U^2}{R}$.\par
\textbf{Fonte real:} $U = \EMF - r\,I$, onde $r$ é a resistência interna.\par
\textbf{Leis de Kirchhoff:} \emph{dos nós} --- a soma das correntes que entram é
igual à soma das que saem; \emph{das malhas} --- a soma algébrica das tensões ao
longo de uma malha fechada é nula.
\end{conceito}

\legendaniveis

% =====================================================================
\subsection{Corrente elétrica e carga}
% =====================================================================

\blocofixacao

\begin{questao}[1]{}
Se uma corrente de $\SI{3}{\ampere}$ percorre um fio durante $\SI{2}{\minute}$, a
carga que atravessa sua seção reta, em coulombs, vale:
\begin{alternativas}[3]
\item 60
\item 110
\item 360
\item 220
\item 180
\end{alternativas}
\resposta{(c)}
\resolucao{Convertendo o tempo para segundos, $t = 2\cdot 60 = \SI{120}{\second}$:
\[
Q = I\,t = 3\cdot 120 = \SI{360}{\coulomb}.
\]
O erro típico (alternativa e) é usar $t = \SI{2}{\minute}$ diretamente, sem
converter para segundos.}
\end{questao}

\begin{questao}[1]{}
Uma carga de $\SI{120}{\coulomb}$ passa uniformemente pela seção de um fio durante
$\SI{1}{\minute}$. A intensidade da corrente é:
\begin{alternativas}[3]
\item $\tfrac{1}{30}\,\si{\ampere}$
\item $\tfrac{1}{2}\,\si{\ampere}$
\item \SI{2}{\ampere}
\item \SI{30}{\ampere}
\item \SI{120}{\ampere}
\end{alternativas}
\resposta{(c)}
\resolucao{$I = \dfrac{Q}{t} = \dfrac{120}{60} = \SI{2}{\ampere}$.}
\end{questao}

\begin{questao}[1]{}
Uma lâmpada permanece acesa durante $\SI{5}{\minute}$ sob corrente de
$\SI{2}{\ampere}$. A carga total que circula, em coulombs, é:
\begin{alternativas}[3]
\item 0,4
\item 2,5
\item 10
\item 150
\item 600
\end{alternativas}
\resposta{(e)}
\resolucao{$Q = I\,t = 2\cdot(5\cdot60) = 2\cdot 300 = \SI{600}{\coulomb}$.}
\end{questao}

\blocoaplicacao

\begin{questao}[2]{}
Um condutor metálico é percorrido por uma corrente contínua e constante de
$\SI{32}{\milli\ampere}$. Determine:
\begin{itensq}
\item a carga que atravessa uma seção reta por segundo;
\item o número de elétrons correspondente.
\end{itensq}
\dados{$e = \SI{1.6e-19}{\coulomb}$.}
\resposta{(a) \SI{32}{\milli\coulomb}; (b) $n = \pot{2.0}{17}$ elétrons}
\resolucao{(a) Em $t = \SI{1}{\second}$:
$Q = I\,t = \pot{32}{-3}\cdot 1 = \SI{32}{\milli\coulomb}$.\\
(b) $n = \dfrac{Q}{e} = \dfrac{\pot{32}{-3}}{\pot{1.6}{-19}} = \pot{2.0}{17}$
elétrons por segundo.}
\end{questao}

\begin{questao}[2]{}
O gráfico mostra a intensidade da corrente $i$ em um condutor em função do tempo.
Determine a carga que atravessa uma seção do condutor nos cinco primeiros
segundos.

\circuitoqq{%
\begin{tikzpicture}
 \begin{axis}[width=9cm,height=5cm,xlabel={$t$ (s)},ylabel={$i$ (A)},
   xmin=0,xmax=5.5,ymin=0,ymax=5,xtick={0,1,2,3,4,5},ytick={0,1,2,3,4},
   grid=both,axis lines=left]
   \addplot[Icol,very thick] coordinates {(0,4)(5,4)};
   \addplot[Icol!30] coordinates {(0,0)(0,4)};
 \end{axis}
\end{tikzpicture}}{Corrente constante em função do tempo.}

\resposta{$Q = \SI{20}{\coulomb}$}
\resolucao{A carga é a \emph{área} sob o gráfico $i \times t$. Como a corrente é
constante e igual a $\SI{4}{\ampere}$ durante $\SI{5}{\second}$, a área é um
retângulo:
\[
Q = i\,t = 4\cdot 5 = \SI{20}{\coulomb}.
\]
Essa interpretação por área vale mesmo quando a corrente varia --- basta calcular
a área da figura correspondente.}
\end{questao}

\begin{questao}[2]{}
O gráfico representa a corrente elétrica $i$ que atravessa um condutor em função
do tempo. Determine a carga total que passa por uma seção reta entre $t=0$ e
$t=\SI{6}{\second}$.

\circuitoqq{%
\begin{tikzpicture}
 \begin{axis}[width=9.5cm,height=5cm,xlabel={$t$ (s)},ylabel={$i$ (A)},
   xmin=0,xmax=6.5,ymin=0,ymax=5,xtick={0,1,2,3,4,5,6},ytick={0,2,4},
   grid=both,axis lines=left]
   \addplot[Icol,very thick] coordinates {(0,4)(2,4)(6,0)};
   \addplot[Icol!12] coordinates {(0,0)(0,4)};
 \end{axis}
\end{tikzpicture}}{Corrente variável em função do tempo.}

\resposta{$Q = \SI{16}{\coulomb}$}
\resolucao{A carga é a \emph{área} sob o gráfico $i \times t$. A figura é composta
de um retângulo (de $0$ a \SI{2}{\second}) e um triângulo (de $2$ a \SI{6}{\second}):
\[
Q = \underbrace{4\cdot2}_{\text{retângulo}}
  + \underbrace{\tfrac12\cdot4\cdot4}_{\text{triângulo}}
  = 8 + 8 = \SI{16}{\coulomb}.
\]
Mesmo com corrente variável, a área continua valendo como carga total.}
\end{questao}

\begin{questao}[2]{}
Qual é a carga elétrica que atravessa, durante $\SI{10}{\hour}$, a seção de um
condutor percorrido por corrente constante de $\SI{5}{\ampere}$?
\begin{alternativas}[3]
\item \SI{50}{\coulomb}
\item \SI{1000}{\coulomb}
\item \SI{90000}{\coulomb}
\item \SI{180000}{\coulomb}
\item \SI{360000}{\coulomb}
\end{alternativas}
\resposta{(d)}
\resolucao{Convertendo o tempo: $\SI{10}{\hour} = 10\cdot3600 = \SI{36000}{\second}$.
\[
Q = I\,t = 5\cdot36000 = \SI{180000}{\coulomb}.
\]
O erro comum é esquecer de converter horas em segundos.}
\end{questao}

\blocoaprofundamento

\begin{questao}[3]{}
O gráfico mostra a corrente em um condutor variando linearmente de
$\SI{6}{\ampere}$ a $\SI{2}{\ampere}$ entre $t=0$ e $t=\SI{8}{\second}$.

\circuitoqq{%
\begin{tikzpicture}
 \begin{axis}[width=9cm,height=4.6cm,xlabel={$t$ (s)},ylabel={$i$ (A)},
   xmin=0,xmax=9,ymin=0,ymax=7,xtick={0,2,4,6,8},ytick={0,2,4,6},
   grid=both,axis lines=left]
   \addplot[Icol,very thick] coordinates {(0,6)(8,2)};
 \end{axis}
\end{tikzpicture}}{Corrente decrescente em função do tempo.}

\begin{itensq}
\item Determine a carga total que atravessa o condutor nesse intervalo.
\item Calcule a corrente \emph{média} e verifique a coerência com o item (a).
\item Em que instante a corrente instantânea iguala a média?
\end{itensq}
\resposta{(a) $Q = \SI{32}{\coulomb}$; (b) $i_m = \SI{4}{\ampere}$;
(c) $t = \SI{4}{\second}$}
\resolucao{(a) A carga é a área sob o gráfico --- aqui, um trapézio de bases
\SI{6}{} e \SI{2}{\ampere} e altura \SI{8}{\second}:
\[
Q=\frac{(6+2)}{2}\cdot 8 = \SI{32}{\coulomb}.
\]
(b) $i_m = \dfrac{Q}{\Delta t} = \dfrac{32}{8} = \SI{4}{\ampere}$, que coincide
com a média aritmética $\dfrac{6+2}{2}$ --- válido \emph{porque} a variação é
linear.\\
(c) Como a reta decresce uniformemente, ela cruza o valor médio exatamente no
\emph{meio} do intervalo: $t = \SI{4}{\second}$.
Conferindo pela equação da reta $i(t) = 6 - 0{,}5t$:
$i(4) = 6-2 = \SI{4}{\ampere}$. \checkmark\\
\textbf{Cuidado:} se a corrente variasse de forma \emph{não} linear (por exemplo,
exponencialmente), a corrente média \textbf{não} seria a média dos extremos ---
seria preciso calcular a área real sob a curva.}
\end{questao}

\begin{questao}[3]{}
Um fio de cobre de seção $A = \SI{1}{\milli\meter\squared}$ conduz corrente
constante de $\SI{10}{\ampere}$. O cobre possui aproximadamente
$n = \SI{8.5e28}{}$ elétrons livres por metro cúbico.
\begin{itensq}
\item Determine a \emph{velocidade de deriva} dos elétrons.
\item Quanto tempo um elétron leva para percorrer $\SI{1}{\meter}$ de fio?
\item Como conciliar esse resultado com o fato de uma lâmpada acender
      \emph{instantaneamente} ao acionarmos o interruptor?
\end{itensq}
\dados{$e = \SI{1.6e-19}{\coulomb}$.}
\resposta{(a) $v_d \approx \SI{0.74}{\milli\meter\per\second}$;
(b) $\approx\SI{1360}{\second}$ ($\approx\SI{23}{\minute}$); (c) ver comentário}
\resolucao{\textbf{(a)} A corrente relaciona-se com a velocidade de deriva por
$I = n\,e\,A\,v_d$. Isolando:
\[
v_d=\frac{I}{n\,e\,A}
=\frac{10}{(\pot{8.5}{28})(\pot{1.6}{-19})(\pot{1}{-6})}
\approx\pot{7.4}{-4}\,\si{\meter\per\second}=\SI{0.74}{\milli\meter\per\second}.
\]
\textbf{(b)} $t = \dfrac{1}{\pot{7.4}{-4}} \approx \SI{1360}{\second}
\approx \SI{23}{\minute}$ --- um elétron leva aproximadamente 23 minutos para percorrer um metro.

\textbf{(c) A resolução do aparente paradoxo.} A lâmpada acende
instantaneamente porque \textbf{não é preciso que um elétron viaje da chave até a
lâmpada}. O que se propaga rapidamente é o \emph{campo elétrico} ao longo do
condutor, a uma velocidade próxima à da luz ($\sim\SI{2e8}{\meter\per\second}$
em cabos reais). Esse campo põe em movimento, quase simultaneamente,
\emph{todos} os elétrons livres ao longo de todo o fio --- inclusive os que já
estavam dentro do filamento da lâmpada.\\
\textbf{Analogia:} é como um tubo já cheio de água. Ao empurrar na entrada, sai
água na saída imediatamente, embora cada molécula individual se desloque
lentamente. A informação (pressão) viaja rápido; a matéria, devagar.\\
\textbf{Ordens de grandeza envolvidas:} deriva $\sim\SI{e-4}{\meter\per\second}$;
agitação térmica dos elétrons $\sim\SI{e6}{\meter\per\second}$; sinal elétrico
$\sim\SI{e8}{\meter\per\second}$ --- doze ordens de grandeza separando a mais
lenta da mais rápida.}
\end{questao}

% BEGIN BANCO COMPLEMENTAR Corrente elétrica e carga
% =====================================================================
%  Banco complementar --- Corrente Eletrica e Carga  (REVISADO)
%  Substitui a versao gerada por template (8 clones em N2, 10 em N3 e
%  8 questoes conceituais marcadas como N4).
%  O cap02 ja cobre: Q=It em tres questoes N1, condutor com 32 mA,
%  DOIS graficos i x t, carga em 10 h, rampa linear de 6 A a 2 A e a
%  velocidade de deriva em fio de cobre de 1 mm2 com 10 A.
%  Este banco explora: capacidade em A.h e autonomia sob corrente
%  variavel, formas de onda (triangular, trapezoidal, pulsada),
%  carga liquida x carga total em correntes bipolares, densidade de
%  corrente e limite termico, eletrolise, corrente de fuga, o erro de
%  estimar carga pelo pico e a distincao entre velocidade de deriva e
%  velocidade de propagacao do sinal.
%  Distribuicao mantida: 5 N1 + 8 N2 + 10 N3 + 8 N4 = 31
% =====================================================================

\subsubsection*{Banco complementar --- Corrente elétrica e carga}

\begin{atencao}[Organização do banco]
A relação $Q=It$ vale apenas para corrente \textbf{constante}. No caso geral,
$Q=\int i\,\mathrm{d}t$ --- ou seja, a carga é a \textbf{área} sob a curva
$i\times t$, com sinal. O N3 explora formas de onda, eletrólise e limites
térmicos; o N4 trata de autonomia de baterias, erros de estimativa por pico,
correntes bipolares e a diferença entre velocidade de deriva e velocidade do
sinal.
\end{atencao}

\blocofixacao

\begin{questao}[1]{}
Uma corrente constante de $\SI{1.5}{\ampere}$ circula durante
$\SI{40}{\second}$. Determine a carga transportada.
\resposta{$Q=\SI{60}{\coulomb}$}
\resolucao{
\[
Q=It=1{,}5\times40=\SI{60}{\coulomb}.
\]
\textbf{Fundamento:} a corrente é a taxa de passagem de carga,
$i=\mathrm{d}q/\mathrm{d}t$. Sendo constante, a integral vira produto simples.
Um ampère é, por definição, um coulomb por segundo.}
\end{questao}

\begin{questao}[1]{}
Uma carga de $\SI{90}{\coulomb}$ atravessa a seção de um condutor em
$\SI{45}{\second}$. Determine a corrente média.
\resposta{$I=\SI{2}{\ampere}$}
\resolucao{
\[
I=\frac{Q}{t}=\frac{90}{45}=\SI{2}{\ampere}.
\]
\textbf{Atenção ao termo "média":} o resultado é a corrente \emph{média} no
intervalo. Ela pode ter variado --- inclusive muito --- sem que isso altere a
carga total. Duas formas de onda completamente diferentes podem ter a mesma
média.}
\end{questao}

\begin{questao}[1]{}
Determine o tempo necessário para transportar $\SI{60}{\coulomb}$ com corrente
de $\SI{250}{\milli\ampere}$.
\resposta{$t=\SI{240}{\second}=\SI{4}{\minute}$}
\resolucao{
\[
t=\frac{Q}{I}=\frac{60}{0{,}250}=\SI{240}{\second}.
\]
\textbf{Erro comum:} usar $I=\SI{250}{\ampere}$ e obter
$\SI{0.24}{\second}$. Converta o prefixo \emph{antes} de dividir --- o
resultado errado é mil vezes menor e ainda assim parece plausível.}
\end{questao}

\begin{questao}[1]{}
Determine o número de elétrons que atravessa uma seção quando circula
$\SI{16}{\coulomb}$.
\dados{$e=\pot{1.6}{-19}\,\si{\coulomb}$}
\resposta{$n=\pot{1}{20}$ elétrons}
\resolucao{
\[
n=\frac{Q}{e}=\frac{16}{\pot{1.6}{-19}}=\pot{1}{20}.
\]
Cem quintilhões de elétrons --- e essa carga passa em apenas
$\SI{8}{\second}$ com corrente de $\SI{2}{\ampere}$.\\
\textbf{Ordem de grandeza:} o coulomb é uma unidade \emph{enorme} em termos de
partículas. É por isso que a quantização da carga é imperceptível em circuitos:
um erro de um elétron em $10^{20}$ não se mede.}
\end{questao}

\begin{questao}[1]{}
Sobre o sentido convencional da corrente, é correto afirmar:
\begin{alternativas}[2]
\item coincide com o movimento dos elétrons em um metal
\item é oposto ao movimento dos elétrons em um metal
\item existe apenas em semicondutores
\item depende da intensidade da corrente
\end{alternativas}
\resposta{(b)}
\resolucao{O sentido \textbf{convencional} é o do movimento de cargas
\emph{positivas} --- escolha feita antes de se conhecer o elétron. Em metais,
os portadores são elétrons (negativos), que se movem no sentido
\textbf{oposto}.\\
\textbf{Por que a convenção não atrapalha:} uma carga $-q$ movendo-se para a
esquerda é eletricamente equivalente a $+q$ movendo-se para a direita. Todas as
leis de circuitos funcionam identicamente.\\
\textbf{Onde a distinção importa:} em efeitos que dependem do \emph{portador
real}, como o efeito Hall --- cujo sinal revela se os portadores são elétrons ou
lacunas, e é assim que se distingue um semicondutor tipo N de um tipo P.}
\end{questao}

\blocoaplicacao

\begin{questao}[2]{}
Por um condutor passam $\SI{240}{\coulomb}$ em $\SI{120}{\second}$.
\begin{itensq}
\item Determine a corrente média.
\item Determine o número de elétrons.
\end{itensq}
\resposta{$I=\SI{2}{\ampere}$; $n=\pot{1.5}{21}$}
\resolucao{
\[
I=\frac{240}{120}=\SI{2}{\ampere};
\qquad
n=\frac{240}{\pot{1.6}{-19}}=\pot{1.5}{21}\ \text{elétrons}.
\]
\textbf{Taxa por segundo:} $\pot{1.25}{19}$ elétrons atravessam a seção a cada
segundo. Se fosse possível contá-los um por vez, a $\SI{1}{\per\second}$,
levaria cerca de $400$ bilhões de anos --- trinta vezes a idade do universo.}
\end{questao}

\begin{questao}[2]{}
Uma bateria de celular tem capacidade de $\SI{2000}{\milli\ampere\hour}$.
\begin{itensq}
\item Converta para coulombs.
\item Determine por quanto tempo ela alimenta uma carga de
      $\SI{250}{\milli\ampere}$.
\end{itensq}
\resposta{$Q=\SI{7200}{\coulomb}$; $t=\SI{8}{\hour}$}
\resolucao{(a)
\[
Q=2000\,\si{\milli\ampere\hour}=2\,\si{\ampere}\times3600\,\si{\second}
=\SI{7200}{\coulomb}.
\]
(b)
\[
t=\frac{2000\,\si{\milli\ampere\hour}}{250\,\si{\milli\ampere}}=\SI{8}{\hour}.
\]
\textbf{Por que se usa $\si{\ampere\hour}$:} a conta acima sai em uma linha,
enquanto com coulombs seria preciso converter duas vezes. A unidade é prática,
mas é preciso lembrar que ela mede \emph{carga}, não energia.\\
\textbf{Ressalva:} a capacidade nominal pressupõe uma corrente de descarga
específica. Descargas muito rápidas entregam menos carga total --- efeito
Peukert.}
\end{questao}

\begin{questao}[2]{}
Uma corrente pulsada vale $\SI{5}{\ampere}$ durante
$\SI{20}{\milli\second}$ e zero durante os $\SI{80}{\milli\second}$
seguintes, repetindo-se.
\begin{itensq}
\item Determine a carga por ciclo.
\item Determine a corrente média.
\end{itensq}
\resposta{$Q=\SI{100}{\milli\coulomb}$; $I_{méd}=\SI{1}{\ampere}$}
\resolucao{(a) $Q=5\times0{,}020=\SI{0.1}{\coulomb}$.\\
(b) O período é $\SI{100}{\milli\second}$:
\[
I_{méd}=\frac{Q}{T}=\frac{0{,}1}{0{,}1}=\SI{1}{\ampere},
\]
ou, equivalentemente, $I_{pico}\times D=5(0{,}20)$.\\
\textbf{Advertência que vale repetir:} a corrente média de
$\SI{1}{\ampere}$ descreve corretamente a \emph{carga} transportada, mas
\emph{não} o aquecimento. Para efeito Joule, o que vale é o valor eficaz,
$I_{rms}=5\sqrt{0{,}2}=\SI{2.24}{\ampere}$ --- que dissipa cinco vezes mais que
$\SI{1}{\ampere}$ contínuo.}
\end{questao}

\begin{questao}[2]{}
Um condutor de seção $\SI{1.5}{\milli\meter\squared}$ conduz
$\SI{15}{\ampere}$. Determine a densidade de corrente.
\resposta{$J=\SI{10}{\ampere\per\milli\meter\squared}=\pot{1}{7}\,
\si{\ampere\per\meter\squared}$}
\resolucao{
\[
J=\frac{I}{A}=\frac{15}{1{,}5}=\SI{10}{\ampere\per\milli\meter\squared}
=\frac{15}{\pot{1.5}{-6}}=\pot{1}{7}\,\si{\ampere\per\meter\squared}.
\]
\textbf{Por que $J$ importa mais que $I$:} o aquecimento e a durabilidade de um
condutor dependem da \emph{concentração} de corrente, não do total. Valores
usuais em instalações prediais ficam entre $4$ e
$\SI{6}{\ampere\per\milli\meter\squared}$ --- este condutor está bem acima, e
aqueceria demais.\\
\textbf{Comparação:} trilhas de circuito impresso admitem
$\SI{20}{\ampere\per\milli\meter\squared}$ ou mais, porque dissipam calor
diretamente para o ar e para a placa.}
\end{questao}

\begin{questao}[2]{}
Um capacitor de $\SI{100}{\micro\farad}$, inicialmente descarregado, recebe
corrente constante de $\SI{2}{\milli\ampere}$ durante $\SI{5}{\second}$.
\begin{itensq}
\item Determine a carga acumulada.
\item Determine a tensão final.
\end{itensq}
\resposta{$Q=\SI{10}{\milli\coulomb}$; $U=\SI{100}{\volt}$}
\resolucao{(a) $Q=It=\pot{2}{-3}\times5=\SI{10}{\milli\coulomb}$.\\
(b) $U=\dfrac{Q}{C}=\dfrac{\pot{1}{-2}}{\pot{1}{-4}}=\SI{100}{\volt}$.\\
\textbf{Observação:} com corrente \emph{constante}, a tensão cresce
\textbf{linearmente} no tempo --- uma rampa de
$\SI{20}{\volt\per\second}$. É exatamente o oposto do carregamento por fonte de
tensão através de resistor, que produz crescimento exponencial.\\
\textbf{Aplicação:} é assim que se geram rampas precisas em geradores de sinal e
em conversores analógico-digitais de dupla rampa.}
\end{questao}

\begin{questao}[2]{}
Em uma solução de $\mathrm{NaCl}$, íons $\mathrm{Na^{+}}$ movem-se para a
direita transportando $\SI{3}{\coulomb}$ por segundo, enquanto
$\mathrm{Cl^{-}}$ movem-se para a esquerda transportando
$\SI{2}{\coulomb}$ por segundo em módulo. Determine a corrente resultante.
\resposta{$I=\SI{5}{\ampere}$ para a direita}
\resolucao{Cargas positivas indo para a direita e cargas negativas indo para a
esquerda produzem corrente convencional \textbf{no mesmo sentido} --- elas se
\emph{somam}:
\[
I=3+2=\SI{5}{\ampere}\ \text{para a direita}.
\]
\textbf{Erro comum:} subtrair, obtendo $\SI{1}{\ampere}$. O sinal negativo do
íon \emph{e} o sentido oposto do movimento se cancelam, resultando em
contribuição positiva.\\
\textbf{Regra:} corrente convencional $=$ (cargas positivas no sentido adotado)
$+$ (cargas negativas no sentido contrário). Em eletrólitos e em plasmas, os
dois tipos de portador contribuem simultaneamente.}
\end{questao}

\begin{questao}[2]{}
Uma bateria descarregada de $\SI{1200}{\milli\ampere\hour}$ é recarregada com
corrente constante de $\SI{400}{\milli\ampere}$.
\begin{itensq}
\item Determine o tempo teórico de recarga.
\item Explique por que o tempo real é maior.
\end{itensq}
\resposta{$t=\SI{3}{\hour}$ teórico; na prática, $20$ a $40\%$ mais}
\resolucao{(a)
\[
t=\frac{1200}{400}=\SI{3}{\hour}.
\]
(b) \textbf{Três razões para o tempo real ser maior.}
\begin{itemize}
\item \textbf{Rendimento de carga:} parte da carga injetada não fica armazenada
--- ela se perde em reações secundárias e em aquecimento. O rendimento
coulômbico típico fica entre $85$ e $95\%$.
\item \textbf{Fase de tensão constante:} carregadores modernos aplicam corrente
constante até certa tensão e depois \emph{reduzem} a corrente
progressivamente. Essa fase final é lenta e pode responder por metade do tempo
total.
\item \textbf{Limitação térmica:} se a bateria aquecer, o carregador reduz a
corrente para protegê-la.
\end{itemize}
\textbf{Estimativa prática:} multiplique o tempo teórico por $1{,}3$ a
$1{,}4$.}
\end{questao}

\begin{questao}[2]{}
Uma corrente de $\SI{8}{\ampere}$ atravessa um fusível cuja especificação é
$I^{2}t=\SI{50}{\ampere\squared\second}$. Determine o tempo até a fusão.
\resposta{$t\approx\SI{0.78}{\second}$}
\resolucao{
\[
t=\frac{I^{2}t}{I^{2}}=\frac{50}{64}=\SI{0.78}{\second}.
\]
\textbf{Por que a especificação usa $I^{2}t$:} a energia dissipada no elemento
fusível é $RI^{2}t$, e a fusão ocorre quando essa energia atinge o valor
necessário para fundir o metal. Como $R$ é fixo, o parâmetro característico é
$I^{2}t$.\\
\textbf{Consequência:} dobrar a corrente reduz o tempo de atuação por
\textbf{quatro}. Com $\SI{16}{\ampere}$, o mesmo fusível abriria em
$\SI{0.20}{\second}$; com $\SI{80}{\ampere}$, em
$\SI{7.8}{\milli\second}$.\\
\textbf{Note a diferença conceitual:} aqui a grandeza relevante é
$I^{2}t$ (energia), e não $It$ (carga). São perguntas distintas sobre a mesma
corrente.}
\end{questao}

\blocoaprofundamento

\begin{questao}[3]{}
A corrente em um circuito varia segundo $i(t)=1+2t$ (ampères, com $t$ em
segundos), no intervalo $0\le t\le\SI{2}{\second}$.
\begin{itensq}
\item Determine a carga total transferida.
\item Determine a corrente média e compare com a média dos extremos.
\end{itensq}
\resposta{$Q=\SI{6}{\coulomb}$; $I_{méd}=\SI{3}{\ampere}$}
\resolucao{(a)
\[
Q=\int_0^{2}(1+2t)\,\mathrm{d}t=\Big[t+t^{2}\Big]_0^{2}=2+4=\SI{6}{\coulomb}.
\]
(b) $I_{méd}=Q/T=6/2=\SI{3}{\ampere}$.\\
\textbf{Média dos extremos:} $i(0)=1$ e $i(2)=\SI{5}{\ampere}$, cuja média é
$\SI{3}{\ampere}$ --- \emph{coincide}.\\
\textbf{Mas cuidado:} a coincidência ocorre porque a função é \textbf{linear}.
Para $i(t)=t^{2}$ no mesmo intervalo, a média verdadeira seria
$\frac{1}{2}\int_0^{2}t^{2}\mathrm{d}t=\SI{1.33}{\ampere}$, enquanto a média
dos extremos daria $\SI{2}{\ampere}$ --- erro de $50\%$.\\
\textbf{Regra:} a média dos extremos só vale para variação linear. No caso
geral, é preciso integrar.}
\end{questao}

\begin{questao}[3]{}
Um pulso de corrente cresce linearmente de zero a $\SI{6}{\ampere}$ em
$\SI{2}{\milli\second}$ e retorna a zero em mais $\SI{3}{\milli\second}$.
\begin{itensq}
\item Determine a carga transferida.
\item Determine a corrente média no pulso.
\end{itensq}
\resposta{$Q=\SI{15}{\milli\coulomb}$; $I_{méd}=\SI{3}{\ampere}$}
\resolucao{(a) A carga é a \textbf{área} do triângulo:
\[
Q=\frac{1}{2}\,I_{pico}\,T=\frac{1}{2}(6)(\pot{5}{-3})
=\pot{1.5}{-2}\,\si{\coulomb}=\SI{15}{\milli\coulomb}.
\]
\textbf{Observe:} a repartição $2+3$ do tempo \emph{não} importa para a área ---
qualquer triângulo de base $\SI{5}{\milli\second}$ e altura
$\SI{6}{\ampere}$ tem a mesma área. A assimetria afetaria o valor eficaz, não a
carga.\\
(b) $I_{méd}=Q/T=0{,}015/0{,}005=\SI{3}{\ampere}$, exatamente metade do pico
--- característica de qualquer forma triangular.}
\end{questao}

\begin{questao}[3]{}
Uma corrente trapezoidal sobe de $0$ a $\SI{4}{\ampere}$ em
$\SI{1}{\second}$, permanece constante por $\SI{3}{\second}$ e cai a zero em
$\SI{2}{\second}$.
\begin{itensq}
\item Determine a carga total.
\item Determine a corrente média.
\end{itensq}
\resposta{$Q=\SI{18}{\coulomb}$; $I_{méd}=\SI{3}{\ampere}$}
\resolucao{(a) Some as três áreas:
\[
Q=\underbrace{\tfrac12(1)(4)}_{2}+\underbrace{(3)(4)}_{12}
+\underbrace{\tfrac12(2)(4)}_{4}=\SI{18}{\coulomb}.
\]
(b) O intervalo total é $\SI{6}{\second}$:
$I_{méd}=18/6=\SI{3}{\ampere}$.\\
\textbf{Atalho para trapézios:} a área vale
$I_{pico}\times\dfrac{B+b}{2}$, com $B=\SI{6}{\second}$ (base maior) e
$b=\SI{3}{\second}$ (base menor):
$4\times4{,}5=\SI{18}{\coulomb}$ \checkmark.}
\end{questao}

\begin{questao}[3]{}
Uma corrente alterna entre $+\SI{4}{\ampere}$ por $\SI{3}{\second}$ e
$-\SI{6}{\ampere}$ por $\SI{2}{\second}$, repetindo-se.
\begin{itensq}
\item Determine a carga líquida por ciclo.
\item Determine a carga total \emph{transportada} em módulo.
\item Interprete a diferença.
\end{itensq}
\resposta{Líquida: $\SI{0}{\coulomb}$; total: $\SI{24}{\coulomb}$}
\resolucao{(a)
\[
Q_{líq}=4(3)+(-6)(2)=12-12=\SI{0}{\coulomb}.
\]
(b) Em módulo, passaram $12+12=\SI{24}{\coulomb}$.\\
(c) \textbf{As duas respostas são corretas para perguntas diferentes.}
\begin{itemize}
\item A \textbf{carga líquida} nula significa que nada se acumula: um capacitor
nesse circuito voltaria à tensão inicial a cada ciclo, e uma bateria não se
carregaria nem descarregaria.
\item A \textbf{carga transportada} de $\SI{24}{\coulomb}$ é o que
efetivamente atravessou a seção. Em eletrólise, ela é que produziria depósito
--- se o processo fosse irreversível.
\item O \textbf{aquecimento} depende de outra grandeza ainda:
$I_{rms}=\sqrt{\frac{4^{2}(3)+6^{2}(2)}{5}}=\SI{4.9}{\ampere}$.
\end{itemize}
\textbf{Moral:} "quanta corrente passou" é uma pergunta ambígua. Especifique se
o interesse é acúmulo, transporte ou dissipação.}
\end{questao}

\begin{questao}[3]{}
Uma corrente de descarga vale $i(t)=I_0e^{-t/\tau}$, com
$I_0=\SI{2}{\ampere}$ e $\tau=\SI{50}{\milli\second}$.
\begin{itensq}
\item Determine a carga total transferida.
\item Determine a fração transferida no primeiro $\tau$.
\end{itensq}
\resposta{$Q=\SI{0.1}{\coulomb}$; $63{,}2\%$}
\resolucao{(a)
\[
Q=\int_0^{\infty}I_0e^{-t/\tau}\,\mathrm{d}t
=I_0\tau=2(0{,}050)=\SI{0.1}{\coulomb}.
\]
\textbf{Interpretação geométrica:} a carga equivale à corrente inicial mantida
por exatamente $\tau$ --- o retângulo $I_0\times\tau$ tem a mesma área da
exponencial completa.\\
(b)
\[
\frac{Q(\tau)}{Q}=1-e^{-1}=63{,}2\%.
\]
\textbf{Compare com a energia:} na mesma descarga, a energia decai com
constante $\tau/2$ (pois $p\propto i^{2}$), e $86{,}5\%$ dela é dissipada em um
$\tau$. \textbf{Carga e energia têm ritmos diferentes} --- a energia se esgota
mais depressa.}
\end{questao}

\begin{questao}[3]{}
Uma corrente de $\SI{2}{\ampere}$ atravessa uma célula eletrolítica de cobre
durante $\SI{30}{\minute}$.
\begin{itensq}
\item Determine a carga transferida.
\item Determine a massa de cobre depositada.
\end{itensq}
\dados{$M_{Cu}=\SI{63.5}{\gram\per\mole}$; valência $z=2$;
$F=\SI{96500}{\coulomb\per\mole}$}
\resposta{$Q=\SI{3600}{\coulomb}$; $m=\SI{1.18}{\gram}$}
\resolucao{(a) $Q=2\times1800=\SI{3600}{\coulomb}$.\\
(b) Pela lei de Faraday da eletrólise:
\[
m=\frac{QM}{zF}=\frac{3600(63{,}5)}{2(96500)}
=\frac{228600}{193000}=\SI{1.18}{\gram}.
\]
\textbf{Significado de $z=2$:} cada íon $\mathrm{Cu^{2+}}$ precisa de
\emph{dois} elétrons para se neutralizar. São necessários, portanto,
$2F$ coulombs por mol depositado.\\
\textbf{Conexão conceitual:} esta é a medida mais direta de que a corrente
transporta carga em quantidades proporcionais à matéria depositada --- foi
justamente por experimentos assim que Faraday inferiu a existência de uma
unidade elementar de carga, décadas antes da descoberta do elétron.}
\end{questao}

\begin{questao}[3]{}
Um condutor de $\SI{2.5}{\milli\meter\squared}$ deve conduzir
$\SI{20}{\ampere}$.
\begin{itensq}
\item Determine a densidade de corrente.
\item Compare com o limite usual de $\SI{5}{\ampere\per\milli\meter\squared}$ e
      determine a seção adequada.
\end{itensq}
\resposta{$J=\SI{8}{\ampere\per\milli\meter\squared}$ --- excessiva; adotar
$\SI{4}{\milli\meter\squared}$ ou mais}
\resolucao{(a) $J=20/2{,}5=\SI{8}{\ampere\per\milli\meter\squared}$.\\
(b) O valor está $60\%$ acima do limite usual. A seção mínima seria
\[
A=\frac{20}{5}=\SI{4}{\milli\meter\squared},
\]
e a bitola comercial adequada é $\SI{4}{\milli\meter\squared}$ (ou
$\SI{6}{\milli\meter\squared}$, com folga).\\
\textbf{O que acontece se mantivermos $\SI{2.5}{\milli\meter\squared}$:} a
potência dissipada por metro é $\rho J^{2}A=\rho I^{2}/A$, que com a seção
menor é $1{,}6$ vezes maior. A temperatura do condutor sobe, a isolação
envelhece mais depressa e a resistência aumenta --- realimentando o
aquecimento.\\
\textbf{Ressalva:} o limite de $\SI{5}{\ampere\per\milli\meter\squared}$ é
apenas uma regra de bolso. O valor normativo depende do tipo de isolação, do
método de instalação (eletroduto, bandeja, ao ar livre), do agrupamento de
circuitos e da temperatura ambiente.}
\end{questao}

\begin{questao}[3]{}
Um capacitor de $\SI{470}{\micro\farad}$ carregado a $\SI{100}{\volt}$
apresenta corrente de fuga aproximadamente constante de
$\SI{20}{\micro\ampere}$.
\begin{itensq}
\item Determine a carga armazenada.
\item Estime o tempo para descarga completa.
\item Comente a validade da hipótese de corrente constante.
\end{itensq}
\resposta{$Q=\SI{47}{\milli\coulomb}$; $t\approx\SI{39}{\minute}$}
\resolucao{(a) $Q=CU=\pot{470}{-6}(100)=\SI{47}{\milli\coulomb}$.\\
(b) Com $I$ constante:
\[
t=\frac{Q}{I}=\frac{\pot{4.7}{-2}}{\pot{2}{-5}}=\SI{2350}{\second}
\approx\SI{39}{\minute}.
\]
(c) \textbf{A hipótese é otimista.} A corrente de fuga é aproximadamente
proporcional à tensão ($I=U/R_{fuga}$), de modo que ela \emph{diminui} durante
a descarga --- e o processo é exponencial, não linear.\\
Com $R_{fuga}=100/\pot{2}{-5}=\SI{5}{\mega\ohm}$, tem-se
$\tau=R_{fuga}C=\SI{2350}{\second}$ --- o \emph{mesmo} número, mas com
significado diferente: após $\SI{39}{\minute}$ resta $37\%$ da tensão, não
zero. A descarga prática ($<1\%$) leva $5\tau\approx\SI{3.3}{\hour}$.\\
\textbf{Implicação de segurança:} um banco de capacitores pode reter tensão
letal por horas. Resistores de descarga permanentes não são opcionais.}
\end{questao}

\begin{questao}[3]{}
Compare a carga transportada e o aquecimento produzido por duas correntes de
mesma média: (A) $\SI{2}{\ampere}$ constantes; (B) $\SI{10}{\ampere}$ durante
$20\%$ do tempo e zero no restante.
\begin{itensq}
\item Determine a carga em $\SI{10}{\second}$ para cada uma.
\item Determine o valor eficaz de cada uma.
\item Compare o aquecimento em um resistor de $\SI{1}{\ohm}$.
\end{itensq}
\resposta{Mesma carga ($\SI{20}{\coulomb}$); $I_{rms}=2$ e
$\SI{4.47}{\ampere}$; aquecimento $5$ vezes maior em (B)}
\resolucao{(a) Ambas transportam $Q=I_{méd}t=2(10)=\SI{20}{\coulomb}$ ---
por construção.\\
(b) \textbf{(A)} $I_{rms}=\SI{2}{\ampere}$ (constante).\\
\textbf{(B)} $I_{rms}=\sqrt{0{,}2(10)^{2}}=\sqrt{20}=\SI{4.47}{\ampere}$.\\
(c)
\[
P_A=1(2)^{2}=\SI{4}{\watt};
\qquad
P_B=1(4{,}47)^{2}=\SI{20}{\watt}.
\]
\textbf{Cinco vezes mais}, exatamente o inverso do ciclo de trabalho.\\
\textbf{Síntese fundamental:} \emph{mesma carga, aquecimento completamente
diferente}. A carga depende da média; o aquecimento, do valor eficaz. Duas
grandezas independentes de uma mesma forma de onda.\\
\textbf{Consequência prática:} circuitos pulsados exigem condutores
dimensionados pelo valor \emph{eficaz}, ainda que a carga transportada seja
modesta.}
\end{questao}

\begin{questao}[3]{}
Um circuito recebe pulsos de $\SI{50}{\ampere}$ com duração de
$\SI{100}{\micro\second}$, a uma taxa de $\SI{200}{\per\second}$.
\begin{itensq}
\item Determine a carga por pulso e a corrente média.
\item Determine o ciclo de trabalho.
\item Determine o valor eficaz.
\end{itensq}
\resposta{$\SI{5}{\milli\coulomb}$ por pulso; $I_{méd}=\SI{1}{\ampere}$;
$D=2\%$; $I_{rms}=\SI{7.07}{\ampere}$}
\resolucao{(a) $Q=50\times\pot{100}{-6}=\SI{5}{\milli\coulomb}$.
Com $200$ pulsos por segundo,
\[
I_{méd}=200\times\pot{5}{-3}=\SI{1}{\ampere}.
\]
(b) $D=\dfrac{\pot{100}{-6}}{1/200}=\dfrac{\pot{100}{-6}}{\pot{5}{-3}}
=0{,}02=2\%$.\\
(c) $I_{rms}=I_{pico}\sqrt{D}=50\sqrt{0{,}02}=\SI{7.07}{\ampere}$.\\
\textbf{Três números, três finalidades.}
\begin{itemize}
\item $I_{méd}=\SI{1}{\ampere}$ dimensiona a \textbf{fonte} e a autonomia da
bateria.
\item $I_{rms}=\SI{7.07}{\ampere}$ dimensiona o \textbf{aquecimento} de
condutores e componentes.
\item $I_{pico}=\SI{50}{\ampere}$ dimensiona a \textbf{corrente máxima}
suportável pelos semicondutores e a queda instantânea na fonte.
\end{itemize}
Usar o número errado leva a projetos ou superdimensionados ou destruídos.}
\end{questao}

\blocodesafio

\begin{questao}[4]{}
\textbf{(Autonomia sob corrente variável.)} Uma bateria de
$\SI{2.4}{\ampere\hour}$ alimenta uma carga cuja corrente alterna entre
$\SI{0.4}{\ampere}$ por $\SI{20}{\minute}$ e $\SI{1.2}{\ampere}$ por
$\SI{5}{\minute}$, repetidamente.
\begin{itensq}
\item Determine a carga consumida por ciclo e a corrente média.
\item Determine a autonomia.
\item Discuta por que a autonomia real será menor.
\end{itensq}
\resposta{$\SI{0.233}{\ampere\hour}$ por ciclo; $I_{méd}=\SI{0.56}{\ampere}$;
autonomia $\approx\SI{4.3}{\hour}$}
\resolucao{(a) Convertendo os tempos para horas:
\[
Q_{ciclo}=0{,}4\left(\tfrac{20}{60}\right)+1{,}2\left(\tfrac{5}{60}\right)
=0{,}1333+0{,}1000=\SI{0.2333}{\ampere\hour}.
\]
O ciclo dura $\SI{25}{\minute}$, logo
\[
I_{méd}=\frac{0{,}2333}{25/60}=\SI{0.56}{\ampere}.
\]
(b)
\[
n=\frac{2{,}4}{0{,}2333}=10{,}3\ \text{ciclos}
\;\Longrightarrow\;
t=10{,}3(25)=\SI{257}{\minute}=\SI{4.3}{\hour}.
\]
\emph{(Conferindo: $2{,}4/0{,}56=\SI{4.3}{\hour}$ \checkmark.)}\\
\textbf{Observe:} o patamar de $\SI{1.2}{\ampere}$ dura apenas um quinto do
tempo do outro, mas consome $43\%$ da carga.\\
(c) \textbf{Por que a autonomia real é menor.}
\begin{itemize}
\item \textbf{Efeito Peukert:} a capacidade útil diminui com correntes mais
altas. Os picos de $\SI{1.2}{\ampere}$ extraem menos carga útil do que a
proporção sugere.
\item \textbf{Tensão de corte:} o equipamento para de funcionar antes de a
bateria se esgotar, ao atingir sua tensão mínima --- e isso ocorre \emph{mais
cedo} durante um pico de corrente, por causa da queda em $r$.
\item \textbf{Resistência interna crescente} ao longo da descarga, o que agrava
o item anterior.
\item \textbf{Temperatura:} capacidade menor no frio.
\end{itemize}
\textbf{Estimativa realista:} $\SI{4.3}{\hour}$ é um teto teórico; espere de
$3{,}5$ a $\SI{4}{\hour}$.}
\end{questao}

\begin{questao}[4]{}
\textbf{(Erro de estimativa pelo pico.)} Um técnico estima a carga de um pulso
triangular ($0\to\SI{6}{\ampere}\to0$ em $\SI{5}{\milli\second}$) como
$Q\approx I_{pico}T$.
\begin{itensq}
\item Determine o erro cometido.
\item Generalize para formas de onda quaisquer.
\item Proponha uma regra prática.
\end{itensq}
\resposta{Erro de $+100\%$ --- o dobro do valor real}
\resolucao{(a) \textbf{Valor real:}
$Q=\frac12(6)(0{,}005)=\SI{15}{\milli\coulomb}$.\\
\textbf{Estimativa:} $6(0{,}005)=\SI{30}{\milli\coulomb}$.\\
\[
\text{erro}=\frac{30-15}{15}=+100\%.
\]
(b) \textbf{Generalização.} Definindo o \emph{fator de forma}
$k=I_{méd}/I_{pico}$, a estimativa por pico erra sempre por um fator $1/k$:
\begin{center}
\begin{tabular}{lcc}
\hline
Forma de onda & $k$ & erro de $I_{pico}T$\\
\hline
Retangular (contínua) & $1{,}00$ & $0\%$\\
Trapezoidal (exemplo anterior) & $0{,}75$ & $+33\%$\\
Triangular & $0{,}50$ & $+100\%$\\
Exponencial ($T=5\tau$) & $0{,}20$ & $+400\%$\\
Pulsada ($D=2\%$) & $0{,}02$ & $+4900\%$\\
\hline
\end{tabular}
\end{center}
\textbf{O erro é sempre para mais} --- a estimativa por pico é
\emph{conservadora}.\\
(c) \textbf{Regra prática.}
\begin{itemize}
\item Se a intenção é \textbf{garantir margem} (dimensionar uma bateria, por
exemplo), a estimativa por pico é segura, mas pode superdimensionar
absurdamente.
\item Se a intenção é \textbf{prever} a carga, integre --- ou ao menos use o
fator de forma da onda em questão.
\item \textbf{Nunca} use $I_{pico}T$ para estimar autonomia: o
superdimensionamento chega a ordens de grandeza em sinais pulsados.
\end{itemize}}
\end{questao}

\begin{questao}[4]{}
\textbf{(Projeto de forma de onda.)} Projete uma corrente em \textbf{dois
patamares} que transfira exatamente $\SI{10}{\coulomb}$ em
$\SI{4}{\second}$, com valores comerciais simples.
\begin{itensq}
\item Estabeleça a condição geral.
\item Proponha duas soluções e compare o valor eficaz.
\item Indique qual é preferível e por quê.
\end{itensq}
\resposta{Condição: $I_1t_1+I_2t_2=10$ com $t_1+t_2=4$; a solução mais
uniforme minimiza o aquecimento}
\resolucao{(a) \textbf{Condição:}
\[
I_1t_1+I_2t_2=\SI{10}{\coulomb};
\qquad
t_1+t_2=\SI{4}{\second}.
\]
Duas equações e quatro incógnitas --- o problema tem \textbf{dois graus de
liberdade}, e admite infinitas soluções. É preciso um critério adicional.\\
(b) \textbf{Solução A:} $\SI{4}{\ampere}$ por $\SI{1}{\second}$ e
$\SI{2}{\ampere}$ por $\SI{3}{\second}$.
\[
Q=4+6=\SI{10}{\coulomb}\ \checkmark;
\qquad
I_{rms}=\sqrt{\frac{16(1)+4(3)}{4}}=\sqrt{7}=\SI{2.65}{\ampere}.
\]
\textbf{Solução B:} $\SI{8}{\ampere}$ por $\SI{1}{\second}$ e
$\SI{0.667}{\ampere}$ por $\SI{3}{\second}$.
\[
Q=8+2=\SI{10}{\coulomb}\ \checkmark;
\qquad
I_{rms}=\sqrt{\frac{64(1)+0{,}444(3)}{4}}=\sqrt{16{,}3}=\SI{4.04}{\ampere}.
\]
(c) \textbf{A solução A é preferível.} Ela tem valor eficaz $34\%$ menor e,
portanto, aquecimento $2{,}3$ vezes menor --- entregando exatamente a mesma
carga.\\
\textbf{Princípio geral:} para uma dada carga em um dado tempo, o valor eficaz é
\textbf{mínimo} quando a corrente é \emph{constante}
($I=10/4=\SI{2.5}{\ampere}$, $I_{rms}=\SI{2.5}{\ampere}$). Qualquer
variação aumenta $I_{rms}$ acima da média --- é a mesma desigualdade
$\langle i^{2}\rangle\ge\langle i\rangle^{2}$ vista antes.\\
\textbf{Conclusão de projeto:} se a aplicação permitir, use corrente constante.
Se exigir patamares, mantenha-os o mais próximos possível entre si.}
\end{questao}

\begin{questao}[4]{}
\textbf{(Corrente bipolar.)} Uma corrente periódica é positiva em parte do
ciclo e negativa no restante.
\begin{itensq}
\item Formule a condição para carga líquida nula.
\item Determine se isso implica ausência de efeitos.
\item Cite duas aplicações em que essa condição é imposta deliberadamente.
\end{itensq}
\resposta{$\int_0^{T}i\,\mathrm{d}t=0$; os efeitos térmicos e mecânicos
permanecem}
\resolucao{(a) \textbf{Condição:}
\[
\int_0^{T}i(t)\,\mathrm{d}t=0,
\]
ou seja, as áreas positiva e negativa devem ser iguais em módulo. Em termos de
patamares, $I_{+}t_{+}=|I_{-}|t_{-}$.\\
(b) \textbf{Não implica ausência de efeitos.} O que se anula é apenas o
\emph{acúmulo} de carga. Permanecem:
\begin{itemize}
\item \textbf{Aquecimento:} $P=RI_{rms}^{2}$, com
$I_{rms}=\sqrt{\frac1T\int i^{2}\mathrm{d}t}>0$ sempre.
\item \textbf{Forças magnéticas} sobre condutores, que dependem de $i$
instantâneo.
\item \textbf{Desgaste eletroquímico}, se as reações direta e inversa não forem
perfeitamente simétricas.
\end{itemize}
(c) \textbf{Duas aplicações da condição imposta deliberadamente.}
\begin{itemize}
\item \textbf{Estimulação elétrica médica} (marca-passos, eletroestimuladores):
pulsos bifásicos com carga líquida rigorosamente nula evitam acúmulo de
subprodutos eletroquímicos no tecido e corrosão dos eletrodos. É requisito
normativo, não preferência.
\item \textbf{Acionamento de transformadores:} carga líquida não nula
introduziria componente contínua, que satura o núcleo magnético e provoca
corrente de magnetização excessiva. Conversores em ponte são projetados para
garantir simetria.
\end{itemize}
\textbf{E há um terceiro caso instrutivo:} em galvanoplastia usa-se
\emph{deliberadamente} carga líquida não nula --- ali o acúmulo é o objetivo.}
\end{questao}

\begin{questao}[4]{}
\textbf{(Medição --- média contra instantâneo.)} Dois instrumentos medem a
mesma corrente pulsada: um indica a média temporal, outro o valor instantâneo.
\begin{itensq}
\item Explique por que os resultados diferem.
\item Determine qual usar para estimar carga, aquecimento e pico.
\item Projete uma medição que forneça os três.
\end{itensq}
\resposta{Média para carga, eficaz para aquecimento, instantâneo para pico}
\resolucao{(a) \textbf{Os instrumentos medem grandezas diferentes de uma mesma
forma de onda.} Um amperímetro de bobina móvel responde à média (por inércia
mecânica); um osciloscópio mostra o valor instantâneo; um multímetro "true RMS"
calcula o valor eficaz.\\
Para a corrente pulsada de $\SI{50}{\ampere}$ com $D=2\%$, os três leriam
$\SI{1}{\ampere}$, $\SI{50}{\ampere}$ e $\SI{7.07}{\ampere}$ --- todos
corretos, todos respondendo perguntas distintas.\\
(b)
\begin{center}
\begin{tabular}{ll}
\hline
Objetivo & Grandeza\\
\hline
Carga transferida, autonomia de bateria & \textbf{média}\\
Aquecimento, dimensionamento de condutor & \textbf{eficaz}\\
Corrente máxima em semicondutor, saturação & \textbf{pico}\\
\hline
\end{tabular}
\end{center}
(c) \textbf{Projeto de medição completa.} Use um \emph{resistor shunt} de baixo
valor e observe a tensão sobre ele com osciloscópio digital:
\begin{enumerate}
\item O \textbf{traço instantâneo} fornece o pico diretamente.
\item A função \textbf{média} do osciloscópio integra a forma de onda e entrega
$I_{méd}$ --- e portanto a carga, multiplicando pelo tempo.
\item A função \textbf{RMS} entrega $I_{rms}$.
\end{enumerate}
\textbf{Cuidados:} a largura de banda do osciloscópio e do shunt deve ser muito
maior que a taxa de subida dos pulsos, sob pena de arredondá-los e subestimar o
pico. E a indutância parasita do shunt introduz erro proporcional a
$\mathrm{d}i/\mathrm{d}t$ --- use shunts coaxiais em pulsos rápidos.}
\end{questao}

\begin{questao}[4]{}
\textbf{(Integração de sinal ruidoso.)} Uma corrente medida contém ruído
aditivo de média zero.
\begin{itensq}
\item Explique por que integrar melhora a estimativa da carga.
\item Determine como o erro decresce com o tempo de integração.
\item Indique em que situação integrar \emph{piora} o resultado.
\end{itensq}
\resposta{O erro relativo decresce como $1/\sqrt{N}$; integrar piora se houver
\emph{deriva} (ruído de média não nula)}
\resolucao{(a) \textbf{Por que integrar ajuda.} A carga é
$Q=\int i\,\mathrm{d}t$. Se o ruído tem média zero, suas contribuições
positivas e negativas tendem a se cancelar na integral, enquanto o sinal
verdadeiro se acumula. A relação sinal-ruído \emph{melhora} com o tempo.\\
(b) \textbf{Decrescimento do erro.} Para $N$ amostras independentes de ruído de
desvio $\sigma$, o desvio da \emph{soma} é $\sigma\sqrt{N}$, enquanto a soma do
sinal cresce com $N$. Logo o erro relativo cai como
\[
\frac{\sigma\sqrt{N}}{SN}=\frac{\sigma}{S\sqrt{N}}\propto\frac{1}{\sqrt{T}}.
\]
\textbf{Quantificando:} integrar por $\SI{100}{\second}$ em vez de
$\SI{1}{\second}$ reduz o erro relativo por um fator $10$. Para ganhar outro
fator $10$, seriam necessários $\SI{10000}{\second}$ --- rendimento
decrescente.\\
(c) \textbf{Quando integrar piora.} Se o ruído \textbf{não} tiver média zero ---
isto é, se houver \emph{offset} ou \emph{deriva} no instrumento --- o erro
\textbf{também se acumula}, e cresce \emph{linearmente} com $T$:
\[
\text{erro}=\varepsilon_{offset}\times T.
\]
Enquanto o ruído aleatório cai com $\sqrt{T}$, o erro sistemático cresce com
$T$. Existe, portanto, um \textbf{tempo ótimo} de integração, além do qual a
deriva domina.\\
\textbf{Prática de laboratório:} meça a leitura com \emph{entrada em curto}
antes do ensaio, para quantificar o \emph{offset}, e subtraia-o. Sem essa
correção, integrações longas são pior que inúteis --- elas produzem números
precisos e errados.}
\end{questao}

\begin{questao}[4]{}
\textbf{(Por que a LKC é tão precisa.)} A lei dos nós afirma que a carga não se
acumula em um nó. Quantifique o quanto ela poderia, em princípio, ser violada.
\begin{itensq}
\item Estime a capacitância parasita de um nó típico e relacione carga acumulada
      com potencial.
\item Determine o desequilíbrio de corrente necessário para elevar o nó em
      $\SI{1}{\volt}$.
\item Conclua sobre o caráter da lei.
\end{itensq}
\resposta{Com $C\approx\SI{1}{\pico\farad}$, basta $\SI{1}{\pico\coulomb}$ para
$\SI{1}{\volt}$ --- um desequilíbrio de $\SI{1}{\micro\ampere}$ por
$\SI{1}{\micro\second}$}
\resolucao{(a) Um nó real (uma solda, um trecho de trilha) tem capacitância
parasita para o ambiente da ordem de $C\approx\SI{1}{\pico\farad}$. Se uma
carga $\Delta Q$ se acumulasse ali, o potencial subiria:
\[
\Delta V=\frac{\Delta Q}{C}.
\]
Para $\Delta V=\SI{1}{\volt}$:
\[
\Delta Q=\pot{1}{-12}\,\si{\coulomb}
=\frac{\pot{1}{-12}}{\pot{1.6}{-19}}=\pot{6.25}{6}\ \text{elétrons}.
\]
Apenas seis milhões de elétrons --- número minúsculo diante dos
$10^{19}$ que atravessam a seção a cada segundo com $\SI{2}{\ampere}$.\\
(b) O desequilíbrio necessário é irrisório:
\[
\Delta I\,\Delta t=\SI{1}{\pico\coulomb}
\;\Longrightarrow\;
\SI{1}{\micro\ampere}\ \text{durante}\ \SI{1}{\micro\second},
\]
ou $\SI{1}{\nano\ampere}$ durante $\SI{1}{\milli\second}$. Em um circuito de
$\SI{2}{\ampere}$, isso corresponde a um desequilíbrio relativo de
$\pot{5}{-7}$ --- meio milionésimo.\\
(c) \textbf{Conclusão: a LKC não é uma aproximação, é uma consequência
autorregulada.} Qualquer tentativa de acumular carga em um nó eleva
imediatamente seu potencial, e esse potencial \emph{expulsa} o excesso --- as
correntes se reajustam em nanossegundos.\\
\textbf{Onde a lei realmente falha:} em frequências altíssimas, quando o tempo
de propagação pelo circuito deixa de ser desprezível, correntes de deslocamento
por capacitâncias parasitas passam a importar e a LKC "de nó concentrado" perde
sentido. É preciso então tratar o circuito por linhas de transmissão ou por
campos.\\
\textbf{Critério prático:} a LKC vale enquanto as dimensões do circuito forem
muito menores que o comprimento de onda do sinal --- a chamada aproximação de
\emph{parâmetros concentrados}.}
\end{questao}

\begin{questao}[4]{}
\textbf{(Deriva contra propagação.)} Um fio de cobre de
$\SI{2.5}{\milli\meter\squared}$ conduz $\SI{10}{\ampere}$, com densidade de
portadores $n=\pot{8.5}{28}\,\si{\per\meter\cubed}$.
\begin{itensq}
\item Determine a velocidade de deriva dos elétrons.
\item Determine o tempo para um elétron percorrer $\SI{1}{\meter}$.
\item Compare com o tempo de propagação do sinal e explique o paradoxo
      aparente.
\end{itensq}
\resposta{$v_d=\SI{0.29}{\milli\meter\per\second}$; $\SI{57}{\minute}$ para
$\SI{1}{\meter}$; o sinal leva $\SI{5}{\nano\second}$}
\resolucao{(a)
\[
v_d=\frac{I}{nAe}
=\frac{10}{(\pot{8.5}{28})(\pot{2.5}{-6})(\pot{1.6}{-19})}
=\frac{10}{\pot{3.4}{4}}=\pot{2.94}{-4}\,\si{\meter\per\second}.
\]
Menos de um terço de milímetro por segundo.\\
(b)
\[
t=\frac{1}{\pot{2.94}{-4}}=\SI{3400}{\second}\approx\SI{57}{\minute}.
\]
(c) \textbf{O sinal viaja a cerca de $\pot{2}{8}\,\si{\meter\per\second}$}
(dois terços de $c$, em cabo típico), percorrendo o metro em
$\SI{5}{\nano\second}$ --- \textbf{$\pot{6.8}{11}$ vezes mais rápido} que os
elétrons.\\
\textbf{Resolução do paradoxo aparente.} O que se propaga rapidamente não são os
elétrons, mas o \textbf{campo eletromagnético} que os põe em movimento. Ao
fechar o interruptor, o campo se estabelece ao longo de todo o condutor quase
instantaneamente, e \emph{todos} os elétrons começam a se mover praticamente ao
mesmo tempo --- cada um localmente, e devagar.\\
\textbf{Analogia:} um cano cheio de água. Ao abrir a torneira, a água sai
imediatamente na outra ponta, embora nenhuma molécula tenha percorrido o cano.
O que se propaga é a \emph{pressão}, não a matéria.\\
\textbf{Consequência prática:} a lâmpada acende no instante em que se aciona o
interruptor, mas os elétrons que estavam junto dele levariam quase uma hora para
chegar até ela.}
\end{questao}

% END BANCO COMPLEMENTAR

\gabarito[2.1 Corrente elétrica e carga]

% =====================================================================
\subsection{Lei de Ohm e resistividade}
% =====================================================================

\blocofixacao

\begin{questao}[1]{}
De acordo com a segunda lei de Ohm, a resistência de um fio condutor:
\begin{alternativas}
\item aumenta com a área da seção e diminui com o comprimento.
\item aumenta com o comprimento e diminui com a área da seção.
\item depende apenas do material.
\item independe da temperatura.
\item é sempre igual para todos os materiais.
\end{alternativas}
\resposta{(b)}
\resolucao{De $R = \rho\dfrac{L}{A}$: a resistência é \emph{diretamente}
proporcional ao comprimento $L$ e \emph{inversamente} proporcional à área $A$.
O material entra pela resistividade $\rho$, que por sua vez varia com a
temperatura --- o que torna (c) e (d) incorretas.}
\end{questao}

\begin{questao}[1]{}
O gráfico mostra a resistência de um fio em função de seu comprimento, mantida a
área constante.

\circuitoqq{%
\begin{tikzpicture}
 \begin{axis}[width=9cm,height=5cm,xlabel={Comprimento (m)},
   ylabel={$R$ ($\Omega$)},xmin=0,xmax=4,ymin=0,ymax=7,
   xtick={0,1,2,3},ytick={0,2,4,6},grid=major,axis lines=left]
   \addplot[Rcol,very thick,domain=0:3] {2*x};
   \addplot[only marks,mark=*,Rcol] coordinates {(1,2)(2,4)(3,6)};
 \end{axis}
\end{tikzpicture}}{Resistência proporcional ao comprimento.}

\begin{itensq}
\item Qual é a constante de proporcionalidade $k$ da reta $R = kL$?
\item Qual seria a resistência de um fio de $\SI{5}{\meter}$ do mesmo tipo?
\item Que lei essa proporcionalidade confirma?
\end{itensq}
\resposta{(a) $k=\SI{2}{\ohm\per\meter}$; (b) \SI{10}{\ohm}; (c) 2ª lei de Ohm}
\resolucao{(a) A reta passa por $(3, 6)$, logo
$k = \dfrac{R}{L} = \dfrac{6}{3} = \SI{2}{\ohm\per\meter}$.\\
(b) $R = kL = 2\cdot 5 = \SI{10}{\ohm}$.\\
(c) A proporcionalidade direta entre $R$ e $L$ (com $A$ e material fixos) confirma
a segunda lei de Ohm, $R = \rho\dfrac{L}{A}$, em que $k = \dfrac{\rho}{A}$.}
\end{questao}

\blocoaplicacao

\begin{questao}[2]{}
Um estudante quer construir um resistor de $\SI{10}{\ohm}$ com fio de
níquel-cromo, de resistividade $\rho = \SI{1.1e-6}{\ohm\meter}$ e seção
$A = \SI{0.5}{\milli\meter\squared}$. Qual é o comprimento necessário?
\resposta{$L \approx \SI{4.55}{\meter}$}
\resolucao{Isolando $L$ na segunda lei de Ohm (com
$A = \SI{0.5}{\milli\meter\squared} = \pot{0.5}{-6}\,\si{\meter\squared}$):
\[
L=\frac{R\,A}{\rho}=\frac{10\cdot\pot{0.5}{-6}}{\pot{1.1}{-6}}
=\frac{\pot{5}{-6}}{\pot{1.1}{-6}}\approx\SI{4.55}{\meter}.
\]}
\end{questao}

\begin{questao}[2]{}
Um fio condutor de comprimento $L$, diâmetro $D$ e resistência $R$ tem seu
comprimento \textbf{e} seu diâmetro \emph{duplicados}. A nova resistência será:
\begin{alternativas}[3]
\item $R$
\item $2R$
\item $R/2$
\item $4R$
\item $R/4$
\end{alternativas}
\resposta{(c)}
\resolucao{A área da seção é proporcional ao quadrado do diâmetro, $A \propto D^2$.
Ao dobrar $D$, a área quadruplica ($A' = 4A$); ao dobrar $L$, o comprimento
duplica ($L' = 2L$):
\[
R'=\rho\frac{L'}{A'}=\rho\frac{2L}{4A}=\frac{1}{2}\,\rho\frac{L}{A}=\frac{R}{2}.
\]
A armadilha (alternativa b) é considerar só o efeito do comprimento e esquecer que
a área cresce com o \emph{quadrado} do diâmetro.}
\end{questao}

\begin{questao}[2]{}
O gráfico representa a diferença de potencial $U$ entre dois pontos de um fio em
função da corrente $i$. A resistência do trecho, em ohms, vale:

\circuitoqq{%
\begin{tikzpicture}
 \begin{axis}[width=9cm,height=5cm,xlabel={$i$ (A)},ylabel={$U$ (V)},
   xmin=0,xmax=5,ymin=0,ymax=22,xtick={0,1,2,3,4},ytick={0,4,8,12,16,20},
   grid=major,axis lines=left]
   \addplot[Vcol,very thick,domain=0:4] {4*x};
   \addplot[only marks,mark=*,Vcol] coordinates {(1,4)(2,8)(3,12)(4,16)};
 \end{axis}
\end{tikzpicture}}{Curva característica $U \times i$ de um resistor ôhmico.}

\begin{alternativas}[3]
\item 0,05
\item 4
\item 20
\item 80
\item 160
\end{alternativas}
\resposta{(b)}
\resolucao{Para um condutor ôhmico, $U = R\,i$: a resistência é a \emph{inclinação}
da reta. Tomando o ponto $(4, 16)$:
\[
R=\frac{U}{i}=\frac{16}{4}=\SI{4}{\ohm}.
\]
A reta passar pela origem confirma que o resistor obedece à lei de Ohm.}
\end{questao}

\begin{questao}[2]{}
Dois fios do mesmo material têm resistências $R_1 = \SI{6}{\ohm}$ e
$R_2 = \SI{3}{\ohm}$; o fio 1 tem o dobro do comprimento do fio 2. Qual é a relação
entre as áreas das seções transversais, $A_1/A_2$?
\resposta{$A_1/A_2 = 1$ (áreas iguais)}
\resolucao{Como o material é o mesmo, $\rho$ é comum. De $R = \rho L/A$:
\[
\frac{R_1}{R_2}=\frac{L_1/A_1}{L_2/A_2}=\frac{L_1}{L_2}\cdot\frac{A_2}{A_1}.
\]
Com $\dfrac{R_1}{R_2}=\dfrac{6}{3}=2$ e $\dfrac{L_1}{L_2}=2$:
\[
2 = 2\cdot\frac{A_2}{A_1}\ \Rightarrow\ \frac{A_2}{A_1}=1
\ \Rightarrow\ A_1 = A_2.
\]
Ou seja, dobrar o comprimento já explica a resistência dobrada; as seções são
iguais.}
\end{questao}

\begin{questao}[2]{}
A tabela apresenta a resistividade de alguns materiais a $\SI{20}{\celsius}$.

\begin{table}[H]
\centering\small\sffamily
\begin{tabular}{@{}l c@{}}
\toprule
\textbf{Material} & \textbf{$\rho$ ($\times10^{-8}\,\si{\ohm\meter}$)}\\
\midrule
Prata      & 1,6\\
Cobre      & 1,7\\
Alumínio   & 2,8\\
Ferro      & 10\\
Níquel-cromo & 110\\
\bottomrule
\end{tabular}
\caption{Resistividade de materiais condutores.}
\end{table}

Dois fios de mesmo comprimento e mesma seção, um de cobre e outro de alumínio,
têm resistências $R_{\text{Cu}}$ e $R_{\text{Al}}$. A razão
$R_{\text{Al}}/R_{\text{Cu}}$ vale aproximadamente:
\begin{alternativas}[3]
\item 0,6
\item 1,0
\item 1,6
\item 2,8
\item 4,4
\end{alternativas}
\resposta{(c)}
\resolucao{Com $L$ e $A$ iguais, $R = \rho\dfrac{L}{A}$ é proporcional a $\rho$:
\[
\frac{R_{\text{Al}}}{R_{\text{Cu}}}=\frac{\rho_{\text{Al}}}{\rho_{\text{Cu}}}
=\frac{2{,}8}{1{,}7}\approx1{,}6.
\]
Por isso, para a mesma resistência, um fio de alumínio precisa ser mais grosso
que um de cobre --- o que se compensa por ele ser mais leve e barato, motivo do
seu uso em linhas de transmissão.}
\end{questao}

\begin{questao}[2]{}
Um fio de cobre tem comprimento $L = \SI{2}{\meter}$ e seção transversal
$A = \SI{1}{\milli\meter\squared}$. Calcule sua resistência.
\dados{$\rho_{\text{Cu}} = \SI{1.7e-8}{\ohm\meter}$.}
\resposta{$R = \SI{34}{\milli\ohm}$}
\resolucao{Com $A = \SI{1}{\milli\meter\squared} = \pot{1}{-6}\,\si{\meter\squared}$:
\[
R=\rho\frac{L}{A}=\pot{1.7}{-8}\cdot\frac{2}{\pot{1}{-6}}
=\pot{3.4}{-2}\,\si{\ohm}=\SI{34}{\milli\ohm}.
\]
Resistência bem baixa, como esperado de um bom condutor em curto comprimento.}
\end{questao}

\begin{questao}[2]{Apostila original revisada}
Um fio possui resistência elétrica $R$. Mantendo o mesmo material, seu comprimento é
duplicado e sua área de seção transversal é reduzida à metade. Determine a nova
resistência em função de $R$.
\resposta{$R'=4R$}
\resolucao{Pela segunda lei de Ohm, $R=\rho L/A$. Logo,
\[
R'=\rho\frac{2L}{A/2}=4\rho\frac{L}{A}=4R.
\]}
\end{questao}

\blocoaprofundamento

\begin{questao}[3]{}
Um fio de aço e um de cobre, de mesmo comprimento e mesma área, são ligados
individualmente à mesma tensão. Sabendo que
$\rho_{\text{aço}} = 10\,\rho_{\text{cobre}}$, compare:
\begin{itensq}
\item as correntes em cada fio;
\item as potências dissipadas em cada um.
\end{itensq}
\resposta{(a) $I_{\text{cobre}} = 10\,I_{\text{aço}}$;
(b) $P_{\text{cobre}} = 10\,P_{\text{aço}}$}
\resolucao{Como $L$ e $A$ são iguais, $R = \rho L/A \Rightarrow
R_{\text{aço}} = 10\,R_{\text{cobre}}$.\\
(a) Sob a mesma tensão, $I = U/R$, logo a corrente é inversamente proporcional a
$R$:
\[
\frac{I_{\text{cobre}}}{I_{\text{aço}}}
=\frac{R_{\text{aço}}}{R_{\text{cobre}}}=10
\ \Rightarrow\ I_{\text{cobre}}=10\,I_{\text{aço}}.
\]
(b) $P = U^2/R$; com $U$ igual, a potência também é inversamente proporcional a
$R$:
\[
\frac{P_{\text{cobre}}}{P_{\text{aço}}}=\frac{R_{\text{aço}}}{R_{\text{cobre}}}=10.
\]
O cobre, menos resistivo, conduz mais corrente e dissipa mais potência sob a mesma
tensão --- por isso é preferido em condutores.}
\end{questao}

\begin{questao}[3]{}
Dois fios de \emph{mesmo comprimento} e \emph{mesma seção}, um de cobre
($\rho_{\text{Cu}}=\SI{1.7e-8}{\ohm\meter}$) e outro de alumínio
($\rho_{\text{Al}}=\SI{2.8e-8}{\ohm\meter}$), são emendados em série e ligados a
$\SI{12}{\volt}$.
\begin{itensq}
\item Determine a razão entre as tensões em cada trecho.
\item Calcule a tensão em cada fio.
\item Em qual dos dois há maior dissipação de calor? Justifique.
\end{itensq}
\resposta{(a) $U_{\text{Al}}/U_{\text{Cu}}\approx1{,}65$;
(b) $\SI{4.53}{\volt}$ e $\SI{7.47}{\volt}$; (c) no alumínio}
\resolucao{(a) Com $L$ e $A$ iguais, $R \propto \rho$; e como estão em série
(mesma corrente), $U = RI \propto \rho$:
\[
\frac{U_{\text{Al}}}{U_{\text{Cu}}}=\frac{\rho_{\text{Al}}}{\rho_{\text{Cu}}}
=\frac{2{,}8}{1{,}7}\approx1{,}65.
\]
(b) Pelo divisor de tensão, a repartição segue as resistividades:
\[
U_{\text{Cu}}=12\cdot\frac{1{,}7}{4{,}5}\approx\SI{4.53}{\volt},
\qquad
U_{\text{Al}}=12\cdot\frac{2{,}8}{4{,}5}\approx\SI{7.47}{\volt}.
\]
(c) A potência é $P = UI$ com $I$ comum, logo \emph{maior tensão implica maior
potência}: o \textbf{alumínio} aquece mais. Fisicamente, sendo mais resistivo,
ele opõe mais resistência à mesma corrente.\\
\textbf{Implicação prática:} emendas entre metais diferentes são pontos críticos
em instalações --- além do aquecimento desigual, os coeficientes de dilatação
distintos afrouxam o contato com o tempo, aumentando a resistência local e
realimentando o aquecimento. É a causa clássica de incêndios em conexões
Cu--Al, razão pela qual normas exigem conectores bimetálicos específicos.}
\end{questao}

\begin{questao}[3]{}
A resistência de um condutor metálico varia com a temperatura segundo
$R = R_0\big[1+\alpha(T-T_0)\big]$, onde $\alpha$ é o coeficiente de temperatura.
Para o cobre, $\alpha \approx \SI{4e-3}{\per\celsius}$.

Um enrolamento de cobre de um motor mede $\SI{10}{\ohm}$ a $\SI{20}{\celsius}$
(motor frio). Após operar em regime, sua resistência sobe para
$\SI{12.4}{\ohm}$.
\begin{itensq}
\item Determine a temperatura do enrolamento em operação.
\item Se a tensão aplicada for mantida em $\SI{220}{\volt}$, qual é a variação
      percentual da potência dissipada entre o motor frio e o quente?
\item Explique como esse efeito é usado, na prática, para medir a temperatura
      interna de máquinas elétricas sem sensores.
\end{itensq}
\resposta{(a) $T = \SI{80}{\celsius}$; (b) a potência cai $\approx\SI{19}{\percent}$;
(c) ver comentário}
\resolucao{\textbf{(a)} Isolando a temperatura:
\[
12{,}4 = 10\big[1+\pot{4}{-3}(T-20)\big]
\;\Rightarrow\;
1{,}24 = 1+\pot{4}{-3}(T-20)
\;\Rightarrow\;
T-20 = \frac{0{,}24}{\pot{4}{-3}} = 60,
\]
logo $T = \SI{80}{\celsius}$.

\textbf{(b)} Com tensão fixa, $P = \dfrac{U^2}{R}$ é inversamente proporcional a
$R$:
\[
\frac{P_{\text{quente}}}{P_{\text{frio}}}=\frac{R_{\text{frio}}}{R_{\text{quente}}}
=\frac{10}{12{,}4}=0{,}806,
\]
ou seja, a potência dissipada \emph{cai} cerca de \SI{19.4}{\percent}.

\textbf{(c) Método da variação de resistência.} Como $\alpha$ do cobre é bem
conhecido e estável, mede-se a resistência do enrolamento a frio (com o motor
parado e em temperatura ambiente conhecida) e novamente logo após o desligamento.
Invertendo a fórmula:
\[
T = T_0 + \frac{1}{\alpha}\left(\frac{R}{R_0}-1\right),
\]
obtém-se a temperatura \emph{média} de todo o enrolamento --- justamente o que
importa para a vida útil do isolamento --- sem instalar nenhum sensor interno.
Esse é o procedimento normalizado em ensaios de elevação de temperatura de
motores e transformadores.\\
\textbf{Observação crítica:} o método fornece a média, não o \emph{ponto quente}.
Como a degradação do isolante segue uma lei exponencial com a temperatura, o ponto
mais quente pode falhar antes do que a média sugeriria --- por isso as normas
aplicam margens de segurança sobre esse valor.}
\end{questao}

\begin{questao}[3]{Apostila original revisada}
A resistividade de um material varia aproximadamente de forma linear com a
temperatura. Foram medidos os valores
$\rho(\SI{0}{\celsius})=\SI{1.5e-8}{\ohm\metre}$ e
$\rho(\SI{100}{\celsius})=\SI{2.5e-8}{\ohm\metre}$.
\begin{itensq}
\item Determine a taxa de variação da resistividade por grau Celsius.
\item Estime a resistividade a $\SI{150}{\celsius}$, mantendo o modelo linear.
\end{itensq}
\resposta{(a) $\SI{1.0e-10}{\ohm\metre\per\celsius}$; (b) $\SI{3.0e-8}{\ohm\metre}$}
\resolucao{A inclinação é
\[
\frac{\Delta\rho}{\Delta T}=\frac{(2.5-1.5)\times10^{-8}}{100}
=\SI{1.0e-10}{\ohm\metre\per\celsius}.
\]
Extrapolando até $\SI{150}{\celsius}$:
$\rho=1.5\times10^{-8}+150\times10^{-10}=\SI{3.0e-8}{\ohm\metre}$.}
\end{questao}

% BEGIN BANCO COMPLEMENTAR Lei de Ohm e resistividade
% =====================================================================
%  Banco complementar --- Lei de Ohm e Resistividade  (REVISADO)
%  Substitui a versao gerada por template (8 clones em N2, 11 em N3 e
%  10 questoes conceituais marcadas como N4).
%  A secao correspondente do cap02 e densa e ja cobre: 2a lei de Ohm
%  (MC), grafico R x L, resistor de 10 ohm com niquel-cromo, fio com
%  L e D dobrados, grafico U x i, dois fios de mesmo material, tabela
%  de resistividades, fio de cobre 2 m / 1 mm2, L dobrado com A
%  dobrada, aco x cobre, cobre x aluminio, R=R0[1+alfa(T-T0)] e
%  resistividade linear com T a partir de dados.
%  Este banco evita esses casos e explora: projeto de resistor de fio,
%  faixa total de R (tolerancia + deriva + autoaquecimento), R estatica
%  x dinamica, secao de cabo por queda de tensao, fio ESTICADO a volume
%  constante, secao variavel por integracao, sensor a corrente
%  constante, medicao a quatro fios e estabilidade termica.
%  Distribuicao mantida: 8 N1 + 8 N2 + 11 N3 + 10 N4 = 37
% =====================================================================

\subsubsection*{Banco complementar --- Lei de Ohm e resistividade}

\begin{atencao}[Organização do banco]
Duas leis distintas convivem aqui. A \textbf{primeira} ($U=RI$) é uma relação
\emph{de circuito} e vale apenas para componentes ôhmicos. A \textbf{segunda}
($R=\rho L/A$) é uma relação \emph{de geometria e material}, e vale sempre. O N3
trata de projeto de resistores, seção de cabos, sensores e medição; o N4 exige
integração, demonstrações e análise de estabilidade térmica.
\end{atencao}

\blocofixacao

\begin{questao}[1]{}
Um resistor de $\SI{470}{\ohm}$ é percorrido por $\SI{15}{\milli\ampere}$.
Determine a tensão em seus terminais.
\resposta{$U=\SI{7.05}{\volt}$}
\resolucao{
\[
U=RI=470\times\pot{15}{-3}=\SI{7.05}{\volt}.
\]
\textbf{Cuidado com o prefixo:} usar $I=\SI{15}{\ampere}$ daria
$\SI{7050}{\volt}$ --- resultado absurdo, mas que passa despercebido se a
conversão não for feita explicitamente.}
\end{questao}

\begin{questao}[1]{}
Uma tensão de $\SI{9}{\volt}$ produz corrente de
$\SI{30}{\milli\ampere}$ em um componente ôhmico. Determine sua resistência.
\resposta{$R=\SI{300}{\ohm}$}
\resolucao{
\[
R=\frac{U}{I}=\frac{9}{0{,}030}=\SI{300}{\ohm}.
\]
\textbf{Atenção à hipótese:} a razão $U/I$ sempre pode ser calculada, mas só se
chama \emph{resistência} no sentido pleno se o componente for ôhmico --- isto é,
se essa razão for a mesma para qualquer tensão aplicada. Em componentes não
lineares ela varia, e recebe o nome de \emph{resistência estática}.}
\end{questao}

\begin{questao}[1]{}
Determine a corrente em um resistor de $\SI{2.2}{\kilo\ohm}$ submetido a
$\SI{12}{\volt}$.
\resposta{$I=\SI{5.45}{\milli\ampere}$}
\resolucao{
\[
I=\frac{U}{R}=\frac{12}{2200}=\pot{5.45}{-3}\,\si{\ampere}
=\SI{5.45}{\milli\ampere}.
\]
\textbf{Atalho útil:} volts divididos por quilohms dão \emph{miliampères}
diretamente --- $12/2{,}2=5{,}45$. Essa combinação de prefixos aparece o tempo
todo em eletrônica e poupa conversões.}
\end{questao}

\begin{questao}[1]{}
Um condutor de cobre ($\rho=\pot{1.7}{-8}\,\si{\ohm\meter}$) tem
$\SI{50}{\meter}$ e seção de $\SI{2.5}{\milli\meter\squared}$. Determine sua
resistência.
\resposta{$R=\SI{0.34}{\ohm}$}
\resolucao{Convertendo a área para $\si{\meter\squared}$:
$A=\SI{2.5}{\milli\meter\squared}=\pot{2.5}{-6}\,\si{\meter\squared}$.
\[
R=\frac{\rho L}{A}=\frac{\pot{1.7}{-8}\times50}{\pot{2.5}{-6}}
=\SI{0.34}{\ohm}.
\]
\textbf{Erro clássico:} esquecer que
$\SI{1}{\milli\meter\squared}=\pot{1}{-6}\,\si{\meter\squared}$, e não
$\pot{1}{-3}$. A área é o \emph{quadrado} de um comprimento, logo o fator de
conversão também é quadrático.}
\end{questao}

\begin{questao}[1]{}
Determine a condutância de um resistor de $\SI{25}{\ohm}$.
\resposta{$G=\SI{0.04}{\ensuremath{\mathrm{S}}}$}
\resolucao{
\[
G=\frac{1}{R}=\frac{1}{25}=\SI{0.04}{\ensuremath{\mathrm{S}}}.
\]
A condutância mede a \emph{facilidade} de conduzir. Ela é a grandeza natural
quando os elementos estão em paralelo (somam-se diretamente) e na análise nodal.
Nas relações de material, define-se analogamente a \textbf{condutividade}
$\sigma=1/\rho$.}
\end{questao}

\begin{questao}[1]{}
Segundo a segunda lei de Ohm, dobrar o \textbf{comprimento} de um fio, mantida
a seção, faz a resistência:
\begin{alternativas}[2]
\item dobrar
\item cair à metade
\item quadruplicar
\item não se alterar
\end{alternativas}
\resposta{(a)}
\resolucao{Como $R=\rho L/A$, tem-se $R\propto L$: dobrar $L$ dobra $R$.\\
\textbf{Interpretação física:} o dobro do comprimento equivale a dois fios
idênticos em \emph{série} --- e resistências em série se somam. A analogia
hidráulica ajuda: um cano duas vezes mais longo oferece o dobro de oposição ao
escoamento.}
\end{questao}

\begin{questao}[1]{}
Dobrar a \textbf{área} da seção de um fio, mantido o comprimento, faz a
resistência:
\begin{alternativas}[2]
\item dobrar
\item cair à metade
\item quadruplicar
\item não se alterar
\end{alternativas}
\resposta{(b)}
\resolucao{$R\propto1/A$: dobrar $A$ divide $R$ por dois.\\
\textbf{Interpretação:} o dobro da área equivale a dois fios idênticos em
\emph{paralelo}, e o paralelo de dois iguais vale metade.\\
\textbf{Cuidado com o diâmetro:} dobrar o \emph{diâmetro} quadruplica a área
($A=\pi d^{2}/4$) e divide $R$ por \textbf{quatro} --- confundir diâmetro com
área é o deslize mais comum aqui.}
\end{questao}

\begin{questao}[1]{}
Dois fios de mesmo material e mesmo comprimento têm diâmetros $d$ e $2d$.
Determine a razão entre suas resistências.
\resposta{$R_1/R_2=4$}
\resolucao{A área é proporcional ao quadrado do diâmetro:
\[
\frac{A_2}{A_1}=\left(\frac{2d}{d}\right)^{2}=4
\;\Longrightarrow\;
\frac{R_1}{R_2}=\frac{A_2}{A_1}=4.
\]
O fio mais grosso tem um quarto da resistência. É por isso que a bitola dos
condutores é especificada em $\si{\milli\meter\squared}$ (área) e não em
milímetros de diâmetro --- a área é o que entra linearmente na conta.}
\end{questao}

\blocoaplicacao

\begin{questao}[2]{}
Determine o comprimento de um fio de níquel-cromo
($\rho=\pot{1.1}{-6}\,\si{\ohm\meter}$) de seção
$\SI{0.2}{\milli\meter\squared}$ necessário para obter
$\SI{22}{\ohm}$.
\resposta{$L=\SI{4}{\meter}$}
\resolucao{Isolando $L$:
\[
L=\frac{RA}{\rho}=\frac{22\times\pot{0.2}{-6}}{\pot{1.1}{-6}}
=\frac{\pot{4.4}{-6}}{\pot{1.1}{-6}}=\SI{4}{\meter}.
\]
\textbf{Por que níquel-cromo:} sua resistividade é cerca de $65$ vezes a do
cobre, o que permite resistências úteis com comprimentos manejáveis. Com cobre,
os mesmos $\SI{22}{\ohm}$ exigiriam $\SI{259}{\meter}$ de fio --- inviável.}
\end{questao}

\begin{questao}[2]{}
Um cabo de alumínio ($\rho=\pot{2.8}{-8}\,\si{\ohm\meter}$) de
$\SI{200}{\meter}$ deve ter resistência de $\SI{0.5}{\ohm}$. Determine a seção
e o diâmetro.
\resposta{$A=\SI{11.2}{\milli\meter\squared}$; $d=\SI{3.78}{\milli\meter}$}
\resolucao{
\[
A=\frac{\rho L}{R}=\frac{\pot{2.8}{-8}\times200}{0{,}5}
=\pot{1.12}{-5}\,\si{\meter\squared}=\SI{11.2}{\milli\meter\squared}.
\]
\[
d=\sqrt{\frac{4A}{\pi}}=\sqrt{\frac{4\times\pot{1.12}{-5}}{\pi}}
=\SI{3.78}{\milli\meter}.
\]
\textbf{Valor comercial:} a bitola imediatamente superior é
$\SI{16}{\milli\meter\squared}$, que daria $R=\SI{0.35}{\ohm}$ --- com folga.
Adotar $\SI{10}{\milli\meter\squared}$ resultaria em $\SI{0.56}{\ohm}$,
acima do especificado.}
\end{questao}

\begin{questao}[2]{}
Mediu-se $R=\SI{2.5}{\ohm}$ em um fio de $\SI{30}{\meter}$ e seção
$\SI{0.5}{\milli\meter\squared}$. Determine a resistividade e identifique o
material provável.
\resposta{$\rho=\pot{4.2}{-8}\,\si{\ohm\meter}$ --- compatível com latão ou
zinco}
\resolucao{
\[
\rho=\frac{RA}{L}=\frac{2{,}5\times\pot{0.5}{-6}}{30}
=\pot{4.17}{-8}\,\si{\ohm\meter}.
\]
\textbf{Comparação:} cobre vale $\pot{1.7}{-8}$, alumínio
$\pot{2.8}{-8}$, zinco $\pot{5.9}{-8}$ e latão em torno de
$\pot{7}{-8}\,\si{\ohm\meter}$. O valor obtido está entre alumínio e zinco.\\
\textbf{Ressalva honesta:} a resistividade sozinha \emph{não} identifica o
material com segurança --- ligas de composições distintas podem coincidir. Ela
serve para descartar hipóteses (não é cobre puro, certamente) e como um entre
vários indícios.}
\end{questao}

\begin{questao}[2]{}
Um fio cilíndrico é \textbf{esticado} até dobrar seu comprimento, conservando o
volume e o material. Determine a nova resistência.
\resposta{$R'=4R$}
\resolucao{\textbf{Volume constante:} $V=AL=A'L'$. Com $L'=2L$:
\[
A'=\frac{AL}{2L}=\frac{A}{2}.
\]
\textbf{Nova resistência:}
\[
R'=\frac{\rho L'}{A'}=\frac{\rho(2L)}{A/2}=4\,\frac{\rho L}{A}=4R.
\]
\textbf{Dois efeitos que se multiplicam:} o comprimento dobra (fator $2$) e a
área cai à metade (outro fator $2$).\\
\textbf{Não confunda} com o caso em que $L$ e $A$ são alterados
\emph{independentemente} --- ali os fatores podem se cancelar. No estiramento,
eles se reforçam, porque o volume amarra as duas dimensões.}
\end{questao}

\begin{questao}[2]{}
Um resistor de fio tem $\SI{50}{\ohm}$ a $\SI{20}{\celsius}$, com
$\alpha=\SI{0.0039}{\per\celsius}$. Determine sua resistência a
$\SI{80}{\celsius}$.
\resposta{$R=\SI{61.7}{\ohm}$}
\resolucao{
\[
R=R_0\left[1+\alpha(T-T_0)\right]=50\left[1+0{,}0039(60)\right]
=50(1{,}234)=\SI{61.7}{\ohm}.
\]
Aumento de $23{,}4\%$ --- nada desprezível.\\
\textbf{Consequência para medição:} um ohmímetro que aplique corrente suficiente
para aquecer o componente lerá um valor \emph{maior} que o real a temperatura
ambiente. É por isso que medições de precisão especificam a corrente de teste e
a temperatura.}
\end{questao}

\begin{questao}[2]{}
Um estudante mediu, em um componente, os pares
$(\SI{2}{\volt};\SI{20}{\milli\ampere})$,
$(\SI{4}{\volt};\SI{40}{\milli\ampere})$ e
$(\SI{6}{\volt};\SI{57}{\milli\ampere})$. O componente é ôhmico?
\resposta{Não --- a razão $U/I$ cresce de $100$ para $\SI{105}{\ohm}$}
\resolucao{Calculando $R=U/I$ em cada ponto:
\[
\frac{2}{0{,}020}=100;
\qquad
\frac{4}{0{,}040}=100;
\qquad
\frac{6}{0{,}057}=\SI{105}{\ohm}.
\]
Os dois primeiros pontos são consistentes; o terceiro destoa em $5\%$.\\
\textbf{Diagnóstico provável: autoaquecimento.} Em $\SI{6}{\volt}$ a potência é
$\SI{0.34}{\watt}$, contra $\SI{0.04}{\watt}$ no primeiro ponto --- oito vezes
mais. Se o material tiver coeficiente positivo, a resistência sobe com a
temperatura.\\
\textbf{Como confirmar:} refaça o terceiro ponto \emph{rapidamente}, antes que o
componente aqueça. Se a leitura voltar a $\SI{100}{\ohm}$, o material é ôhmico
--- e o desvio era térmico, não uma propriedade elétrica intrínseca.}
\end{questao}

\begin{questao}[2]{}
Compare a seção necessária para obter a mesma resistência com cobre
($\rho=\pot{1.7}{-8}$) e com alumínio ($\pot{2.8}{-8}\,\si{\ohm\meter}$),
mantido o comprimento.
\resposta{$A_{Al}/A_{Cu}=1{,}65$}
\resolucao{Para $R$ e $L$ iguais, $A\propto\rho$:
\[
\frac{A_{Al}}{A_{Cu}}=\frac{\rho_{Al}}{\rho_{Cu}}=\frac{2{,}8}{1{,}7}=1{,}65.
\]
O alumínio exige $65\%$ mais seção.\\
\textbf{Por que ainda assim se usa alumínio:} sua densidade é
$\SI{2.7}{\gram\per\centi\meter\cubed}$ contra $\SI{8.9}{}$ do cobre. Para a
mesma resistência, o cabo de alumínio pesa
$1{,}65\times(2{,}7/8{,}9)=0{,}50$ --- \textbf{metade} do peso, e custa bem
menos. É por isso que linhas de transmissão aéreas são de alumínio, onde o peso
é crítico, enquanto instalações prediais são de cobre, onde o espaço nos
eletrodutos é que limita.}
\end{questao}

\begin{questao}[2]{}
Um fio de cobre de $\SI{100}{\meter}$ e $\SI{1.5}{\milli\meter\squared}$
conduz $\SI{10}{\ampere}$. Determine a resistência, a queda de tensão e a
potência dissipada.
\resposta{$R=\SI{1.13}{\ohm}$; $\Delta U=\SI{11.3}{\volt}$;
$P=\SI{113}{\watt}$}
\resolucao{
\[
R=\frac{\pot{1.7}{-8}\times100}{\pot{1.5}{-6}}=\SI{1.13}{\ohm};
\qquad
\Delta U=RI=\SI{11.3}{\volt};
\qquad
P=RI^{2}=\SI{113}{\watt}.
\]
\textbf{Avaliação:} em uma rede de $\SI{220}{\volt}$, essa queda representa
$5{,}2\%$ --- acima do limite usual de $4\%$ para circuitos terminais. E
$\SI{113}{\watt}$ dissipados ao longo do cabo é aquecimento real, que exige
verificação de capacidade de condução.\\
\textbf{Conclusão:} a bitola está subdimensionada para esse comprimento. O
critério de queda de tensão, e não o de corrente admissível, é que domina em
circuitos longos.}
\end{questao}

\blocoaprofundamento

\begin{questao}[3]{}
Projete um resistor de fio de $\SI{10}{\ohm}$ usando material de
$\rho=\pot{5}{-7}\,\si{\ohm\meter}$ e fio de $\SI{0.5}{\milli\meter}$ de
diâmetro.
\begin{itensq}
\item Determine o comprimento necessário.
\item Determine a corrente máxima para $\SI{5}{\watt}$ de dissipação.
\item Comente a construção.
\end{itensq}
\resposta{(a) $L=\SI{3.93}{\meter}$; (b) $I=\SI{0.707}{\ampere}$}
\resolucao{(a) Seção:
$A=\dfrac{\pi d^{2}}{4}=\dfrac{\pi(\pot{5}{-4})^{2}}{4}
=\pot{1.963}{-7}\,\si{\meter\squared}$.
\[
L=\frac{RA}{\rho}=\frac{10\times\pot{1.963}{-7}}{\pot{5}{-7}}
=\SI{3.93}{\meter}.
\]
(b) $I=\sqrt{P/R}=\sqrt{5/10}=\SI{0.707}{\ampere}$, com
$U=\SI{7.07}{\volt}$.\\
(c) \textbf{Construção.} Quase quatro metros de fio precisam ser enrolados
sobre um corpo cerâmico. Três cuidados:
\begin{itemize}
\item \textbf{Espaçamento entre espiras}, para evitar curtos e permitir
circulação de ar.
\item \textbf{Área de dissipação suficiente} --- $\SI{5}{\watt}$ exigem corpo de
alguns centímetros, e a temperatura de superfície pode passar de
$\SI{100}{\celsius}$.
\item \textbf{Indutância parasita}: um enrolamento simples de $\SI{3.93}{\meter}$
comporta-se como bobina. Em corrente contínua isso é irrelevante, mas em sinais
rápidos é preciso enrolamento \emph{bifilar}, em que a ida e a volta se
cancelam.
\end{itemize}}
\end{questao}

\begin{questao}[3]{}
Um resistor nominal de $\SI{100}{\ohm}$ tem tolerância de $\SI{5}{\percent}$ e
coeficiente térmico de $400\,\mathrm{ppm/^{\circ}C}$, especificado a
$\SI{20}{\celsius}$. Determine a faixa total de resistência na operação entre
$\SI{0}{\celsius}$ e $\SI{70}{\celsius}$.
\resposta{de $\SI{94.2}{\ohm}$ a $\SI{107.1}{\ohm}$ ($-5{,}8\%$ a
$+7{,}1\%$)}
\resolucao{\textbf{Deriva térmica:} $400\,\mathrm{ppm}=\pot{4}{-4}$ por
$\si{\celsius}$.
\[
\SI{70}{\celsius}:\ \Delta T=+50 \Rightarrow +2{,}0\%;
\qquad
\SI{0}{\celsius}:\ \Delta T=-20 \Rightarrow -0{,}8\%.
\]
\textbf{Combinação de pior caso} (os dois efeitos no mesmo sentido):
\[
R_{\max}=100(1{,}05)(1{,}020)=\SI{107.1}{\ohm};
\qquad
R_{\min}=100(0{,}95)(0{,}992)=\SI{94.2}{\ohm}.
\]
\textbf{Observações de projeto.}
\begin{itemize}
\item A tolerância de fabricação domina --- ela responde por $5$ dos
$7{,}1$ pontos.
\item A faixa é \emph{assimétrica}, porque a temperatura de operação se afasta
mais para cima que para baixo do valor de referência.
\item Se $7\%$ for demais, apertar a tolerância para $1\%$ traria a faixa a
$-1{,}8\%$/$+3{,}0\%$; trocar o coeficiente para
$25\,\mathrm{ppm}$ traria a $-5{,}0\%$/$+5{,}1\%$. \textbf{Ataque primeiro a
maior parcela.}
\end{itemize}}
\end{questao}

\begin{questao}[3]{}
Uma lâmpada incandescente apresenta $\SI{2.00}{\ampere}$ a
$\SI{12.0}{\volt}$ e $\SI{2.02}{\ampere}$ a $\SI{12.5}{\volt}$.
\begin{itensq}
\item Determine a resistência estática no primeiro ponto.
\item Determine a resistência dinâmica.
\item Explique a diferença.
\end{itensq}
\resposta{$R_{est}=\SI{6}{\ohm}$; $r_{din}=\SI{25}{\ohm}$}
\resolucao{(a) $R_{est}=\dfrac{U}{I}=\dfrac{12{,}0}{2{,}00}=\SI{6}{\ohm}$.\\
(b) $r_{din}=\dfrac{\Delta U}{\Delta I}=\dfrac{0{,}5}{0{,}02}=\SI{25}{\ohm}$.\\
(c) \textbf{A dinâmica é mais de quatro vezes maior.} A razão é o
\textbf{autoaquecimento}: aumentar a tensão eleva a potência, o filamento
esquenta e sua resistência sobe --- de modo que a corrente cresce muito menos
que proporcionalmente.\\
\textbf{Cada grandeza responde a uma pergunta diferente:}
\begin{itemize}
\item $R_{est}$ governa o \emph{balanço de potência} no ponto de operação:
$P=U^{2}/R_{est}=\SI{24}{\watt}$ \checkmark.
\item $r_{din}$ governa a \emph{resposta a pequenas variações} --- é ela que
entra em qualquer modelo incremental.
\end{itemize}
\textbf{A lâmpada não é ôhmica em regime}, mas é aproximadamente ôhmica a
temperatura constante. É o acoplamento térmico, e não a física da condução, que
quebra a linearidade.}
\end{questao}

\begin{questao}[3]{}
Um cabo de cobre deve alimentar uma carga de $\SI{20}{\ampere}$ a
$\SI{220}{\volt}$, a $\SI{40}{\meter}$ de distância, com queda de tensão
inferior a $\SI{3}{\percent}$.
\begin{itensq}
\item Determine a resistência máxima admissível do cabo.
\item Determine a seção mínima.
\item Escolha a bitola comercial.
\end{itensq}
\resposta{$R\le\SI{0.33}{\ohm}$; $A\ge\SI{4.12}{\milli\meter\squared}$;
adotar $\SI{6}{\milli\meter\squared}$}
\resolucao{(a) Queda admissível: $0{,}03\times220=\SI{6.6}{\volt}$.
\[
R_{\max}=\frac{6{,}6}{20}=\SI{0.33}{\ohm}.
\]
(b) \textbf{Atenção ao comprimento:} a corrente percorre \emph{ida e volta},
logo $L=2\times40=\SI{80}{\meter}$.
\[
A\ge\frac{\rho L}{R_{\max}}=\frac{\pot{1.7}{-8}\times80}{0{,}33}
=\pot{4.12}{-6}\,\si{\meter\squared}=\SI{4.12}{\milli\meter\squared}.
\]
(c) A bitola comercial imediatamente superior é
$\SI{6}{\milli\meter\squared}$, que dá $R=\SI{0.227}{\ohm}$ e queda de
$\SI{4.5}{\volt}$ ($2{,}1\%$) --- dentro da especificação com folga.\\
\textbf{Verificação obrigatória do outro critério:} $\SI{6}{\milli\meter\squared}$
suporta cerca de $\SI{41}{\ampere}$ em eletroduto, bem acima dos
$\SI{20}{\ampere}$ da carga. \textbf{Aqui a queda de tensão é o critério
dominante} --- e é o que costuma ocorrer em circuitos longos. Em circuitos
curtos, a corrente admissível é que decide.}
\end{questao}

\begin{questao}[3]{}
Um fio é esticado até que seu comprimento fique $n$ vezes maior, conservando o
volume.
\begin{itensq}
\item Deduza a nova resistência em função de $n$.
\item Avalie para $n=3$ e $n=5$.
\item Determine a redução do diâmetro em cada caso.
\end{itensq}
\resposta{$R'=n^{2}R$; $9R$ e $25R$; $d'/d=1/\sqrt{n}$}
\resolucao{(a) Volume constante: $AL=A'L'$ com $L'=nL$, logo
$A'=A/n$.
\[
R'=\frac{\rho(nL)}{A/n}=n^{2}\frac{\rho L}{A}=n^{2}R.
\]
(b) $n=3\Rightarrow9R$; $n=5\Rightarrow25R$.\\
(c) Como $A\propto d^{2}$ e $A'=A/n$:
\[
\frac{d'}{d}=\frac{1}{\sqrt{n}}
\Rightarrow
n=3:\ 0{,}577;
\qquad
n=5:\ 0{,}447.
\]
\textbf{Aplicação prática --- o extensômetro.} Um fio colado a uma peça que se
deforma sofre exatamente esse efeito. Para pequenas deformações
$\varepsilon=\Delta L/L$, a expansão de $n^{2}$ dá
\[
\frac{\Delta R}{R}\approx2\varepsilon,
\]
ou seja, o fio funciona como sensor de deformação com \emph{fator de
sensibilidade} próximo de $2$. É o princípio das células de carga e das balanças
eletrônicas.}
\end{questao}

\begin{questao}[3]{}
Uma barra condutora de comprimento $L=\SI{1}{\meter}$ tem seção que varia
linearmente segundo $A(x)=A_0\left(1+x/L\right)$, com
$A_0=\SI{1}{\milli\meter\squared}$ e $\rho=\pot{1.7}{-8}\,\si{\ohm\meter}$.
\begin{itensq}
\item Determine a resistência total.
\item Compare com uma barra uniforme de seção $A_0$ e com uma de seção média.
\end{itensq}
\resposta{$R=\dfrac{\rho L\ln2}{A_0}=\SI{11.8}{\milli\ohm}$}
\resolucao{(a) Somando as resistências de fatias infinitesimais em série:
\[
R=\int_0^{L}\frac{\rho\,\mathrm{d}x}{A(x)}
=\frac{\rho}{A_0}\int_0^{L}\frac{\mathrm{d}x}{1+x/L}
=\frac{\rho L}{A_0}\Big[\ln(1+x/L)\Big]_0^{L}
=\frac{\rho L\ln2}{A_0}.
\]
\[
R=\frac{\pot{1.7}{-8}\times1\times0{,}693}{\pot{1}{-6}}
=\SI{11.8}{\milli\ohm}.
\]
(b)
\begin{center}
\begin{tabular}{lc}
\hline
Modelo & $R$\\
\hline
uniforme com $A_0$ (menor seção) & $\SI{17.0}{\milli\ohm}$\\
\textbf{seção variável (exato)} & $\SI{11.8}{\milli\ohm}$\\
uniforme com $\bar A=1{,}5A_0$ & $\SI{11.3}{\milli\ohm}$\\
\hline
\end{tabular}
\end{center}
\textbf{O valor exato é maior que o da seção média} --- porque a resistência
depende de $1/A$, e a média de $1/A$ é maior que $1/\bar A$ (desigualdade das
médias). O erro de usar a seção média é de $4{,}3\%$, para menos.\\
\textbf{Regra:} em seções variáveis, integrar $\rho\,\mathrm{d}x/A(x)$. Usar a
seção média subestima; usar a seção mínima superestima. As duas aproximações
servem como limites.}
\end{questao}

\begin{questao}[3]{}
Um sensor Pt100 tem $R_0=\SI{100}{\ohm}$ a $\SI{0}{\celsius}$ e
$\alpha=\SI{0.00385}{\per\celsius}$, alimentado por corrente constante de
$\SI{1}{\milli\ampere}$.
\begin{itensq}
\item Determine a tensão a $\SI{0}{\celsius}$ e a sensibilidade em
      $\si{\milli\volt\per\celsius}$.
\item Explique a vantagem da alimentação por corrente constante.
\item Verifique o autoaquecimento.
\end{itensq}
\resposta{$U_0=\SI{100}{\milli\volt}$;
$\SI{0.385}{\milli\volt\per\celsius}$}
\resolucao{(a) $U_0=IR_0=\pot{1}{-3}(100)=\SI{100}{\milli\volt}$.
\[
\frac{\mathrm{d}U}{\mathrm{d}T}=I\,R_0\,\alpha
=\pot{1}{-3}(100)(0{,}00385)=\SI{0.385}{\milli\volt\per\celsius}.
\]
(b) \textbf{Vantagem decisiva: linearidade exata.} Com $I$ constante,
$U=IR$ é \emph{proporcional} a $R$, e como $R$ é linear em $T$, a tensão também
é. Não há distorção alguma.\\
Com \emph{tensão} constante em um divisor, a saída seria
$U\,R/(R+R_{fixo})$ --- uma função \textbf{não linear} de $R$, exigindo
linearização posterior.\\
(c) \textbf{Autoaquecimento:} $P=RI^{2}=100(\pot{1}{-6})=\SI{0.1}{\milli\watt}$.
Com resistência térmica típica de $\SI{100}{\celsius\per\watt}$ ao ar, a
elevação é de $\SI{0.01}{\celsius}$ --- desprezível.\\
\textbf{Mas cuidado:} elevar a corrente para $\SI{10}{\milli\ampere}$ (para
ganhar sinal) multiplicaria a potência por $100$, levando o erro a
$\SI{1}{\celsius}$. \textbf{A corrente de excitação é um compromisso entre
relação sinal-ruído e erro de autoaquecimento} --- e $\SI{1}{\milli\ampere}$ é
o valor de consenso para Pt100.}
\end{questao}

\begin{questao}[3]{}
Mede-se um resistor de $\SI{0.1}{\ohm}$ com um ohmímetro cujos cabos têm
$\SI{0.2}{\ohm}$ cada.
\begin{itensq}
\item Determine a leitura pelo método de \textbf{dois fios}.
\item Determine o erro.
\item Explique como o método de \textbf{quatro fios} o elimina.
\end{itensq}
\resposta{Leitura de $\SI{0.5}{\ohm}$ --- erro de $400\%$}
\resolucao{(a) O instrumento mede tudo o que está em série:
\[
R_{lido}=0{,}1+2(0{,}2)=\SI{0.5}{\ohm}.
\]
(b) Erro de $\SI{0.4}{\ohm}$, ou $400\%$ --- a leitura é \emph{cinco vezes} o
valor real. Para resistências pequenas, o método de dois fios é simplesmente
inutilizável.\\
(c) \textbf{Método de quatro fios (Kelvin).} Usam-se dois pares independentes:
\begin{itemize}
\item Um par \textbf{injeta} a corrente conhecida $I$. A queda nesses cabos
existe, mas é irrelevante --- ela apenas exige um pouco mais de tensão da fonte.
\item O outro par \textbf{mede} a tensão diretamente sobre o resistor, com um
voltímetro de altíssima impedância. Como praticamente não circula corrente
nesses cabos, sua resistência \emph{não produz queda}.
\end{itemize}
Assim $R=U_{medido}/I$ devolve exatamente $\SI{0.1}{\ohm}$.\\
\textbf{Quando usar:} sempre que $R$ for comparável ou menor que a resistência
dos cabos --- tipicamente abaixo de alguns ohms. É o método obrigatório para
shunts, enrolamentos de motor e contatos.}
\end{questao}

\begin{questao}[3]{}
Dois resistores têm o mesmo valor nominal de $\SI{100}{\ohm}$, mas potências de
$\frac14\,\si{\watt}$ e $\SI{5}{\watt}$.
\begin{itensq}
\item Determine a tensão e a corrente máximas de cada um.
\item Ambos obedecem igualmente à lei de Ohm?
\item Explique o que a especificação de potência realmente limita.
\end{itensq}
\resposta{$\SI{5}{\volt}$/$\SI{50}{\milli\ampere}$ e
$\SI{22.4}{\volt}$/$\SI{224}{\milli\ampere}$; ambos são ôhmicos}
\resolucao{(a)
\[
U_{\max}=\sqrt{PR}:\quad \sqrt{0{,}25(100)}=\SI{5}{\volt};
\qquad \sqrt{5(100)}=\SI{22.4}{\volt};
\]
\[
I_{\max}=\sqrt{P/R}:\quad \SI{50}{\milli\ampere}
\quad\text{e}\quad \SI{224}{\milli\ampere}.
\]
(b) \textbf{Sim, ambos obedecem igualmente.} A lei de Ohm relaciona $U$ e $I$
através de $R$ --- e $R$ é o mesmo. A potência nominal \emph{não} aparece na
lei.\\
(c) \textbf{A potência nominal limita a faixa de validade}, não a relação em
si. Ela informa até onde o componente pode operar \emph{antes de se
autodestruir ou sair da linearidade por aquecimento}.\\
\textbf{Analogia útil:} a lei de Ohm é como a relação entre força e deformação
em uma mola; a potência nominal é o limite de escoamento do material. Duas
molas de mesma constante e materiais diferentes obedecem à mesma lei --- até que
uma delas se rompa.\\
\textbf{Diferença física:} o resistor de $\SI{5}{\watt}$ é maior, com mais área
de troca térmica. Mesma resistência, mesma lei, capacidade de dissipação
$20$ vezes maior.}
\end{questao}

\begin{questao}[3]{}
Um resistor de $\SI{100}{\ohm}$ com $\alpha=\SI{+0.004}{\per\celsius}$ dissipa
$\SI{1}{\watt}$ e tem resistência térmica de
$\SI{100}{\celsius\per\watt}$.
\begin{itensq}
\item Determine a elevação de temperatura e a nova resistência.
\item Analise o efeito sob tensão constante e sob corrente constante.
\end{itensq}
\resposta{$\Delta T=\SI{100}{\celsius}$; $R=\SI{140}{\ohm}$}
\resolucao{(a)
\[
\Delta T=R_{\text{térm}}P=100(1)=\SI{100}{\celsius};
\qquad
R=100\left[1+0{,}004(100)\right]=\SI{140}{\ohm}.
\]
Aumento de $40\%$ apenas por autoaquecimento --- muito acima da tolerância
típica de fabricação.\\
(b) \textbf{Sob tensão constante:} $P=U^{2}/R$ \emph{diminui} quando $R$ sobe.
Menos potência, menos aquecimento, $R$ estabiliza. \textbf{Realimentação
negativa --- estável.}\\
\textbf{Sob corrente constante:} $P=RI^{2}$ \emph{aumenta} quando $R$ sobe.
Mais potência, mais aquecimento, $R$ cresce ainda mais.
\textbf{Realimentação positiva.}\\
\textbf{Condição de estabilidade} no caso de corrente constante: o incremento de
potência por grau, $\alpha R I^{2}$, deve ser menor que a capacidade de
escoar calor, $1/R_{\text{térm}}$:
\[
\alpha R I^{2}R_{\text{térm}}<1.
\]
Aqui, $0{,}004(100)(0{,}01)(100)=0{,}4<1$ \checkmark --- estável, mas com apenas
$2{,}5$ vezes de margem. Dobrar a corrente levaria o produto a $1{,}6$ e
provocaria fuga térmica.\\
\textbf{Materiais com $\alpha<0$} (termistores NTC, carbono) invertem os dois
casos --- e são instáveis justamente sob \emph{tensão} constante.}
\end{questao}

\begin{questao}[3]{}
Um resistor de $\SI{100}{\ohm}$ e $\frac14\,\si{\watt}$ deve ser testado quanto
à ohmicidade. Projete o ensaio.
\begin{itensq}
\item Determine a faixa segura de tensão.
\item Proponha os pontos de medição e o critério de decisão.
\item Indique os cuidados experimentais.
\end{itensq}
\resposta{Até $\SI{5}{\volt}$; usar $20\%$ da potência nominal como teto
prático}
\resolucao{(a) $U_{\max}=\sqrt{PR}=\SI{5}{\volt}$ no limite absoluto. Para
evitar autoaquecimento --- que é justamente o que falsearia o ensaio --- adote
$20\%$ da potência nominal:
\[
P_{ensaio}=\SI{50}{\milli\watt}
\;\Longrightarrow\;
U_{\max}=\sqrt{0{,}05(100)}=\SI{2.24}{\volt}.
\]
(b) \textbf{Pontos:} cinco valores igualmente espaçados entre $0{,}5$ e
$\SI{2.2}{\volt}$, medindo $U$ e $I$ em cada um.\\
\textbf{Critério:} calcular $R=U/I$ em todos os pontos. O componente é ôhmico se
a dispersão dos valores ficar dentro da incerteza combinada dos instrumentos
(tipicamente $1$ a $2\%$). Alternativamente, ajustar uma reta a $U\times I$ e
verificar se o intercepto é nulo dentro da incerteza --- intercepto não nulo
indica tensão de limiar, como em diodos.\\
(c) \textbf{Cuidados.}
\begin{itemize}
\item \textbf{Medir rápido} e desligar entre pontos, para o componente não
aquecer cumulativamente.
\item \textbf{Ordem aleatória} dos pontos, e não crescente --- assim um eventual
aquecimento progressivo aparece como dispersão, e não como tendência que
poderia ser confundida com não linearidade real.
\item \textbf{Método de quatro fios} se $R$ for pequeno.
\item \textbf{Repetir o primeiro ponto ao final:} se o valor mudou, houve
aquecimento e o ensaio precisa ser refeito com potência menor.
\end{itemize}}
\end{questao}

\blocodesafio

\begin{questao}[4]{}
\textbf{(Demonstração --- deformação e resistência.)} Um fio de comprimento $L$
e seção $A$ sofre pequena deformação axial $\varepsilon=\Delta L/L$, a volume
constante.
\begin{itensq}
\item Deduza $\Delta R/R$ em primeira ordem.
\item Determine o fator de sensibilidade.
\item Discuta o efeito de a resistividade também variar com a deformação.
\end{itensq}
\resposta{$\Delta R/R\approx2\varepsilon$; fator $\approx2$}
\resolucao{(a) Com volume constante, $L'=L(1+\varepsilon)$ e
$A'=A/(1+\varepsilon)$:
\[
R'=\frac{\rho L(1+\varepsilon)}{A/(1+\varepsilon)}
=R(1+\varepsilon)^{2}\approx R(1+2\varepsilon),
\]
para $\varepsilon\ll1$. Logo
\[
\frac{\Delta R}{R}\approx2\varepsilon.
\]
(b) O \textbf{fator de sensibilidade} (\emph{gauge factor}) é
$k=\dfrac{\Delta R/R}{\varepsilon}\approx2$.\\
\textbf{Ordem de grandeza prática:} uma deformação de
$\SI{1000}{\micro\meter\per\meter}$ ($\varepsilon=\pot{1}{-3}$) produz
$\Delta R/R=0{,}2\%$ --- em um extensômetro de $\SI{350}{\ohm}$, apenas
$\SI{0.7}{\ohm}$. Daí a necessidade de ponte de Wheatstone e amplificação de
instrumentação.\\
(c) \textbf{Efeito piezorresistivo.} Em metais, $\rho$ varia pouco com a
deformação e $k\approx2$ decorre quase só da geometria. Em
\textbf{semicondutores}, a deformação altera fortemente a mobilidade dos
portadores, e $k$ chega a $100$ ou mais --- ganho enorme de sensibilidade.\\
\textbf{Contrapartida:} extensômetros semicondutores são muito mais sensíveis à
temperatura e menos lineares. \textbf{A escolha entre metal e semicondutor é um
compromisso entre sensibilidade e estabilidade} --- e por isso células de carga
de precisão ainda usam extensômetros metálicos.}
\end{questao}

\begin{questao}[4]{}
\textbf{(Integração --- seção variável.)} Uma barra tem seção
$A(x)=A_0\left(1+x/L\right)$.
\begin{itensq}
\item Deduza $R$ por integração.
\item Generalize para $A(x)=A_0(1+kx/L)$.
\item Analise os limites $k\to0$ e $k\to\infty$.
\end{itensq}
\resposta{$R=\dfrac{\rho L}{A_0}\cdot\dfrac{\ln(1+k)}{k}$; para $k=1$,
$\ln2$}
\resolucao{(a) Fatias de espessura $\mathrm{d}x$ estão em série:
\[
R=\int_0^{L}\frac{\rho\,\mathrm{d}x}{A_0(1+x/L)}
=\frac{\rho L}{A_0}\ln 2.
\]
(b) Com $A(x)=A_0(1+kx/L)$:
\[
R=\frac{\rho}{A_0}\int_0^{L}\frac{\mathrm{d}x}{1+kx/L}
=\frac{\rho L}{A_0k}\Big[\ln(1+kx/L)\Big]_0^{L}
=\frac{\rho L}{A_0}\cdot\frac{\ln(1+k)}{k}.
\]
(c) \textbf{Limite $k\to0$} (seção uniforme). Usando
$\ln(1+k)\approx k-k^{2}/2$:
\[
\frac{\ln(1+k)}{k}\to1
\;\Longrightarrow\;
R\to\frac{\rho L}{A_0}.\ \checkmark
\]
Recupera-se a fórmula elementar --- teste obrigatório de qualquer generalização.\\
\textbf{Limite $k\to\infty$} (seção crescendo muito rápido):
\[
\frac{\ln(1+k)}{k}\to0
\;\Longrightarrow\;
R\to0.
\]
A barra fica tão larga na maior parte de seu comprimento que sua resistência se
concentra apenas na região inicial estreita --- e o total tende a zero, embora
\emph{logaritmicamente devagar}. Mesmo com $k=1000$, o fator ainda vale
$0{,}0069$, e não zero.\\
\textbf{Leitura geral:} a resistência é dominada pela região de \emph{menor}
seção. É por isso que um estrangulamento local em um condutor concentra
praticamente toda a queda de tensão --- e todo o aquecimento.}
\end{questao}

\begin{questao}[4]{}
\textbf{(Ponto de operação com fonte não ideal.)} Um resistor de
$\SI{100}{\ohm}$ é ligado a uma fonte cuja característica é
$U=24-8I$ (com $U$ em volts e $I$ em ampères).
\begin{itensq}
\item Determine o ponto de operação.
\item Determine a potência no resistor e o rendimento.
\item Repita para um resistor de $\SI{8}{\ohm}$ e compare.
\end{itensq}
\resposta{$I=\SI{0.222}{\ampere}$, $U=\SI{22.2}{\volt}$;
$P=\SI{4.94}{\watt}$, $\eta=92{,}6\%$}
\resolucao{(a) A fonte impõe $U=24-8I$ e o resistor impõe $U=100I$. Igualando:
\[
100I=24-8I
\;\Longrightarrow\;
108I=24
\;\Longrightarrow\;
I=\SI{0.222}{\ampere};
\qquad
U=\SI{22.2}{\volt}.
\]
(b) $P_L=UI=22{,}2(0{,}222)=\SI{4.94}{\watt}$.\\
Rendimento: $\eta=\dfrac{R_L}{R_L+R_{Th}}=\dfrac{100}{108}=92{,}6\%$.\\
(c) \textbf{Com $\SI{8}{\ohm}$:} $8I=24-8I\Rightarrow I=\SI{1.5}{\ampere}$ e
$U=\SI{12}{\volt}$.
\[
P_L=12(1{,}5)=\SI{18}{\watt};
\qquad
\eta=\frac{8}{16}=50\%.
\]
\textbf{Comparação reveladora.} O resistor de $\SI{8}{\ohm}$ está \emph{casado}
com a fonte e recebe a potência máxima ($\SI{18}{\watt}$, quase quatro vezes
mais), mas desperdiça metade da energia internamente. O de
$\SI{100}{\ohm}$ recebe pouco, porém com rendimento de $92{,}6\%$.\\
\textbf{Ponto de método:} o problema é a interseção de duas características
$U\times I$ --- uma reta de fonte e uma reta de carga. Se a carga fosse não
linear (um diodo, por exemplo), o procedimento seria idêntico: basta trocar a
segunda reta pela curva do dispositivo.}
\end{questao}

\begin{questao}[4]{}
\textbf{(Estabilidade térmica.)} Um resistor de resistência $R$, coeficiente
$\alpha$ e resistência térmica $R_{\text{térm}}$ é alimentado por corrente constante
$I$.
\begin{itensq}
\item Deduza a condição de estabilidade térmica.
\item Avalie para $R=\SI{100}{\ohm}$, $\alpha=\SI{0.004}{\per\celsius}$,
      $R_{\text{térm}}=\SI{100}{\celsius\per\watt}$ e $I=\SI{100}{\milli\ampere}$.
\item Determine a corrente crítica.
\end{itensq}
\resposta{Condição: $\alpha R I^{2}R_{\text{térm}}<1$; aqui $0{,}4$ --- estável;
$I_c=\SI{158}{\milli\ampere}$}
\resolucao{(a) Em equilíbrio, $\Delta T=R_{\text{térm}}P=R_{\text{térm}}RI^{2}$ e
$R=R_0(1+\alpha\Delta T)$. Substituindo,
\[
R=R_0\left(1+\alpha R_{\text{térm}}RI^{2}\right)
\;\Longrightarrow\;
R\left(1-\alpha R_{\text{térm}}R_0I^{2}\right)=R_0.
\]
Existe solução finita e estável se, e somente se, o parênteses for positivo:
\[
\boxed{\alpha\,R_0\,I^{2}\,R_{\text{térm}}<1}
\]
Fisicamente: o acréscimo de potência por grau de aquecimento deve ser
\emph{menor} que a capacidade de escoar calor. Caso contrário, cada grau
adicional gera mais calor do que consegue dissipar --- \textbf{fuga térmica}.\\
(b)
\[
0{,}004\times100\times(0{,}1)^{2}\times100=0{,}4<1\ \checkmark
\]
Estável, com margem de $2{,}5$ vezes. A resistência de equilíbrio é
$R=100/(1-0{,}4)=\SI{167}{\ohm}$, e a elevação, $\SI{167}{\celsius}$ ---
elevada, mas estável.\\
(c) Corrente crítica:
\[
I_c=\sqrt{\frac{1}{\alpha R_0R_{\text{térm}}}}=\sqrt{\frac{1}{0{,}004(100)(100)}}
=\sqrt{0{,}025}=\SI{158}{\milli\ampere}.
\]
\textbf{Margem estreita:} apenas $58\%$ acima da corrente de operação. E note
que a estabilidade \emph{matemática} não garante a \emph{sobrevivência} --- a
$\SI{167}{\celsius}$ o componente já pode estar além de sua temperatura máxima.\\
\textbf{Sob tensão constante} o problema não existe: $P=U^{2}/R$ decresce com
$R$, e a realimentação é sempre negativa.}
\end{questao}

\begin{questao}[4]{}
\textbf{(Dois critérios de dimensionamento.)} Um cabo deve alimentar
$\SI{30}{\ampere}$ a $\SI{220}{\volt}$, com queda máxima de
$\SI{4}{\percent}$.
\begin{itensq}
\item Determine a seção exigida pelo critério de queda de tensão, para
      $\SI{20}{\meter}$ e para $\SI{80}{\meter}$.
\item Sabendo que $\SI{6}{\milli\meter\squared}$ suporta cerca de
      $\SI{41}{\ampere}$, determine qual critério domina em cada caso.
\item Formule a regra geral.
\end{itensq}
\resposta{$\SI{3.87}{\milli\meter\squared}$ e
$\SI{15.5}{\milli\meter\squared}$; a corrente domina no cabo curto, a queda no
longo}
\resolucao{(a) Queda admissível: $0{,}04(220)=\SI{8.8}{\volt}$, logo
$R_{\max}=8{,}8/30=\SI{0.293}{\ohm}$.
\[
\SI{20}{\meter}:\ A\ge\frac{\pot{1.7}{-8}(40)}{0{,}293}
=\SI{2.32}{\milli\meter\squared};
\]
\[
\SI{80}{\meter}:\ A\ge\frac{\pot{1.7}{-8}(160)}{0{,}293}
=\SI{9.28}{\milli\meter\squared}.
\]
(lembrando o fator $2$ de ida e volta).\\
(b) \textbf{Cabo de $\SI{20}{\meter}$:} a queda exige
$\SI{2.32}{\milli\meter\squared}$, mas a corrente de $\SI{30}{\ampere}$ exige
pelo menos $\SI{6}{\milli\meter\squared}$. \textbf{O critério de corrente
domina.}\\
\textbf{Cabo de $\SI{80}{\meter}$:} a queda exige
$\SI{9.28}{\milli\meter\squared}$, acima do que a corrente pediria.
\textbf{O critério de queda domina}, e adota-se
$\SI{10}{\milli\meter\squared}$.\\
(c) \textbf{Regra geral.} Os dois critérios escalam de formas diferentes:
\begin{itemize}
\item \textbf{Corrente admissível:} depende só de $I$ --- \emph{independe do
comprimento}. É um critério de segurança térmica do próprio cabo.
\item \textbf{Queda de tensão:} $A\propto IL/\Delta U$ --- cresce
\emph{linearmente} com o comprimento. É um critério de qualidade da energia
entregue.
\end{itemize}
Existe um \textbf{comprimento de cruzamento} em que os dois coincidem; abaixo
dele manda a corrente, acima manda a queda. Aqui, com
$\SI{6}{\milli\meter\squared}$, o cruzamento ocorre por volta de
$\SI{52}{\meter}$.\\
\textbf{Procedimento correto:} calcular os \emph{dois} e adotar a maior seção.
Verificar apenas um deles é a origem mais comum de instalações que "funcionam"
mas apresentam tensão baixa na ponta.}
\end{questao}

\begin{questao}[4]{}
\textbf{(Sensor a corrente constante.)} Compare duas formas de converter a
variação de um sensor resistivo em tensão.
\begin{itensq}
\item Obtenha $U(R)$ para excitação por corrente constante e para divisor de
      tensão.
\item Compare a linearidade das duas.
\item Determine a sensibilidade de cada uma no ponto de operação
      $R=R_{fixo}$.
\end{itensq}
\resposta{Corrente constante: $U=IR$, exatamente linear; divisor:
$U=U_0R/(R+R_f)$, não linear}
\resolucao{(a) \textbf{Corrente constante:} $U=IR$ --- proporcional a $R$.\\
\textbf{Divisor:} $U=U_0\dfrac{R}{R+R_f}$ --- função racional de $R$.\\
(b) \textbf{Linearidade.} A primeira é \emph{exatamente} linear, sempre. A
segunda só é aproximadamente linear para pequenas variações; expandindo em torno
de $R=R_f$:
\[
U\approx\frac{U_0}{2}\left[1+\frac{\Delta R}{2R_f}
-\frac{1}{4}\left(\frac{\Delta R}{R_f}\right)^{2}+\cdots\right].
\]
O termo quadrático produz erro de linearidade. Para
$\Delta R/R_f=10\%$, ele vale $0{,}25\%$ da leitura --- pequeno, mas suficiente
para limitar a exatidão em medições finas.\\
(c) \textbf{Sensibilidades.}
\[
\text{corrente constante:}\quad \frac{\mathrm{d}U}{\mathrm{d}R}=I;
\qquad
\text{divisor:}\quad \frac{\mathrm{d}U}{\mathrm{d}R}
=\frac{U_0R_f}{(R+R_f)^{2}}\bigg|_{R=R_f}=\frac{U_0}{4R_f}.
\]
Ambas podem ser ajustadas ao mesmo valor escolhendo $I=U_0/(4R_f)$ --- a
sensibilidade não distingue os métodos no ponto de operação.\\
\textbf{O que decide é o resto.}
\begin{itemize}
\item \textbf{Corrente constante:} linearidade exata, mas exige fonte de
corrente estável e o sinal parte de zero apenas se $R\to0$.
\item \textbf{Divisor:} circuito trivial e barato, e a saída pode ser centrada
em meia alimentação --- conveniente para conversores A/D. Em contrapartida,
carrega o erro de linearidade e a deriva de $R_f$ e de $U_0$.
\end{itemize}
\textbf{Solução de compromisso muito usada:} a \emph{ponte} de Wheatstone, que
cancela a componente comum e entrega diretamente o desequilíbrio --- combinando
boa linearidade com circuito passivo.}
\end{questao}

\begin{questao}[4]{}
\textbf{(Resistividade e temperatura.)} A resistividade de um metal varia como
$\rho(T)=\rho_0\left[1+\alpha(T-T_0)\right]$.
\begin{itensq}
\item Mostre que a mesma lei vale para $R$, e identifique a hipótese oculta.
\item Determine o erro cometido ao ignorar a dilatação térmica das dimensões.
\item Avalie para cobre entre $\SI{20}{\celsius}$ e
      $\SI{120}{\celsius}$.
\end{itensq}
\resposta{A hipótese oculta é dimensões constantes; o erro é da ordem de
$\SI{0.05}{\percent}$ --- desprezível}
\resolucao{(a) De $R=\rho L/A$, se $L$ e $A$ forem constantes, então
$R\propto\rho$ e a lei se transfere diretamente. \textbf{A hipótese oculta é
justamente essa}: que a geometria não muda com a temperatura --- o que é falso,
pois todo metal se dilata.\\
(b) Com coeficiente de dilatação linear $\lambda$:
$L'=L(1+\lambda\Delta T)$ e $A'=A(1+\lambda\Delta T)^{2}$, logo
\[
\frac{L'}{A'}=\frac{L}{A}\cdot\frac{1}{1+\lambda\Delta T}.
\]
Assim
\[
R=\frac{\rho_0(1+\alpha\Delta T)}{1+\lambda\Delta T}\cdot\frac{L}{A}
\approx R_0\left[1+(\alpha-\lambda)\Delta T\right].
\]
\textbf{O efeito da dilatação é subtrair $\lambda$ de $\alpha$} --- a peça fica
mais longa (aumentando $R$) mas também mais grossa (reduzindo $R$), e o segundo
efeito predomina por ser quadrático na área.\\
(c) Para o cobre, $\alpha\approx\SI{3.9e-3}{\per\celsius}$ e
$\lambda\approx\SI{1.7e-5}{\per\celsius}$:
\[
\frac{\lambda}{\alpha}=\frac{1{,}7\times10^{-5}}{3{,}9\times10^{-3}}
=0{,}0044=0{,}44\%.
\]
Com $\Delta T=\SI{100}{\celsius}$, $R$ aumenta $39\%$ pela resistividade e
diminui $0{,}17\%$ pela dilatação --- efeito líquido de $38{,}8\%$.\\
\textbf{Conclusão:} ignorar a dilatação introduz erro de cerca de
$0{,}4\%$ \emph{do efeito térmico}, ou $0{,}17\%$ do valor de $R$. Desprezível
em quase toda aplicação --- mas não em padrões de resistência de laboratório,
onde se buscam partes por milhão.}
\end{questao}

\begin{questao}[4]{}
\textbf{(Medição a quatro fios.)} Quantifique a vantagem do método de Kelvin.
\begin{itensq}
\item Deduza o erro relativo do método de dois fios em função de
      $R_{cabo}/R_x$.
\item Determine a partir de qual valor de $R_x$ o método de dois fios dá erro
      inferior a $\SI{1}{\percent}$, com cabos de $\SI{0.2}{\ohm}$ cada.
\item Explique por que o método de quatro fios não elimina \emph{todos} os
      erros.
\end{itensq}
\resposta{Erro $=2R_{cabo}/R_x$; é preciso $R_x\ge\SI{40}{\ohm}$}
\resolucao{(a) O instrumento lê $R_x+2R_{cabo}$, logo
\[
\varepsilon=\frac{2R_{cabo}}{R_x}.
\]
(b) Impondo $\varepsilon\le0{,}01$ com $R_{cabo}=\SI{0.2}{\ohm}$:
\[
R_x\ge\frac{2(0{,}2)}{0{,}01}=\SI{40}{\ohm}.
\]
\textbf{Regra prática:} o método de dois fios só é aceitável para resistências
de dezenas de ohms ou mais. Abaixo disso, o erro cresce rapidamente --- em
$\SI{1}{\ohm}$ ele já é de $40\%$.\\
(c) \textbf{O que o método de quatro fios \emph{não} elimina.}
\begin{itemize}
\item \textbf{Tensões termoelétricas.} Junções entre metais diferentes nos
contatos geram forças eletromotrizes de dezenas de microvolts --- comparáveis ao
sinal quando se mede miliohms. Combate-se invertendo o sentido da corrente e
tomando a média das duas leituras.
\item \textbf{Autoaquecimento} pela corrente de excitação, que altera o próprio
$R_x$.
\item \textbf{Localização dos contatos de tensão.} Eles devem estar
\emph{entre} os de corrente, e a distribuição de corrente na peça influencia o
que se mede. Em amostras curtas, a definição do que é "o resistor" torna-se
ambígua.
\item \textbf{Ruído e deriva} do amplificador, agravados porque o sinal é de
poucos milivolts.
\end{itemize}
\textbf{Síntese:} o método de quatro fios elimina o erro \emph{sistemático e
dominante} --- a resistência dos cabos. Os erros remanescentes são de outra
natureza e exigem outras técnicas.}
\end{questao}

\begin{questao}[4]{}
\textbf{(Condutor composto.)} Um cabo de alta tensão tem alma de aço
($\rho_{\text{aço}}=\pot{1.5}{-7}\,\si{\ohm\meter}$, $A=\SI{20}{\milli\meter\squared}$)
envolvida por alumínio ($\rho_{Al}=\pot{2.8}{-8}\,\si{\ohm\meter}$,
$A=\SI{100}{\milli\meter\squared}$), com $\SI{1}{\kilo\meter}$ de comprimento.
\begin{itensq}
\item Determine a resistência de cada material e a do conjunto.
\item Determine a fração de corrente em cada um.
\item Explique a função da alma de aço.
\end{itensq}
\resposta{$R_{\text{aço}}=\SI{7.5}{\ohm}$, $R_{Al}=\SI{0.28}{\ohm}$,
$R_{eq}=\SI{0.270}{\ohm}$; o aço conduz $3{,}6\%$}
\resolucao{(a) Os dois materiais estão em \textbf{paralelo} (mesmo comprimento,
mesmos terminais):
\[
R_{\text{aço}}=\frac{\pot{1.5}{-7}(1000)}{\pot{20}{-6}}=\SI{7.5}{\ohm};
\qquad
R_{Al}=\frac{\pot{2.8}{-8}(1000)}{\pot{100}{-6}}=\SI{0.28}{\ohm};
\]
\[
R_{eq}=\frac{7{,}5(0{,}28)}{7{,}78}=\SI{0.270}{\ohm}.
\]
(b) Pelo divisor de corrente em condutâncias:
\[
\frac{I_{Al}}{I_T}=\frac{1/0{,}28}{1/0{,}28+1/7{,}5}=96{,}4\%;
\qquad
\frac{I_{\text{aço}}}{I_T}=3{,}6\%.
\]
(c) \textbf{A alma de aço praticamente não conduz --- e não é essa a sua
função.} Ela existe por razões \textbf{mecânicas}:
\begin{itemize}
\item O alumínio tem resistência à tração baixa e escoa sob carga permanente. O
aço sustenta o peso próprio do cabo entre torres, que podem distar centenas de
metros.
\item Reduz a flecha (o "barrigamento") do vão e, com ela, a altura necessária
das torres.
\item Suporta os esforços de vento e de acúmulo de gelo.
\end{itemize}
\textbf{Observe a divisão de papéis:} o alumínio conduz $96\%$ da corrente e o
aço sustenta praticamente toda a carga mecânica. É o princípio do cabo
\textbf{ACSR} (\emph{aluminium conductor steel-reinforced}), padrão em linhas de
transmissão aéreas --- um caso em que a análise elétrica revela que o componente
"elétrico" está ali por outro motivo.}
\end{questao}

\begin{questao}[4]{}
\textbf{(Limites da lei de Ohm.)} Discuta situações em que a relação
$U=RI$ deixa de descrever o comportamento de um condutor.
\begin{itensq}
\item Enumere ao menos quatro mecanismos, com exemplos.
\item Explique por que a lei de Ohm não é uma lei fundamental.
\item Indique como se procede quando ela falha.
\end{itensq}
\resposta{Autoaquecimento, não linearidade intrínseca, efeitos de alta
frequência e ruptura do meio}
\resolucao{(a) \textbf{Quatro mecanismos.}
\begin{enumerate}
\item \textbf{Autoaquecimento.} A resistência depende da temperatura, que
depende da potência dissipada. Exemplo: lâmpada incandescente, cuja resistência
a frio é cerca de dez vezes menor que a quente.
\item \textbf{Não linearidade intrínseca.} Junções semicondutoras seguem
relação exponencial, não linear. Exemplo: diodo, cuja "resistência" varia
ordens de grandeza em uma década de tensão.
\item \textbf{Efeitos de frequência.} O efeito pelicular concentra a corrente na
superfície do condutor, aumentando a resistência efetiva; somam-se indutância e
capacitância parasitas. Exemplo: um fio que tem $\SI{1}{\milli\ohm}$ em CC pode
apresentar dezenas de miliohms em megahertz.
\item \textbf{Ruptura do meio.} Acima da rigidez dielétrica, um isolante
torna-se condutor abruptamente. Exemplo: arco elétrico, cuja característica
$U\times I$ tem inclinação \emph{negativa}.
\end{enumerate}
\emph{(Acrescente-se ainda a saturação de velocidade dos portadores em campos
muito intensos e a supercondutividade, em que $R$ colapsa a zero.)}\\
(b) \textbf{A lei de Ohm é uma lei empírica e aproximada}, não fundamental. Ela
descreve o comportamento de \emph{certos} materiais (metais, em faixa limitada
de temperatura e frequência) e emerge de um modelo estatístico de colisões dos
portadores --- não de um princípio de conservação.\\
Compare com as leis de Kirchhoff, que decorrem da conservação de carga e de
energia e valem \emph{sempre}, para qualquer componente, linear ou não.\\
(c) \textbf{Quando a lei falha, procede-se assim:}
\begin{itemize}
\item Substituir $R$ pela \textbf{característica $U\times I$} completa do
dispositivo, obtida por medição ou por modelo.
\item Resolver o circuito por \textbf{interseção de características} (reta de
carga), como se faz com diodos.
\item Para pequenos sinais, linearizar em torno do ponto de operação usando a
\textbf{resistência dinâmica} --- e aí toda a maquinaria linear volta a valer.
\end{itemize}
\textbf{As leis de Kirchhoff continuam válidas em todos esses casos} --- é
sempre por elas que se começa.}
\end{questao}

% END BANCO COMPLEMENTAR

\gabarito[2.2 Lei de Ohm e resistividade]

% =====================================================================
\subsection{Potência elétrica e energia}
% =====================================================================

\blocofixacao

\begin{questao}[1]{}
A tensão nos terminais de um condutor é $\SI{12}{\volt}$ e ele dissipa
$\SI{24}{\watt}$.
\begin{itensq}
\item Qual é a corrente que o percorre?
\item Qual é o valor de sua resistência?
\end{itensq}
\resposta{(a) \SI{2}{\ampere}; (b) \SI{6}{\ohm}}
\resolucao{(a) De $P = U\,I$: $I = \dfrac{P}{U} = \dfrac{24}{12} = \SI{2}{\ampere}$.\\
(b) $R = \dfrac{U}{I} = \dfrac{12}{2} = \SI{6}{\ohm}$
(ou $R = \dfrac{U^2}{P} = \dfrac{144}{24} = \SI{6}{\ohm}$).}
\end{questao}

\begin{questao}[1]{}
Um aparelho opera em $\SI{220}{\volt}$ e consome $\SI{440}{\watt}$.
\begin{itensq}
\item Qual é a corrente que ele consome?
\item Qual sua resistência equivalente?
\end{itensq}
\resposta{(a) \SI{2}{\ampere}; (b) \SI{110}{\ohm}}
\resolucao{(a) $I = \dfrac{P}{U} = \dfrac{440}{220} = \SI{2}{\ampere}$.\\
(b) $R = \dfrac{U^2}{P} = \dfrac{220^2}{440} = \SI{110}{\ohm}$.}
\end{questao}

\begin{questao}[1]{Apostila original revisada}
Uma fonte fornece $\SI{100}{\watt}$ a um circuito, mas somente
$\SI{95}{\watt}$ são aproveitados pela carga. Determine o rendimento do circuito
e a potência dissipada em perdas.
\resposta{$\eta=95\%$; perdas de $\SI{5}{\watt}$}
\resolucao{$\eta=P_{\text{útil}}/P_{\mathrm{entrada}}=95/100=0{,}95$.
A diferença, $100-95=\SI{5}{\watt}$, é dissipada em fios, contatos e outros elementos.}
\end{questao}

\blocoaplicacao

\begin{questao}[2]{}
Uma lâmpada de especificação $\SI{60}{\watt}$/$\SI{220}{\volt}$ é ligada, por
engano, a uma rede de $\SI{110}{\volt}$. Sua potência dissipada passa a ser:
\begin{alternativas}
\item inalterada.
\item reduzida à metade.
\item duplicada.
\item reduzida à quarta parte.
\item aumentada quatro vezes.
\end{alternativas}
\resposta{(d)}
\resolucao{A resistência da lâmpada é fixa, determinada por seus valores nominais:
$R = \dfrac{U_{\text{nom}}^2}{P_{\text{nom}}} = \dfrac{220^2}{60} \approx
\SI{807}{\ohm}$. Sob a nova tensão, $P' = \dfrac{U'^2}{R}$. Como a tensão caiu à
metade e $P \propto U^2$:
\[
P' = \frac{(110)^2}{R} = \frac{60}{4} = \SI{15}{\watt},
\]
ou seja, a quarta parte. Meia tensão dissipa um quarto da potência.}
\end{questao}

\begin{questao}[2]{}
Um chuveiro elétrico tem especificações $\SI{220}{\volt}$/$\SI{4400}{\watt}$. Se
for ligado a uma rede de $\SI{110}{\volt}$, a potência dissipada será:
\begin{alternativas}[2]
\item \SI{550}{\watt}
\item \SI{1100}{\watt}
\item \SI{2200}{\watt}
\item \SI{4400}{\watt}
\item \SI{8800}{\watt}
\end{alternativas}
\resposta{(b)}
\resolucao{Mesma ideia da questão anterior: metade da tensão dissipa um quarto da
potência:
\[
P' = \frac{4400}{4} = \SI{1100}{\watt}.
\]
Por isso um chuveiro projetado para \SI{220}{\volt} "esquenta pouco" em rede de
\SI{110}{\volt}.}
\end{questao}

\begin{questao}[2]{}
Dois resistores $R_1 = \SI{10}{\ohm}$ e $R_2 = \SI{20}{\ohm}$ estão em série sob
uma tensão total de $\SI{60}{\volt}$.
\begin{itensq}
\item Qual é a corrente do circuito?
\item Qual é a potência dissipada em cada resistor?
\end{itensq}
\resposta{(a) \SI{2}{\ampere}; (b) $P_1=\SI{40}{\watt}$, $P_2=\SI{80}{\watt}$}
\resolucao{(a) $I = \dfrac{60}{10+20} = \SI{2}{\ampere}$.\\
(b) Com a mesma corrente, $P = R\,I^2$:
\[
P_1 = 10\cdot 2^2 = \SI{40}{\watt},\qquad
P_2 = 20\cdot 2^2 = \SI{80}{\watt}.
\]
\textbf{Verificação:} $P_{\text{total}} = U\,I = 60\cdot2 = \SI{120}{\watt}
= 40+80$. Na série, o maior resistor dissipa a maior potência.}
\end{questao}

\begin{questao}[2]{}
Um chuveiro elétrico de $\SI{5500}{\watt}$ opera em $\SI{220}{\volt}$.
\begin{itensq}
\item Qual é a corrente que ele exige da rede?
\item Qual é a resistência da sua resistência de aquecimento?
\item O disjuntor do circuito é de $\SI{20}{\ampere}$. Ele suporta esse chuveiro?
\end{itensq}
\resposta{(a) \SI{25}{\ampere}; (b) \SI{8.8}{\ohm}; (c) não}
\resolucao{(a) $I = \dfrac{P}{U} = \dfrac{5500}{220} = \SI{25}{\ampere}$.\\
(b) $R = \dfrac{U^2}{P} = \dfrac{220^2}{5500} = \SI{8.8}{\ohm}$.\\
(c) A corrente exigida (\SI{25}{\ampere}) \emph{excede} os \SI{20}{\ampere} do
disjuntor, que desarmaria. Seria necessário um disjuntor de pelo menos
\SI{25}{\ampere} (na prática, \SI{32}{\ampere}) e fiação compatível --- por isso
chuveiros exigem circuito exclusivo.}
\end{questao}

\begin{questao}[2]{Apostila original revisada}
Em uma residência, a tarifa média foi obtida de uma conta de
R\$~80,00 para um consumo de $\SI{200}{\kilo\watt\hour}$. Um chuveiro de
$\SI{3}{\kilo\watt}$ é utilizado durante $\SI{30}{\minute}$ por dia, ao longo de
30 dias. Calcule a energia mensal consumida pelo chuveiro e seu custo.
\resposta{$\SI{45}{\kilo\watt\hour}$; R\$~18,00}
\resolucao{A tarifa é $80/200=\text{R\$ }0{,}40/\si{\kilo\watt\hour}$.
O tempo mensal é $0{,}5\times30=\SI{15}{\hour}$, então
$E=3\times15=\SI{45}{\kilo\watt\hour}$ e o custo é $45\times0{,}40=\text{R\$ }18{,}00$.}
\end{questao}

\blocoaprofundamento

\begin{questao}[3]{}
Em uma residência permanecem ligados, durante $\SI{2}{\hour}$, um ferro elétrico
de $\SI{1440}{\watt}$/$\SI{120}{\volt}$ e duas lâmpadas de
$\SI{60}{\watt}$/$\SI{120}{\volt}$. Mantida a tensão em $\SI{120}{\volt}$,
determine a corrente total no fusível e a energia elétrica consumida.
\resposta{$I = \SI{13}{\ampere}$; $E = \SI{3.12}{\kilo\watt\hour}$}
\resolucao{Potência total: $P = 1440 + 2\cdot60 = \SI{1560}{\watt}$.\\
Corrente no fusível (todos em paralelo, sob \SI{120}{\volt}):
\[
I = \frac{P}{U} = \frac{1560}{120} = \SI{13}{\ampere}.
\]
Energia em $\SI{2}{\hour}$:
\[
E = P\,t = 1{,}560\,\si{\kilo\watt}\cdot2\,\si{\hour} = \SI{3.12}{\kilo\watt\hour}.
\]
O kWh é a unidade da conta de luz: potência (em kW) vezes tempo (em h).}
\end{questao}

\begin{questao}[3]{}
Uma residência utiliza, por dia: uma lâmpada de $\SI{100}{\watt}$ acesa
$\SI{5}{\hour}$, um chuveiro de $\SI{4400}{\watt}$ ligado $\SI{0.5}{\hour}$ e uma
geladeira de $\SI{200}{\watt}$ que funciona, em média, $\SI{8}{\hour}$ por dia.
Calcule o consumo mensal de energia (30 dias) e o custo, considerando a tarifa de
$\SI{0.80}{\text{R\$}\per\kilo\watt\hour}$.
\resposta{$E \approx \SI{127.5}{\kilo\watt\hour}$; custo $\approx$ R\$\,102,00}
\resolucao{Energia diária de cada aparelho ($E = P\,t$, em kWh):
\[
\text{lâmpada: } 0{,}1\cdot5 = \SI{0.5}{\kilo\watt\hour};\quad
\text{chuveiro: } 4{,}4\cdot0{,}5 = \SI{2.2}{\kilo\watt\hour};\quad
\text{geladeira: } 0{,}2\cdot8 = \SI{1.6}{\kilo\watt\hour}.
\]
Total diário: $0{,}5 + 2{,}2 + 1{,}6 = \SI{4.3}{\kilo\watt\hour}$. Em 30 dias:
\[
E = 4{,}3\cdot30 = \SI{129}{\kilo\watt\hour}
\quad\Rightarrow\quad
\text{custo} = 129\cdot0{,}80 \approx \text{R\$}\,103{,}00.
\]
(Valores próximos de \SI{127.5}{\kilo\watt\hour} conforme arredondamentos.) É
assim que se estima a conta de luz --- e fica claro que o chuveiro, apesar de
usado poucos minutos, é um dos maiores vilões pelo seu alto \emph{potência}.}
\end{questao}

\begin{questao}[3]{}
Duas lâmpadas incandescentes de especificações $\SI{60}{\watt}$/$\SI{220}{\volt}$
e $\SI{100}{\watt}$/$\SI{220}{\volt}$ são ligadas \textbf{em série} à rede de
$\SI{220}{\volt}$.
\begin{itensq}
\item Calcule a resistência de cada lâmpada (suposta constante).
\item Determine a potência dissipada por cada uma nessa ligação.
\item \textbf{Qual delas brilha mais?} O resultado contraria a intuição?
      Explique.
\end{itensq}
\resposta{(a) $\SI{806.7}{\ohm}$ e $\SI{484}{\ohm}$;
(b) $\SI{23.4}{\watt}$ e $\SI{14.1}{\watt}$; (c) a de \SI{60}{\watt}}
\resolucao{\textbf{(a)} Das especificações nominais, $R = \dfrac{U^2}{P}$:
\[
R_{60}=\frac{220^2}{60}\approx\SI{806.7}{\ohm},
\qquad
R_{100}=\frac{220^2}{100}=\SI{484}{\ohm}.
\]
Note: \emph{a lâmpada de menor potência tem MAIOR resistência.}

\textbf{(b)} Em série, a corrente é comum:
\[
I=\frac{220}{806{,}7+484}=\frac{220}{1290{,}7}\approx\SI{0.170}{\ampere}.
\]
Com $P = R I^2$:
\[
P_{60}=806{,}7\,(0{,}170)^2\approx\SI{23.4}{\watt},
\qquad
P_{100}=484\,(0{,}170)^2\approx\SI{14.1}{\watt}.
\]
\textbf{Verificação:} $23{,}4+14{,}1 = \SI{37.5}{\watt} = 220\cdot0{,}170$.
\checkmark

\textbf{(c) A lâmpada de \SI{60}{\watt} brilha MAIS} --- resultado que contraria
a intuição de quem associa "100 W" a "mais brilhante".\\
\textbf{Por quê:} as especificações valem \emph{apenas} na tensão nominal, com
cada lâmpada recebendo os \SI{220}{\volt} completos. Em série, elas
\textbf{compartilham} a mesma corrente, e quem dita a potência passa a ser a
resistência: $P = RI^2$. A de \SI{60}{\watt}, tendo maior $R$, fica com maior
parcela da tensão e dissipa mais.\\
\textbf{Regra geral:}
\begin{center}\small\sffamily
\begin{tabular}{@{}l l l@{}}
\toprule
\textbf{Ligação} & \textbf{Grandeza comum} & \textbf{Brilha mais}\\
\midrule
Série    & corrente $I$ & a de \emph{maior} $R$ (menor potência nominal)\\
Paralelo & tensão $U$   & a de \emph{menor} $R$ (maior potência nominal)\\
\bottomrule
\end{tabular}
\end{center}
Observe ainda que \emph{ambas} ficam muito abaixo do nominal (\SI{37.5}{\watt} no
total contra \SI{160}{\watt} se ligadas corretamente em paralelo).}
\end{questao}

\begin{questao}[3]{}
Um chuveiro elétrico de $\SI{220}{\volt}$ possui uma resistência de aquecimento
com uma \emph{derivação central}, permitindo duas posições: na posição "verão",
usa-se a resistência inteira ($\SI{8.8}{\ohm}$); na posição "inverno", apenas
\emph{metade} dela.
\begin{itensq}
\item Calcule a potência em cada posição.
\item Determine a corrente em cada caso e verifique se um disjuntor de
      $\SI{40}{\ampere}$ é adequado.
\item Sabendo que a água entra a $\SI{20}{\celsius}$ com vazão de
      $\SI{3}{\liter\per\minute}$, estime a temperatura de saída no inverno.
\end{itensq}
\dados{$c_{\text{água}}=\SI{4180}{\joule\per\kilogram\per\celsius}$;
$\rho_{\text{água}}=\SI{1}{\kilogram\per\liter}$.}
\resposta{(a) \SI{5500}{\watt} e \SI{11000}{\watt}; (b) \SI{25}{\ampere} e
\SI{50}{\ampere} --- disjuntor insuficiente; (c) $\approx\SI{72}{\celsius}$}
\resolucao{\textbf{(a)} Com $P = \dfrac{U^2}{R}$:
\[
P_{\text{verão}}=\frac{220^2}{8{,}8}=\SI{5500}{\watt},
\qquad
P_{\text{inverno}}=\frac{220^2}{4{,}4}=\SI{11000}{\watt}.
\]
\emph{Metade da resistência dobra a potência} --- e não a reduz, como a intuição
poderia sugerir.

\textbf{(b)} $I = \dfrac{P}{U}$:
$I_{\text{verão}} = \SI{25}{\ampere}$ e $I_{\text{inverno}} = \SI{50}{\ampere}$.
O disjuntor de \SI{40}{\ampere} \textbf{não é adequado} --- ele desarmaria na
posição inverno. Seria necessário um de \SI{63}{\ampere} com fiação
de $\SI{10}{\milli\meter\squared}$.

\textbf{(c) Balanço térmico.} A vazão mássica é
\[
\dot m = \SI{3}{\liter\per\minute} = \frac{3}{60}=\SI{0.05}{\kilogram\per\second}.
\]
Toda a potência elétrica vira calor na água:
\[
P = \dot m\,c\,\Delta T
\;\Rightarrow\;
\Delta T=\frac{P}{\dot m\,c}=\frac{11000}{0{,}05\cdot4180}\approx\SI{52.6}{\celsius}.
\]
Logo a saída fica a $20+52{,}6 \approx \SI{72}{\celsius}$ --- alta demais para
banho, o que na prática é ajustado aumentando a vazão.\\
\textbf{Síntese de engenharia:} este exercício conecta \emph{três} domínios ---
circuito elétrico (lei de Ohm), proteção (dimensionamento de disjuntor) e
termodinâmica (balanço de energia). É exatamente esse tipo de integração que
caracteriza o trabalho profissional real.}
\end{questao}

% BEGIN BANCO COMPLEMENTAR Potência elétrica e energia
% =====================================================================
%  Banco complementar --- Potencia Eletrica e Energia  (REVISADO)
%  Substitui a versao gerada por template (7 clones em N2, 8 em N3 e
%  8 questoes conceituais marcadas como N4).
%  O conteudo de ENERGIA esta no cap02, que ja cobre: rendimento
%  100 W/95 W, tarifa media a partir de conta de R\$ 80 / 200 kWh,
%  ferro eletrico e outros aparelhos por 2 h, e o consumo diario
%  residencial (lampada, chuveiro etc.). O cap05 cobre P=RI^2, P=VI
%  e maxima transferencia. O banco de circuitos serie revisado ja usa
%  um calculo de custo mensal (300 W, 6 h/dia, 30 dias).
%  Este banco explora: potencia MEDIA x valor eficaz, regime pulsado
%  e constante termica, rendimento em cascata, tensao errada de
%  alimentacao, energia em transitorio exponencial, payback e o
%  dimensionamento economico de condutor.
%  Distribuicao mantida: 5 N1 + 7 N2 + 8 N3 + 8 N4 = 28
% =====================================================================

\subsubsection*{Banco complementar --- Potência elétrica e energia}

\begin{atencao}[Organização do banco]
As três formas de $P$ ($UI$, $RI^{2}$ e $U^{2}/R$) são equivalentes apenas
quando vale a lei de Ohm; escolha a que usa os dados disponíveis sem passos
intermediários. Para \emph{energia}, lembre que $\SI{1}{\kilo\watt\hour}
=\SI{3.6}{\mega\joule}$. O N3 trata de regime pulsado, rendimento em cascata e
tensão incorreta de alimentação; o N4 exige integração, demonstração de que a
potência média usa o valor \emph{eficaz}, e análises de retorno econômico.
\end{atencao}

\blocofixacao

\begin{questao}[1]{}
Um resistor de $\SI{8}{\ohm}$ é percorrido por $\SI{2.5}{\ampere}$. Determine a
potência dissipada.
\resposta{$P=\SI{50}{\watt}$}
\resolucao{
\[
P=RI^{2}=8(2{,}5)^{2}=8(6{,}25)=\SI{50}{\watt}.
\]
\textbf{Escolha da fórmula:} conhecendo $R$ e $I$, use $RI^{2}$ diretamente.
Calcular $U=RI=\SI{20}{\volt}$ e depois $P=UI$ dá o mesmo resultado, mas
acrescenta um passo e uma oportunidade de erro.}
\end{questao}

\begin{questao}[1]{}
Um resistor de $\SI{44}{\ohm}$ é submetido a $\SI{220}{\volt}$. Determine a
potência.
\resposta{$P=\SI{1100}{\watt}$}
\resolucao{
\[
P=\frac{U^{2}}{R}=\frac{220^{2}}{44}=\frac{48400}{44}=\SI{1100}{\watt}.
\]
Note a dependência \emph{quadrática} na tensão: uma queda de $10\%$ na rede
($\SI{198}{\volt}$) reduziria a potência para
$\SI{891}{\watt}$ --- $19\%$ a menos. É por isso que chuveiros "esquentam
menos" em horários de pico.}
\end{questao}

\begin{questao}[1]{}
Converta $\SI{250}{\watt\hour}$ para joules.
\resposta{$\SI{900}{\kilo\joule}$}
\resolucao{
\[
E=Pt=250\,\si{\watt}\times3600\,\si{\second}=\SI{900000}{\joule}
=\SI{900}{\kilo\joule}.
\]
\textbf{Fator de conversão útil:}
$\SI{1}{\kilo\watt\hour}=\SI{3.6}{\mega\joule}$. O joule é pequeno demais para
contas de energia elétrica --- daí o uso do quilowatt-hora, que é uma unidade
de energia, não de potência. Confundir os dois é o erro mais comum do tópico.}
\end{questao}

\begin{questao}[1]{}
Conhecendo apenas a tensão sobre um resistor e o valor de sua resistência, a
expressão adequada para a potência é:
\begin{alternativas}[2]
\item $P=UI$
\item $P=RI^{2}$
\item $P=U^{2}/R$
\item $P=U/R$
\end{alternativas}
\resposta{(c)}
\resolucao{As alternativas (a) e (b) exigiriam calcular $I$ antes. A (d) está
dimensionalmente errada: $\si{\volt}/\si{\ohm}$ é ampère, não watt.\\
\textbf{Verificação dimensional rápida:} $U^{2}/R$ dá
$\si{\volt\squared}/\si{\ohm}=\si{\volt}\cdot\si{\ampere}=\si{\watt}$
\checkmark. Sempre que houver dúvida entre fórmulas parecidas, confira a
dimensão --- ela elimina a maioria das alternativas erradas.}
\end{questao}

\begin{questao}[1]{}
Um dispositivo de $\SI{60}{\watt}$ opera por $\SI{5}{\minute}$. Determine a
energia consumida em joules.
\resposta{$E=\SI{18}{\kilo\joule}$}
\resolucao{Convertendo o tempo para segundos:
\[
E=Pt=60\times(5\times60)=60(300)=\SI{18000}{\joule}=\SI{18}{\kilo\joule}.
\]
\textbf{Cuidado com as unidades:} $P$ em watts exige $t$ em \emph{segundos}
para obter joules. Se o tempo ficar em minutos, o resultado sai
$60\times$ menor --- erro silencioso, porque o número parece plausível.}
\end{questao}

\blocoaplicacao

\begin{questao}[2]{}
Um chuveiro de $\SI{5500}{\watt}$ opera em $\SI{220}{\volt}$. Determine a
corrente e a resistência.
\resposta{$I=\SI{25}{\ampere}$; $R=\SI{8.8}{\ohm}$}
\resolucao{
\[
I=\frac{P}{U}=\frac{5500}{220}=\SI{25}{\ampere};
\qquad
R=\frac{U^{2}}{P}=\frac{48400}{5500}=\SI{8.8}{\ohm}.
\]
\textbf{Conferência:} $R=U/I=220/25=\SI{8.8}{\ohm}$ \checkmark.\\
\textbf{Implicação de instalação:} $\SI{25}{\ampere}$ é uma corrente elevada.
Ela exige circuito exclusivo, condutor de seção adequada e disjuntor de
$\SI{32}{\ampere}$ --- e é por isso que chuveiro nunca compartilha circuito com
tomadas.}
\end{questao}

\begin{questao}[2]{}
Projete a resistência de um aquecedor de $\SI{1200}{\watt}$ para
$\SI{220}{\volt}$ e determine a corrente nominal e o disjuntor adequado.
\resposta{$R=\SI{40.3}{\ohm}$; $I=\SI{5.45}{\ampere}$; disjuntor de
$\SI{10}{\ampere}$}
\resolucao{
\[
R=\frac{U^{2}}{P}=\frac{48400}{1200}=\SI{40.3}{\ohm};
\qquad
I=\frac{1200}{220}=\SI{5.45}{\ampere}.
\]
\textbf{Disjuntor:} o valor comercial imediatamente acima é
$\SI{10}{\ampere}$, que oferece folga de $83\%$ --- adequado, pois aquecedores
resistivos não têm corrente de partida elevada.\\
\textbf{Nota sobre o material:} $\SI{40.3}{\ohm}$ deve ser obtido com fio de
liga resistiva (níquel-cromo) de comprimento e seção calculados. E como a
resistência desses materiais varia pouco com a temperatura, a potência
permanece próxima do nominal em regime.}
\end{questao}

\begin{questao}[2]{}
Duas lâmpadas incandescentes de $\SI{60}{\watt}$ são fabricadas para
$\SI{127}{\volt}$ e para $\SI{220}{\volt}$. Determine a resistência de cada
uma.
\resposta{$\SI{268.8}{\ohm}$ e $\SI{806.7}{\ohm}$}
\resolucao{
\[
R_{127}=\frac{127^{2}}{60}=\SI{268.8}{\ohm};
\qquad
R_{220}=\frac{220^{2}}{60}=\SI{806.7}{\ohm}.
\]
A lâmpada de tensão maior tem resistência \textbf{três vezes} maior --- razão
$(220/127)^{2}=3{,}0$.\\
\textbf{Consequência construtiva:} para a mesma potência, o filamento de
$\SI{220}{\volt}$ é mais longo e mais fino. Ele é, portanto, mais frágil
mecanicamente --- fato que se observa na prática, com lâmpadas de tensão maior
tendendo a durar menos sob vibração.}
\end{questao}

\begin{questao}[2]{}
Um aparelho opera em dois modos: $\SI{1500}{\watt}$ durante
$\SI{20}{\minute}$ e $\SI{600}{\watt}$ durante $\SI{40}{\minute}$, todos os
dias. Determine o consumo mensal ($\SI{30}{\day}$) e o custo a
$\mathrm{R\$}\,0{,}95$ por $\si{\kilo\watt\hour}$.
\resposta{$\SI{27}{\kilo\watt\hour}$; $\mathrm{R\$}\,25{,}65$}
\resolucao{Energia diária, com os tempos em horas:
\[
E=1{,}5\left(\tfrac13\right)+0{,}6\left(\tfrac23\right)
=0{,}50+0{,}40=\SI{0.90}{\kilo\watt\hour}.
\]
\[
E_{\text{mês}}=0{,}90(30)=\SI{27}{\kilo\watt\hour};
\qquad
\text{custo}=27(0{,}95)=\mathrm{R\$}\,25{,}65.
\]
\textbf{Observação:} o modo de $\SI{600}{\watt}$, embora tenha menos da metade
da potência, responde por $44\%$ do consumo --- porque funciona o dobro do
tempo. Em energia, \textbf{potência sozinha não decide}; o produto pelo tempo é
que conta.}
\end{questao}

\begin{questao}[2]{}
Um resistor é submetido a pulsos de $\SI{50}{\watt}$ com ciclo de trabalho de
$\SI{10}{\percent}$. Determine a potência média.
\resposta{$P_{\text{média}}=\SI{5}{\watt}$}
\resolucao{
\[
P_{\text{média}}=P_{pico}\times D=50(0{,}10)=\SI{5}{\watt}.
\]
\textbf{Mas cuidado:} isso \emph{não} significa que um resistor de
$\SI{5}{\watt}$ sirva. Se a duração do pulso for comparável ou maior que a
constante térmica do componente, ele atinge temperaturas próximas às de
$\SI{50}{\watt}$ contínuos e pode falhar.\\
A questão de quando a potência média basta é retomada quantitativamente no bloco
de aprofundamento.}
\end{questao}

\begin{questao}[2]{}
Uma fonte com rendimento de $\SI{90}{\percent}$ alimenta um conversor de
$\SI{85}{\percent}$. Determine o rendimento global e a potência de entrada
necessária para entregar $\SI{100}{\watt}$ úteis.
\resposta{$\eta=76{,}5\%$; $P_{entrada}=\SI{130.7}{\watt}$}
\resolucao{Rendimentos em cascata \textbf{multiplicam-se}:
\[
\eta=0{,}90\times0{,}85=0{,}765=76{,}5\%.
\]
\[
P_{entrada}=\frac{100}{0{,}765}=\SI{130.7}{\watt}.
\]
\textbf{Erro comum:} \emph{somar} ou fazer a média dos rendimentos, o que daria
$87{,}5\%$ --- valor otimista demais. Cada estágio desperdiça uma fração do que
recebe, e as perdas se compõem multiplicativamente.\\
\textbf{Perdas:} $\SI{13.1}{\watt}$ na fonte e $\SI{17.6}{\watt}$ no
conversor, totalizando $\SI{30.7}{\watt}$ dissipados.}
\end{questao}

\begin{questao}[2]{}
Uma carga alterna entre $\SI{500}{\watt}$ por $\SI{5}{\minute}$ e
$\SI{100}{\watt}$ por $\SI{15}{\minute}$, em ciclos de
$\SI{20}{\minute}$. Determine a potência média e a energia por ciclo.
\resposta{$P_{\text{média}}=\SI{200}{\watt}$; $E=\SI{66.7}{\watt\hour}$}
\resolucao{A potência média é a \textbf{média ponderada pelo tempo}:
\[
P_{\text{média}}=\frac{500(5)+100(15)}{20}=\frac{2500+1500}{20}=\SI{200}{\watt}.
\]
\[
E_{ciclo}=P_{\text{média}}\times t_{ciclo}=200\times\frac{20}{60}
=\SI{66.7}{\watt\hour}.
\]
\textbf{Atenção:} a média \emph{aritmética} das potências ($300\,\si{\watt}$)
estaria errada --- ela ignoraria que o patamar baixo dura três vezes mais.
Média de potência é sempre ponderada pelo tempo.}
\end{questao}

\blocoaprofundamento

\begin{questao}[3]{}
Um resistor de $\SI{10}{\ohm}$ conduz $\SI{4}{\ampere}$ durante
$\SI{25}{\percent}$ do tempo e nada nos $\SI{75}{\percent}$ restantes.
\begin{itensq}
\item Determine a corrente média e a potência calculada a partir dela.
\item Determine a potência média verdadeira.
\item Explique a discrepância.
\end{itensq}
\resposta{(a) $\SI{1}{\ampere}$ e $\SI{10}{\watt}$ (errado); (b)
$\SI{40}{\watt}$; (c) a potência depende de $\langle i^{2}\rangle$}
\resolucao{(a) $I_{\text{média}}=0{,}25(4)+0{,}75(0)=\SI{1}{\ampere}$, e
$RI_{\text{média}}^{2}=10(1)=\SI{10}{\watt}$.\\
(b) A potência instantânea vale $\SI{160}{\watt}$ durante $25\%$ do tempo e
zero no resto:
\[
P_{\text{média}}=0{,}25(160)+0{,}75(0)=\SI{40}{\watt}.
\]
\textbf{Quatro vezes} o valor obtido em (a).\\
(c) A potência é \emph{quadrática} na corrente, e a média de um quadrado não é
o quadrado da média:
\[
P_{\text{média}}=R\left\langle i^{2}\right\rangle
\ne R\left\langle i\right\rangle^{2}.
\]
Define-se então o \textbf{valor eficaz}
$I_{rms}=\sqrt{\langle i^{2}\rangle}=\sqrt{0{,}25(16)}=\SI{2}{\ampere}$, com o
qual $P=RI_{rms}^{2}=\SI{40}{\watt}$ \checkmark.\\
\textbf{Definição operacional do valor eficaz:} é a corrente contínua que
produziria o \emph{mesmo aquecimento}. Aqui, $\SI{2}{\ampere}$ contínuos
aqueceriam tanto quanto o regime pulsado --- e é por isso que multímetros de
qualidade medem "true RMS".}
\end{questao}

\begin{questao}[3]{}
Um resistor com potência nominal de $\SI{10}{\watt}$ recebe pulsos de
$\SI{50}{\watt}$ com ciclo de trabalho de $\SI{10}{\percent}$.
\begin{itensq}
\item Determine a potência média e compare com o nominal.
\item Estabeleça o critério para decidir se o componente sobrevive.
\item Analise dois casos: pulsos de $\SI{1}{\milli\second}$ e de
      $\SI{10}{\second}$.
\end{itensq}
\resposta{$P_{\text{média}}=\SI{5}{\watt}$; a decisão depende da constante térmica}
\resolucao{(a) $P_{\text{média}}=\SI{5}{\watt}$, metade do nominal --- aparentemente
seguro.\\
(b) \textbf{O critério é a comparação entre a duração do pulso e a constante
térmica $\tau_t$ do componente} (tipicamente de segundos a dezenas de segundos
para resistores de fio; menos para os de filme).
\begin{itemize}
\item Se $t_{pulso}\ll\tau_t$, a massa térmica \emph{integra} os pulsos: a
temperatura acompanha a potência \textbf{média}, e o critério de
$\SI{5}{\watt}$ vale.
\item Se $t_{pulso}\gg\tau_t$, o resistor atinge o equilíbrio térmico
\emph{durante} cada pulso: a temperatura corresponde a $\SI{50}{\watt}$
contínuos, e o componente queima.
\end{itemize}
(c) \textbf{Pulsos de $\SI{1}{\milli\second}$:} muito menores que qualquer
$\tau_t$ realista --- o resistor "enxerga" $\SI{5}{\watt}$. \textbf{Seguro.}\\
\textbf{Pulsos de $\SI{10}{\second}$:} comparáveis ou maiores que $\tau_t$ --- o
resistor esquenta como se recebesse $\SI{50}{\watt}$ por dez segundos.
\textbf{Falha provável.}\\
\textbf{Consequência prática:} fabricantes publicam curvas de \emph{potência de
pulso} justamente porque o valor nominal contínuo não responde a essa pergunta.
Nunca dimensione regime pulsado apenas pela média.}
\end{questao}

\begin{questao}[3]{}
Um sistema tem uma fonte de $\SI{92}{\percent}$ em cascata com um conversor de
$\SI{85}{\percent}$.
\begin{itensq}
\item Determine o rendimento global.
\item Compare o ganho de melhorar o estágio pior (de $85$ para
      $\SI{92}{\percent}$) com o de melhorar o melhor (de $92$ para
      $\SI{97}{\percent}$).
\item Formule a regra geral.
\end{itensq}
\resposta{(a) $78{,}2\%$; (b) $84{,}6\%$ contra $82{,}5\%$ --- melhorar o pior
rende mais}
\resolucao{(a) $\eta=0{,}92(0{,}85)=0{,}782=78{,}2\%$.\\
(b) \textbf{Melhorando o pior:} $0{,}92(0{,}92)=84{,}6\%$ --- ganho de
$6{,}4$ pontos.\\
\textbf{Melhorando o melhor:} $0{,}97(0{,}85)=82{,}5\%$ --- ganho de
$4{,}3$ pontos.\\
E note que o segundo caso exigiu uma melhoria maior em pontos absolutos
($5$ contra $7$) para render \emph{menos}.\\
(c) \textbf{Regra geral.} Como $\eta=\prod\eta_k$, o ganho relativo global é a
soma dos ganhos relativos:
\[
\frac{\Delta\eta}{\eta}=\sum_k\frac{\Delta\eta_k}{\eta_k}.
\]
Um mesmo acréscimo \emph{absoluto} rende mais no estágio de menor $\eta_k$,
pois ali representa maior acréscimo relativo. \textbf{Ataque sempre o elo mais
fraco} --- e é também nele que a margem de melhoria costuma ser maior e mais
barata.}
\end{questao}

\begin{questao}[3]{}
Um aparelho consome $\SI{4}{\watt}$ em espera, permanentemente.
\begin{itensq}
\item Determine o consumo anual e o custo a
      $\mathrm{R\$}\,0{,}95/\si{\kilo\watt\hour}$.
\item Estime o impacto de dez aparelhos semelhantes.
\item Proponha uma decisão de engenharia.
\end{itensq}
\resposta{$\SI{35}{\kilo\watt\hour}$ e $\mathrm{R\$}\,33{,}29$ por aparelho;
$\mathrm{R\$}\,333$ para dez}
\resolucao{(a) Um ano tem $\SI{8760}{\hour}$:
\[
E=4\times8760=\SI{35040}{\watt\hour}=\SI{35.04}{\kilo\watt\hour};
\]
\[
\text{custo}=35{,}04(0{,}95)=\mathrm{R\$}\,33{,}29.
\]
(b) Dez aparelhos: $\SI{350}{\kilo\watt\hour}$ e cerca de
$\mathrm{R\$}\,333$ por ano --- comparável ao consumo anual de uma geladeira
eficiente.\\
(c) \textbf{Decisões possíveis, com seus custos.}
\begin{itemize}
\item \textbf{Filtro com chave} para desligar grupos de aparelhos: custo de
dezenas de reais, retorno em poucos meses. Melhor relação custo-benefício.
\item \textbf{Redesenho do produto:} normas modernas exigem espera abaixo de
$\SI{0.5}{\watt}$, o que reduziria o custo a menos de
$\mathrm{R\$}\,4{,}20$ por ano.
\item \textbf{Não fazer nada:} defensável para um aparelho isolado, insustentável
em escala --- o consumo agregado de espera responde por vários por cento do
consumo residencial nacional.
\end{itemize}
\textbf{O ponto de engenharia:} $\SI{4}{\watt}$ parecem irrelevantes porque a
potência é pequena. Mas energia é potência \emph{vezes tempo}, e
$\SI{8760}{\hour}$ é um multiplicador enorme.}
\end{questao}

\begin{questao}[3]{}
Um aparelho projetado para $\SI{220}{\volt}$ e $\SI{1000}{\watt}$ é ligado por
engano em $\SI{127}{\volt}$.
\begin{itensq}
\item Determine a potência resultante, supondo resistência constante.
\item Analise o caso inverso: aparelho de $\SI{127}{\volt}$ em
      $\SI{220}{\volt}$.
\end{itensq}
\resposta{(a) $\SI{333}{\watt}$ ($33\%$); (b) potência \emph{triplicada} ---
destruição}
\resolucao{(a) $R=\dfrac{220^{2}}{1000}=\SI{48.4}{\ohm}$, e em
$\SI{127}{\volt}$:
\[
P=\frac{127^{2}}{48{,}4}=\SI{333}{\watt}
=\left(\frac{127}{220}\right)^{2}\times1000=33{,}3\%.
\]
\textbf{Consequência:} o aparelho funciona, mas com um terço da potência --- um
chuveiro esquenta pouco, um aquecedor não aquece. Situação incômoda, porém
\emph{segura}.\\
(b) $P'=\left(\dfrac{220}{127}\right)^{2}P=3{,}0\,P$ --- a potência
\textbf{triplica}, e a corrente quase dobra.\\
\textbf{Resultado:} superaquecimento imediato, fumaça e, muitas vezes, incêndio.
O disjuntor pode nem atuar, pois a corrente ($\SI{13.6}{\ampere}$ para um
aparelho de $\SI{1000}{\watt}$/$\SI{127}{\volt}$) fica abaixo do valor de
disparo --- a proteção existe para o \emph{circuito}, não para o aparelho.\\
\textbf{Assimetria fundamental:} subtensão é inconveniente; sobretensão é
destrutiva. Por isso a etiqueta de tensão merece conferência antes de ligar
qualquer aparelho resistivo.}
\end{questao}

\begin{questao}[3]{}
Um resistor de $\SI{4}{\ohm}$ é percorrido por
$i(t)=\SI{5}{\ampere}\cdot e^{-t/\tau}$, com $\tau=\SI{20}{\milli\second}$.
\begin{itensq}
\item Determine a potência inicial.
\item Determine a energia total dissipada.
\item Determine o instante em que metade da energia já foi dissipada.
\end{itensq}
\resposta{(a) $\SI{100}{\watt}$; (b) $\SI{1}{\joule}$; (c)
$t=\SI{6.93}{\milli\second}$}
\resolucao{(a) $P_0=RI_0^{2}=4(25)=\SI{100}{\watt}$.\\
(b) A potência decai como $e^{-2t/\tau}$:
\[
E=\int_0^{\infty}RI_0^{2}e^{-2t/\tau}\,\mathrm{d}t
=RI_0^{2}\cdot\frac{\tau}{2}
=100\times\frac{0{,}020}{2}=\SI{1}{\joule}.
\]
\textbf{Regra útil:} a energia equivale à potência inicial mantida por
$\tau/2$. A "constante de tempo da potência" é metade da constante de tempo da
corrente, porque a potência vai com o quadrado.\\
(c) Metade da energia:
\[
\int_0^{t}e^{-2s/\tau}\mathrm{d}s=\frac12\int_0^{\infty}
\;\Longrightarrow\;
1-e^{-2t/\tau}=\frac12
\;\Longrightarrow\;
t=\frac{\tau\ln2}{2}=\SI{6.93}{\milli\second}.
\]
\textbf{Interpretação:} metade de toda a energia é liberada em apenas
$0{,}35\,\tau$ --- o transitório é fortemente concentrado no início. É o pico
inicial, e não a duração, que dimensiona o componente.}
\end{questao}

\begin{questao}[3]{}
Compare uma lâmpada incandescente de $\SI{60}{\watt}$ com uma de LED de
$\SI{9}{\watt}$ e mesmo fluxo luminoso, usadas $\SI{5}{\hour}$ por dia.
\begin{itensq}
\item Determine o consumo e o custo anual de cada uma, a
      $\mathrm{R\$}\,0{,}95/\si{\kilo\watt\hour}$.
\item Determine o tempo de retorno se a lâmpada de LED custar
      $\mathrm{R\$}\,15{,}00$ a mais.
\end{itensq}
\resposta{(a) $\mathrm{R\$}\,104{,}02$ contra $\mathrm{R\$}\,15{,}60$;
(b) cerca de $2 meses$}
\resolucao{(a) Horas anuais: $5\times365=\SI{1825}{\hour}$.
\[
E_{inc}=0{,}060(1825)=\SI{109.5}{\kilo\watt\hour}
\Rightarrow \mathrm{R\$}\,104{,}02;
\]
\[
E_{LED}=0{,}009(1825)=\SI{16.4}{\kilo\watt\hour}
\Rightarrow \mathrm{R\$}\,15{,}60.
\]
Economia anual: $\mathrm{R\$}\,88{,}42$.\\
(b)
\[
t=\frac{15{,}00}{88{,}42}=0{,}17\ \text{ano}\approx2 meses.
\]
\textbf{Retorno excepcional} --- poucos investimentos em engenharia se pagam em
dois meses.\\
\textbf{E há mais:} a lâmpada de LED dura tipicamente $\SI{25000}{\hour}$
contra $\SI{1000}{\hour}$ da incandescente, evitando cerca de $24$ trocas ao
longo da vida. Somando o custo das lâmpadas substituídas e da mão de obra, a
vantagem cresce ainda mais.\\
\textbf{Nota física:} a incandescente converte cerca de $95\%$ da energia em
calor e apenas $5\%$ em luz --- ela é, essencialmente, um aquecedor que
ilumina de lado.}
\end{questao}

\begin{questao}[3]{}
Uma instalação alimenta uma carga de $\SI{10}{\kilo\watt}$ por um cabo de
resistência total $\SI{0.2}{\ohm}$, em $\SI{220}{\volt}$.
\begin{itensq}
\item Determine a corrente e a perda no cabo.
\item Determine a fração da energia perdida.
\item Repita para $\SI{440}{\volt}$ e comente.
\end{itensq}
\resposta{(a) $\SI{45.5}{\ampere}$, $\SI{414}{\watt}$; (b) $4{,}1\%$;
(c) $\SI{103}{\watt}$, $1{,}0\%$}
\resolucao{(a) $I=\dfrac{10000}{220}=\SI{45.5}{\ampere}$ e
\[
P_{cabo}=RI^{2}=0{,}2(45{,}5)^{2}=\SI{414}{\watt}.
\]
(b) $414/10000=4{,}1\%$ da potência útil vira calor no cabo.\\
(c) Em $\SI{440}{\volt}$: $I=\SI{22.7}{\ampere}$ e
$P_{cabo}=0{,}2(22{,}7)^{2}=\SI{103}{\watt}$, ou $1{,}0\%$.\\
\textbf{Dobrar a tensão reduziu a perda a um quarto.} A razão é direta: para a
mesma potência, $I\propto1/U$, e a perda vai com $I^{2}$ --- logo
$P_{perda}\propto1/U^{2}$.\\
\textbf{É esta a razão de existir a alta tensão na transmissão.} Linhas de
$\SI{500}{\kilo\volt}$ transportam gigawatts com perdas de poucos por cento; nas
mesmas linhas em baixa tensão, a energia não chegaria ao destino. O preço é o
isolamento e os transformadores nas duas pontas.}
\end{questao}

\blocodesafio

\begin{questao}[4]{}
\textbf{(Demonstração --- potência média e valor eficaz.)} Prove que a potência
média em um resistor depende do valor \emph{eficaz} da corrente.
\begin{itensq}
\item Deduza $P_{\text{média}}=R\,I_{rms}^{2}$ a partir da definição.
\item Mostre que $I_{rms}\ge I_{\text{média}}$, com igualdade apenas para corrente
      constante.
\item Aplique a uma onda quadrada de $\SI{0}{\ampere}$ e
      $\SI{4}{\ampere}$ com ciclo de $\SI{25}{\percent}$.
\end{itensq}
\resposta{$I_{rms}=\SI{2}{\ampere}$ contra $I_{\text{média}}=\SI{1}{\ampere}$;
$P=R I_{rms}^{2}$}
\resolucao{(a) Por definição, a potência média em um período $T$ é
\[
P_{\text{média}}=\frac{1}{T}\int_0^{T}R\,i^{2}(t)\,\mathrm{d}t
=R\left[\frac{1}{T}\int_0^{T}i^{2}\,\mathrm{d}t\right]
=R\,I_{rms}^{2},
\]
onde a definição de valor eficaz é justamente
$I_{rms}=\sqrt{\frac1T\int_0^{T}i^{2}\mathrm{d}t}$ --- a
\emph{raiz} da \emph{média} do \emph{quadrado}.\\
(b) Pela desigualdade de Cauchy--Schwarz (ou pelo fato de a variância ser não
negativa):
\[
\left\langle i^{2}\right\rangle-\left\langle i\right\rangle^{2}
=\operatorname{Var}(i)\ge0
\;\Longrightarrow\;
I_{rms}^{2}\ge I_{\text{média}}^{2}.
\]
A igualdade vale se e somente se a variância é nula, isto é, se $i(t)$ é
\textbf{constante}.\\
\textbf{Leitura:} qualquer flutuação da corrente aumenta o aquecimento em
relação ao que a média sugeriria. Corrente ondulada aquece mais que corrente
contínua de mesmo valor médio --- resultado com consequências diretas em
retificadores e conversores.\\
(c) $\langle i\rangle=0{,}25(4)=\SI{1}{\ampere}$;
$\langle i^{2}\rangle=0{,}25(16)=\SI{4}{\ampere\squared}$, logo
$I_{rms}=\SI{2}{\ampere}$.\\
Com $R=\SI{10}{\ohm}$: $P=10(4)=\SI{40}{\watt}$, e não os
$\SI{10}{\watt}$ que a corrente média sugeriria. Fator $4$ --- exatamente
$(I_{rms}/I_{\text{média}})^{2}$.}
\end{questao}

\begin{questao}[4]{}
\textbf{(Integração --- energia em transitório.)} Um resistor $R$ conduz
$i(t)=I_0e^{-t/\tau}$.
\begin{itensq}
\item Deduza a energia total dissipada.
\item Determine a "potência média equivalente" ao longo de $5\tau$.
\item Compare com a energia de um pulso retangular de mesma amplitude e
      duração $\tau$.
\end{itensq}
\resposta{$E=\dfrac{RI_0^{2}\tau}{2}$; equivale à potência inicial mantida por
$\tau/2$}
\resolucao{(a)
\[
E=\int_0^{\infty}R\,I_0^{2}e^{-2t/\tau}\,\mathrm{d}t
=RI_0^{2}\left[-\frac{\tau}{2}e^{-2t/\tau}\right]_0^{\infty}
=\frac{RI_0^{2}\tau}{2}.
\]
\textbf{Observe o fator $\tau/2$, e não $\tau$:} a potência decai com constante
$\tau/2$ porque é proporcional a $i^{2}$, e o quadrado de uma exponencial tem
metade da constante de tempo.\\
(b) Distribuindo $E$ por $5\tau$:
\[
P_{\text{média}}=\frac{RI_0^{2}\tau/2}{5\tau}=\frac{RI_0^{2}}{10}=\frac{P_0}{10}.
\]
Apenas um décimo da potência inicial --- confirmando que o transitório é curto
diante da janela considerada.\\
(c) \textbf{Pulso retangular} de amplitude $I_0$ e duração $\tau$:
$E_{ret}=RI_0^{2}\tau$, exatamente o \textbf{dobro}.\\
\textbf{Consequência de dimensionamento:} substituir o decaimento real por um
pulso retangular equivalente de duração $\tau$ é uma aproximação
\emph{conservadora} --- ela superestima a energia por um fator $2$. Útil como
estimativa rápida e segura, desde que se saiba que há folga embutida.}
\end{questao}

\begin{questao}[4]{}
\textbf{(Projeto de proteção.)} Um resistor de $\SI{100}{\ohm}$ e
$\frac14\,\si{\watt}$ deve ser protegido contra sobrecarga.
\begin{itensq}
\item Determine a tensão e a corrente máximas admissíveis.
\item Projete um resistor limitador em série, para uma fonte de
      $\SI{24}{\volt}$.
\item Verifique a potência do limitador.
\end{itensq}
\resposta{$U_{\max}=\SI{5}{\volt}$, $I_{\max}=\SI{50}{\milli\ampere}$;
limitador de $\SI{380}{\ohm}$}
\resolucao{(a)
\[
U_{\max}=\sqrt{PR}=\sqrt{0{,}25(100)}=\SI{5}{\volt};
\qquad
I_{\max}=\sqrt{\frac{P}{R}}=\sqrt{0{,}0025}=\SI{50}{\milli\ampere}.
\]
(b) Com $\SI{24}{\volt}$ e a exigência $I\le\SI{50}{\milli\ampere}$:
\[
R_{total}\ge\frac{24}{0{,}050}=\SI{480}{\ohm}
\;\Longrightarrow\;
R_{lim}\ge480-100=\SI{380}{\ohm}.
\]
O valor E24 imediatamente superior é $\SI{390}{\ohm}$, que dá
$I=24/490=\SI{49}{\milli\ampere}$ \checkmark.\\
(c) $P_{lim}=390(0{,}049)^{2}=\SI{0.94}{\watt}$ --- especifique
$\SI{2}{\watt}$.\\
\textbf{Observação incômoda:} o limitador dissipa \emph{quatro vezes} mais que o
resistor protegido, e o conjunto consome $\SI{1.18}{\watt}$ para entregar
$\SI{0.24}{\watt}$ ao componente de interesse.\\
\textbf{Quando isso se justifica:} apenas se o resistor de
$\SI{100}{\ohm}$ tiver função específica (sensor, elemento de calibração) que o
torne insubstituível. Caso contrário, redimensionar o próprio resistor para uma
potência maior é mais simples e mais eficiente.}
\end{questao}

\begin{questao}[4]{}
\textbf{(Análise de retorno.)} Dois equipamentos entregam a mesma energia útil
de $\SI{1000}{\kilo\watt\hour}$ por ano: o modelo A tem rendimento de
$\SI{80}{\percent}$ e o B, de $\SI{92}{\percent}$.
\begin{itensq}
\item Determine a energia de entrada de cada um e a diferença.
\item Determine o tempo de retorno se B custar $\mathrm{R\$}\,800$ a mais, com
      tarifa de $\mathrm{R\$}\,0{,}95/\si{\kilo\watt\hour}$.
\item Discuta os fatores que a análise simples ignora.
\end{itensq}
\resposta{(a) $1250$ e $\SI{1087}{\kilo\watt\hour}$, diferença de
$\SI{163}{\kilo\watt\hour}$; (b) cerca de $\SI{5.2}{year}$}
\resolucao{(a)
\[
E_A=\frac{1000}{0{,}80}=\SI{1250}{\kilo\watt\hour};
\qquad
E_B=\frac{1000}{0{,}92}=\SI{1087}{\kilo\watt\hour};
\]
diferença de $\SI{163}{\kilo\watt\hour}$ por ano.\\
(b) Economia anual: $163(0{,}95)=\mathrm{R\$}\,154{,}89$.
\[
t=\frac{800}{154{,}89}=\SI{5.2}{year}.
\]
(c) \textbf{O que a conta simples ignora.}
\begin{itemize}
\item \textbf{Reajuste da tarifa:} se a energia subir acima da inflação, o
retorno acelera. Historicamente é o que ocorre.
\item \textbf{Custo do capital:} $\mathrm{R\$}\,800$ aplicados renderiam algo
--- o retorno real é mais longo que o nominal.
\item \textbf{Vida útil:} se o equipamento durar $\SI{10}{year}$, B gera lucro
nos últimos cinco. Se durar $\SI{4}{year}$, a troca não se paga.
\item \textbf{Custo do calor rejeitado:} A dissipa $\SI{250}{\kilo\watt\hour}$
contra $\SI{87}{\kilo\watt\hour}$ de B. Em ambiente climatizado, esse calor
extra precisa ser \emph{removido} --- o que pode dobrar a economia efetiva.
\item \textbf{Regime de uso:} rendimentos variam com a carga, e o valor de
catálogo costuma ser o de plena carga.
\end{itemize}
\textbf{Conclusão:} $\SI{5.2}{year}$ é aceitável para equipamento industrial de
vida longa e uso contínuo; é ruim para uso intermitente ou equipamento de
substituição rápida. \textbf{A resposta depende do horizonte, não apenas do
número.}}
\end{questao}

\begin{questao}[4]{}
\textbf{(Restrição orçamentária.)} Deseja-se limitar a
$\mathrm{R\$}\,90{,}00$ o custo mensal de um aquecedor de
$\SI{1.5}{\kilo\watt}$, com tarifa de
$\mathrm{R\$}\,1{,}00/\si{\kilo\watt\hour}$.
\begin{itensq}
\item Determine o tempo de uso admissível.
\item Determine a potência que permitiria uso de $\SI{4}{\hour}$ diárias dentro
      do mesmo orçamento.
\item Discuta qual das duas soluções é preferível.
\end{itensq}
\resposta{(a) $\SI{60}{\hour}$ por mês, ou $\SI{2}{\hour}$ por dia;
(b) $\SI{750}{\watt}$}
\resolucao{(a) O orçamento permite $\SI{90}{\kilo\watt\hour}$ por mês:
\[
t=\frac{90}{1{,}5}=\SI{60}{\hour\per month}
=\SI{2}{\hour\per day}.
\]
(b) Para $4\times30=\SI{120}{\hour}$ mensais dentro dos mesmos
$\SI{90}{\kilo\watt\hour}$:
\[
P=\frac{90}{120}=\SI{0.75}{\kilo\watt}=\SI{750}{\watt}.
\]
(c) \textbf{Depende do objetivo térmico.} As duas soluções entregam a
\emph{mesma energia} --- e portanto o mesmo calor total no mês.
\begin{itemize}
\item \textbf{$\SI{1.5}{\kilo\watt}$ por $\SI{2}{\hour}$:} aquece rápido e
atinge temperatura mais alta, mas o ambiente esfria nas outras
$\SI{22}{\hour}$.
\item \textbf{$\SI{750}{\watt}$ por $\SI{4}{\hour}$:} aquecimento mais suave e
prolongado, com temperatura média mais estável.
\end{itemize}
\textbf{Para conforto térmico, a segunda costuma ser superior} --- o ambiente
tem inércia térmica, e potência menor por mais tempo reduz as oscilações.\\
\textbf{E há uma terceira solução, melhor que ambas:} melhorar o
\emph{isolamento} do ambiente. Isso reduz a energia necessária, e não apenas a
redistribui --- atacando a causa em vez do sintoma.}
\end{questao}

\begin{questao}[4]{}
\textbf{(Otimização econômica de condutor.)} Um cabo alimenta uma carga de
$\SI{40}{\ampere}$ por $\SI{4000}{\hour}$ ao ano. O custo do cabo é
proporcional à seção, e a resistência é inversamente proporcional a ela.
\begin{itensq}
\item Escreva o custo total como função da seção.
\item Determine a condição de mínimo.
\item Interprete o resultado.
\end{itensq}
\resposta{No ótimo, o custo anualizado do cabo iguala o custo anual da energia
perdida}
\resolucao{(a) Sejam $S$ a seção, $k_1$ o custo anualizado do cabo por unidade
de seção e $R=\rho\ell/S$. O custo total anual é
\[
C(S)=\underbrace{k_1S}_{\text{investimento}}
+\underbrace{\frac{\rho\ell}{S}I^{2}\,t\,c_E}_{\text{energia perdida}}
=k_1S+\frac{k_2}{S},
\]
com $k_2=\rho\ell I^{2}t\,c_E$.\\
(b) Derivando e anulando:
\[
\frac{\mathrm{d}C}{\mathrm{d}S}=k_1-\frac{k_2}{S^{2}}=0
\;\Longrightarrow\;
S^{*}=\sqrt{\frac{k_2}{k_1}}.
\]
A segunda derivada, $2k_2/S^{3}>0$, confirma o mínimo.\\
(c) \textbf{Substituindo $S^{*}$ nas duas parcelas:}
\[
k_1S^{*}=\sqrt{k_1k_2}
\qquad\text{e}\qquad
\frac{k_2}{S^{*}}=\sqrt{k_1k_2}.
\]
As duas são \textbf{iguais} no ótimo --- resultado conhecido como \emph{lei de
Kelvin}: \textbf{a seção econômica é aquela em que o custo anualizado do
condutor iguala o custo anual das perdas}.\\
\textbf{Leituras práticas.}
\begin{itemize}
\item Cabos usados \emph{muitas horas} por ano ($t$ grande) justificam seções
maiores --- $S^{*}\propto\sqrt{t}$.
\item Correntes maiores também: $S^{*}\propto I$.
\item Energia mais cara empurra a seção para cima: $S^{*}\propto\sqrt{c_E}$.
\end{itemize}
\textbf{Importante:} a seção econômica é frequentemente \emph{maior} que a
mínima exigida pela norma (que garante apenas segurança térmica e queda de
tensão). Dimensionar pelo mínimo normativo é seguro, mas costuma ser
economicamente subótimo em instalações de uso intenso.}
\end{questao}

\begin{questao}[4]{}
\textbf{(Balanço energético de um chuveiro.)} Um chuveiro de
$\SI{5500}{\watt}$ aquece água de $\SI{20}{\celsius}$ a
$\SI{40}{\celsius}$.
\begin{itensq}
\item Determine a vazão máxima, supondo rendimento de $\SI{100}{\percent}$.
\item Determine o custo de um banho de $\SI{8}{\minute}$ a
      $\mathrm{R\$}\,0{,}95/\si{\kilo\watt\hour}$.
\item Discuta por que o rendimento é tão alto neste caso.
\end{itensq}
\dados{$c_{\text{água}}=\SI{4180}{\joule\per\kilogram\per\kelvin}$}
\resposta{(a) $\SI{0.0658}{\kilogram\per\second}\approx
\SI{3.9}{\liter\per\minute}$; (b) $\mathrm{R\$}\,0{,}70$}
\resolucao{(a) A potência necessária para aquecer uma vazão mássica $\dot m$
por $\Delta T=\SI{20}{\kelvin}$ é $P=\dot m\,c\,\Delta T$:
\[
\dot m=\frac{P}{c\,\Delta T}=\frac{5500}{4180(20)}
=\SI{0.0658}{\kilogram\per\second}\approx\SI{3.9}{\liter\per\minute}.
\]
Vazão modesta --- e é exatamente por isso que abrir demais o registro faz a água
esfriar: a potência é fixa, e mais vazão significa menor $\Delta T$.\\
(b)
\[
E=5{,}5\,\si{\kilo\watt}\times\frac{8}{60}\,\si{\hour}
=\SI{0.733}{\kilo\watt\hour}
\;\Longrightarrow\;
\text{custo}=0{,}733(0{,}95)=\mathrm{R\$}\,0{,}70.
\]
Um banho diário custa cerca de $\mathrm{R\$}\,21$ por mês --- e para uma
família de quatro pessoas, mais de $\mathrm{R\$}\,80$.\\
(c) \textbf{Por que o rendimento é quase $100\%$.} A resistência está
\emph{imersa} na água, e \textbf{toda} a energia dissipada por efeito Joule vira
calor exatamente onde ele é desejado. Não há conversão intermediária, nem calor
"desperdiçado" --- o que seria perda em outro contexto é aqui o produto.\\
\textbf{Contraste instrutivo:} uma lâmpada incandescente tem os mesmos
$\approx100\%$ de conversão em calor, mas rendimento \emph{luminoso} de $5\%$.
O rendimento só se define em relação ao \emph{objetivo}: a mesma física é
excelente para aquecer e péssima para iluminar.}
\end{questao}

\begin{questao}[4]{}
\textbf{(Diagnóstico por consumo.)} A conta de energia de uma residência subiu
de $\SI{180}{\kilo\watt\hour}$ para $\SI{320}{\kilo\watt\hour}$ mensais, sem
mudança de hábitos.
\begin{itensq}
\item Determine a potência contínua equivalente ao acréscimo.
\item Levante hipóteses compatíveis com esse valor.
\item Proponha um procedimento de investigação.
\end{itensq}
\resposta{Acréscimo de $\SI{140}{\kilo\watt\hour}$, equivalente a
$\SI{194}{\watt}$ contínuos}
\resolucao{(a) O mês tem cerca de $\SI{720}{\hour}$:
\[
P=\frac{140000}{720}=\SI{194}{\watt}\ \text{contínuos}.
\]
(b) \textbf{Hipóteses compatíveis com $\approx\SI{200}{\watt}$ permanentes.}
\begin{itemize}
\item \textbf{Fuga para o terra} por isolamento danificado --- especialmente em
resistência de chuveiro ou aquecedor de água. Uma fuga de
$\SI{0.88}{\ampere}$ em $\SI{220}{\volt}$ dá exatamente
$\SI{194}{\watt}$.
\item \textbf{Geladeira com vedação ruim ou gás insuficiente}, operando em
regime quase contínuo em vez de ciclado.
\item \textbf{Bomba ou motor} funcionando sem parar por falha de boia ou
pressostato.
\item \textbf{Ligação clandestina} a jusante do medidor.
\end{itemize}
(c) \textbf{Procedimento de investigação, em ordem.}
\begin{enumerate}
\item \textbf{Desligar o disjuntor geral} e verificar se o medidor continua
girando. Se sim, o problema está \emph{antes} do quadro --- ligação irregular ou
medidor defeituoso.
\item \textbf{Religar circuito por circuito}, observando o medidor a cada
etapa, para isolar o ramo responsável.
\item \textbf{Medir a corrente} de cada circuito suspeito com alicate
amperímetro, com todos os aparelhos desligados. Corrente não nula indica fuga.
\item \textbf{Testar o dispositivo DR}, se houver --- fuga acima de
$\SI{30}{\milli\ampere}$ deveria ter provocado desarme. A ausência de desarme
com fuga confirmada indica DR defeituoso, e é problema de segurança, não apenas
de conta.
\end{enumerate}
\textbf{Prioridade:} uma fuga de $\SI{0.88}{\ampere}$ é risco de choque e de
incêndio. O custo na conta é o sintoma menos grave.}
\end{questao}

% END BANCO COMPLEMENTAR

\gabarito[2.3 Potência elétrica e energia]

% =====================================================================
\subsection{Fontes reais e Leis de Kirchhoff}
% =====================================================================

\blocofixacao

\begin{questao}[1]{}
Um circuito tem duas fontes ideais de $\SI{6}{\volt}$ e $\SI{4}{\volt}$ ligadas em
série e em \textbf{oposição}.
\begin{itensq}
\item Qual é a tensão total fornecida ao circuito?
\item Se a resistência total for $\SI{5}{\ohm}$, qual é a corrente?
\end{itensq}
\resposta{(a) \SI{2}{\volt}; (b) \SI{0.4}{\ampere}}
\resolucao{(a) Em oposição, as fem se subtraem:
$U = 6 - 4 = \SI{2}{\volt}$.\\
(b) $I = \dfrac{U}{R} = \dfrac{2}{5} = \SI{0.4}{\ampere}$.
Se estivessem em série \emph{concordante}, teríamos $6+4 = \SI{10}{\volt}$ e
$I = \SI{2}{\ampere}$.}
\end{questao}

\begin{questao}[1]{}
No nó da figura, as correntes $I_1 = \SI{5}{\ampere}$ e $I_2 = \SI{3}{\ampere}$
entram, e $I_3$ sai. Pela lei dos nós de Kirchhoff, $I_3$ vale:

\circuitoq{%
\begin{circuitikz}[scale=1,transform shape,line width=0.8pt]
 \coordinate (o) at (0,0);
 \draw[->,Icol,line width=1.1pt] (-2,0.6) -- (o) node[midway,above]{$I_1$};
 \draw[->,Icol,line width=1.1pt] (-2,-0.6) -- (o) node[midway,below]{$I_2$};
 \draw[->,Icol,line width=1.1pt] (o) -- (2,0) node[midway,above]{$I_3$};
 \fill (o) circle (0.06);
\end{circuitikz}}

\resposta{$I_3 = \SI{8}{\ampere}$}
\resolucao{A lei dos nós afirma que a soma das correntes que entram é igual à das
que saem:
\[
I_3 = I_1 + I_2 = 5 + 3 = \SI{8}{\ampere}.
\]
Ela expressa a conservação da carga: o nó não acumula carga.}
\end{questao}

\blocoaplicacao

\begin{questao}[2]{}
Uma bateria de fem $\EMF = \SI{12}{\volt}$ e resistência interna
$r = \SI{0.5}{\ohm}$ alimenta um resistor externo $R = \SI{5.5}{\ohm}$. Determine
a corrente e a tensão nos terminais da bateria.
\resposta{$I = \SI{2}{\ampere}$; $U = \SI{11}{\volt}$}
\resolucao{A resistência interna está em série com a carga:
\[
I=\frac{\EMF}{R+r}=\frac{12}{5{,}5+0{,}5}=\SI{2}{\ampere}.
\]
A tensão nos terminais é a fem menos a queda interna:
\[
U=\EMF - r\,I = 12 - 0{,}5\cdot2 = \SI{11}{\volt}.
\]
A diferença de \SI{1}{\volt} "some" dentro da bateria --- é a queda em $r$. Quanto
maior a corrente, menor a tensão entregue: por isso baterias velhas (com $r$
elevado) fornecem menos tensão sob carga.}
\end{questao}

\begin{questao}[2]{}
Uma bateria de fem $\EMF$ e resistência interna $r = \SI{0.4}{\ohm}$ apresenta
tensão de $\SI{12}{\volt}$ em seus terminais quando fornece
$\SI{5}{\ampere}$. Determine a fem e a tensão em vazio (corrente nula).
\resposta{$\EMF = \SI{14}{\volt}$; em vazio, $U = \SI{14}{\volt}$}
\resolucao{Da equação da fonte real, $U = \EMF - r\,I$:
\[
\EMF = U + r\,I = 12 + 0{,}4\cdot5 = \SI{14}{\volt}.
\]
Em vazio ($I = 0$), não há queda interna: $U = \EMF = \SI{14}{\volt}$. A tensão de
circuito aberto de uma fonte é sempre igual à sua fem.}
\end{questao}

\begin{questao}[2]{}
No nó da figura, são conhecidas quatro das cinco correntes. Determine $I_5$,
indicando se entra ou sai do nó.

\circuitoq{%
\begin{tikzpicture}[scale=1,transform shape,line width=0.9pt,>=Stealth]
 \coordinate (o) at (0,0);
 \draw[->,Icol] (-2,0.8) -- (o) node[midway,above]{$\SI{6}{\ampere}$};
 \draw[->,Icol] (-2,-0.8) -- (o) node[midway,below]{$\SI{4}{\ampere}$};
 \draw[->,Icol] (o) -- (2,0.8) node[midway,above]{$\SI{3}{\ampere}$};
 \draw[->,Icol] (o) -- (2,-0.8) node[midway,below]{$\SI{5}{\ampere}$};
 \draw[->,Vcol,thick] (o) -- (0,-1.8) node[midway,right]{$I_5$};
 \fill (o) circle (0.07);
\end{tikzpicture}}{Nó com cinco ramos.}

\resposta{$I_5 = \SI{2}{\ampere}$, saindo do nó}
\resolucao{Pela lei dos nós, a soma das que entram é igual à soma das que saem.
Entram $6 + 4 = \SI{10}{\ampere}$; saem $3 + 5 = \SI{8}{\ampere}$ pelos ramos
conhecidos. Para equilibrar, $I_5$ deve \emph{sair} com:
\[
I_5 = 10 - 8 = \SI{2}{\ampere}.
\]
Como o sentido adotado na figura já era de saída, o valor positivo confirma essa
suposição.}
\end{questao}

\blocoaprofundamento

\begin{questao}[3]{}
No circuito de duas malhas da figura, aplique as leis de Kirchhoff para determinar
as correntes $I_1$, $I_2$ e a corrente $I_3$ no resistor central.

\circuitoq{%
\begin{circuitikz}[scale=1,transform shape,line width=0.8pt]
 % malha esquerda
 \draw (0,0) to[battery1,l_={$E_1=\SI{18}{\volt}$}] (0,3)
       to[R,l={$R_1=\SI{2}{\ohm}$}] (3,3) coordinate(top);
 \draw (top) to[R,l={$R_2=\SI{6}{\ohm}$},i>_={$I_3$}] (3,0) coordinate(bot) -- (0,0);
 % malha direita
 \draw (top) to[R,l={$R_3=\SI{6}{\ohm}$}] (6,3)
       to[battery1,l={$E_2=\SI{6}{\volt}$},invert] (6,0) -- (bot);
 \draw[->,Icol,line width=1pt] (0.6,3.4)--(1.5,3.4) node[midway,above,font=\footnotesize]{$I_1$};
 \draw[->,Icol,line width=1pt] (5.4,3.4)--(4.5,3.4) node[midway,above,font=\footnotesize]{$I_2$};
\end{circuitikz}}

\resposta{$I_1 = \SI{3}{\ampere}$, $I_2 = \SI{1}{\ampere}$, $I_3 = \SI{2}{\ampere}$}
\resolucao{\textbf{Lei dos nós} (no topo): $I_1 = I_2 + I_3$, logo $I_3 = I_1 - I_2$.\\
\textbf{Malha esquerda} ($E_1$, $R_1$, $R_2$):
\[
E_1 = R_1 I_1 + R_2 I_3 \;\Rightarrow\; 18 = 2I_1 + 6(I_1 - I_2).
\]
\textbf{Malha direita} ($E_2$, $R_3$, $R_2$):
\[
E_2 = R_3 I_2 + R_2 I_3 \;\Rightarrow\; 6 = 6I_2 + 6(I_1 - I_2)\cdot(-1)?
\]
Organizando o sistema (com $I_3 = I_1 - I_2$):
\[
\begin{cases}
8I_1 - 6I_2 = 18\\[2pt]
-6I_1 + 12I_2 = -6
\end{cases}
\;\Rightarrow\;
I_1 = \SI{3}{\ampere},\ I_2 = \SI{1}{\ampere}.
\]
Logo $I_3 = 3 - 1 = \SI{2}{\ampere}$.\\
\textbf{Verificação} pela malha esquerda: $2\cdot3 + 6\cdot2 = 6+12 =
\SI{18}{\volt} = E_1$. Resultado consistente.}
\end{questao}

\begin{questao}[3]{}
No circuito de duas malhas da figura, determine as três correntes de ramo
($I_1$, $I_2$, $I_3$) aplicando as leis de Kirchhoff.

\circuitoq{%
\begin{circuitikz}[scale=1,transform shape,line width=0.8pt]
 \draw (0,0) to[battery1,l_={$E_1=\SI{12}{\volt}$}] (0,3)
       to[R,l={$R_1=\SI{2}{\ohm}$}] (3,3) coordinate(top);
 \draw (top) to[R,l={$R_2=\SI{4}{\ohm}$},i>_={$I_3$}] (3,0) coordinate(bot) -- (0,0);
 \draw (top) to[R,l={$R_3=\SI{2}{\ohm}$}] (6,3)
       to[battery1,l={$E_2=\SI{6}{\volt}$}] (6,0) -- (bot);
 \draw[->,Icol,line width=1pt] (0.6,3.4)--(1.5,3.4) node[midway,above,font=\footnotesize]{$I_1$};
 \draw[->,Icol,line width=1pt] (5.4,3.4)--(4.5,3.4) node[midway,above,font=\footnotesize]{$I_2$};
\end{circuitikz}}

\resposta{$I_1 = \SI{2.4}{\ampere}$, $I_2 = \SI{-0.6}{\ampere}$, $I_3 = \SI{1.8}{\ampere}$}
\resolucao{\textbf{Nó superior:} $I_1 + I_2 = I_3$.\\
\textbf{Malha esquerda:} $E_1 = R_1 I_1 + R_2 I_3
\Rightarrow 12 = 2I_1 + 4I_3$.\\
\textbf{Malha direita:} $E_2 = R_3 I_2 + R_2 I_3
\Rightarrow 6 = 2I_2 + 4I_3$.\\
Resolvendo o sistema (com $I_3 = I_1 + I_2$):
\[
I_1 = \SI{2.4}{\ampere},\quad I_2 = \SI{-0.6}{\ampere},\quad
I_3 = \SI{1.8}{\ampere}.
\]
O sinal \emph{negativo} de $I_2$ significa que a corrente real nesse ramo é
\emph{oposta} ao sentido adotado: a fonte $E_2$, mais fraca, na verdade
\textbf{recebe} corrente (está sendo carregada) em vez de fornecer.\\
\textbf{Verificação} (malha esquerda): $2\cdot2{,}4 + 4\cdot1{,}8 = 4{,}8+7{,}2
= \SI{12}{\volt} = E_1$. Resultado consistente. Este é um bom lembrete: adotar um sentido
errado não invalida a solução --- o sinal negativo se encarrega de corrigi-lo.}
\end{questao}

\begin{questao}[3]{}
No circuito da figura, quatro fontes ideais de tensão estão interligadas, e os
terminais $a$, $b$ e $c$ estão \textbf{abertos} (nenhuma carga conectada).
Observe atentamente a \emph{polaridade} indicada dentro de cada fonte.

\circuitoq{%
\begin{circuitikz}[scale=1.05,transform shape,line width=0.9pt]
 % barramento vertical
 \coordinate (N1) at (0,4.2);      % no superior
 \coordinate (N2) at (0,2.2);      % no inferior (referencia)
 \draw (N1) -- (0,4.2);
 % fonte 5V vertical (+ em cima, - embaixo)
 \draw (N2) to[\fonteV,l_={$\SI{5}{\volt}$},invert] (N1);
 % ramo superior -> a  (4V com '-' a esquerda, '+' a direita)
 \draw (N1) -- (0.9,4.2) to[\fonteV,l={$\SI{4}{\volt}$}] (2.9,4.2)
       to[short,-o] (4.0,4.2) node[right,font=\Large]{$a$};
 % barramento inferior
 \draw (N2) -- (0,0.2);
 % ramo do meio -> b  (5V com '+' a esquerda, '-' a direita)
 \draw (0,1.4) -- (0.9,1.4) to[\fonteV,l_={$\SI{5}{\volt}$},invert] (2.9,1.4)
       to[short,-o] (4.0,1.4) node[right,font=\Large]{$b$};
 \fill (0,1.4) circle (0.07);
 % ramo inferior -> c  (3V com '-' a esquerda, '+' a direita)
 \draw (0,0.2) -- (0.9,0.2) to[\fonteV,l_={$\SI{3}{\volt}$}] (2.9,0.2)
       to[short,-o] (4.0,0.2) node[right,font=\Large]{$c$};
\end{circuitikz}}

\begin{itensq}
\item Determine $V_{ac}$.
\item Determine também $V_{ab}$ e $V_{bc}$, e verifique a consistência
      ($V_{ab}+V_{bc}=V_{ac}$).
\item Por que os valores \emph{não} dependem das resistências dos fios?
\end{itensq}
\resposta{(a) $V_{ac}=\SI{6}{\volt}$; (b) $V_{ab}=\SI{14}{\volt}$,
$V_{bc}=\SI{-8}{\volt}$; (c) ver comentário}
\resolucao{\textbf{Estratégia.} Como os terminais estão abertos, \emph{não circula
corrente} em nenhum ramo. Sem corrente, não há queda em resistência alguma: os
potenciais são determinados exclusivamente pelas fem, somadas com o \emph{sinal}
correto. Basta adotar uma referência e "caminhar" pelo circuito.

\medskip
\textbf{Passo 1 --- adotar a referência.} Seja o nó inferior do barramento
(ao qual se ligam os ramos de $b$ e $c$) o nó de referência: $V_{N_2}=0$.

\textbf{Passo 2 --- subir pela fonte de \SI{5}{\volt} vertical.} Ela tem o
terminal $+$ \emph{em cima} e $-$ embaixo. Indo de $N_2$ para $N_1$, entramos
pelo $-$ e saímos pelo $+$: \emph{ganhamos} \SI{5}{\volt}.
\[
V_{N_1}=V_{N_2}+5=\SI{5}{\volt}.
\]

\textbf{Passo 3 --- percorrer cada ramo até seu terminal.}
\begin{itemize}[leftmargin=1.7em,itemsep=2pt,topsep=3pt]
\item \textbf{Ramo de $a$} (fonte \SI{4}{\volt}, com $-$ à esquerda e $+$ à
      direita): indo de $N_1$ para $a$, entramos pelo $-$ e saímos pelo $+$,
      ganhando \SI{4}{\volt}:
      \[V_a = V_{N_1}+4 = 5+4=\SI{9}{\volt}.\]
\item \textbf{Ramo de $b$} (fonte \SI{5}{\volt}, com $+$ à esquerda e $-$ à
      direita): indo de $N_2$ para $b$, entramos pelo $+$ e saímos pelo $-$,
      \emph{perdendo} \SI{5}{\volt}:
      \[V_b = V_{N_2}-5 = \SI{-5}{\volt}.\]
\item \textbf{Ramo de $c$} (fonte \SI{3}{\volt}, com $-$ à esquerda e $+$ à
      direita): ganhamos \SI{3}{\volt}:
      \[V_c = V_{N_2}+3 = \SI{3}{\volt}.\]
\end{itemize}

\textbf{Passo 4 --- calcular as diferenças.}
\[
\boxed{\;V_{ac}=V_a-V_c = 9-3 = \SI{6}{\volt}\;}
\]
\[
V_{ab}=V_a-V_b = 9-(-5)=\SI{14}{\volt},
\qquad
V_{bc}=V_b-V_c = -5-3=\SI{-8}{\volt}.
\]
\textbf{Consistência:} $V_{ab}+V_{bc} = 14+(-8) = \SI{6}{\volt} = V_{ac}$
\checkmark --- como deve ser, pois a diferença de potencial é uma função de
estado (independe do caminho).

\medskip
\textbf{(c) Por que os fios não importam?} A queda em um resistor é $U=RI$. Com
$I=0$ em todos os ramos (terminais abertos), \emph{qualquer} resistência produz
queda nula. Os potenciais dos terminais dependem só das fem e de suas
polaridades.\\
\textbf{Atenção ao erro mais comum:} somar as fontes ignorando a polaridade
(obtendo, por exemplo, $4+5+3=12$). O sinal de cada fonte é determinado pelo
\emph{sentido em que a percorremos}: entrando pelo $-$ e saindo pelo $+$, soma-se;
o contrário, subtrai-se. \emph{Sempre marque a polaridade antes de escrever a
primeira equação.}}
\end{questao}

\blocodesafio

\begin{questao}[4]{}
Duas baterias de fem $\EMF_1 = \SI{12}{\volt}$ ($r_1=\SI{0.2}{\ohm}$) e
$\EMF_2 = \SI{12.6}{\volt}$ ($r_2=\SI{0.1}{\ohm}$) são ligadas em
\textbf{paralelo} para alimentar uma carga $R_L = \SI{2}{\ohm}$ --- situação
comum em bancos de baterias e em sistemas com gerador e bateria em paralelo.

\circuitoq{%
\begin{circuitikz}[scale=0.95,transform shape,line width=0.8pt]
 \draw (0,0) to[battery1,l_={$\EMF_1$}] (0,1.6) to[R,l_={$r_1$},color=Vcol] (0,3);
 \draw (2.4,0) to[battery1,l_={$\EMF_2$}] (2.4,1.6) to[R,l_={$r_2$},color=Vcol] (2.4,3);
 \draw (0,3) -- (2.4,3) -- (5,3) to[R,l={$R_L=\SI{2}{\ohm}$}] (5,0) -- (2.4,0) -- (0,0);
 \node[above,Vcol] at (3.6,3.1) {$V$};
 \fill (2.4,3) circle (0.07); \fill (2.4,0) circle (0.07);
\end{circuitikz}}

\begin{itensq}
\item Determine a tensão $V$ sobre a carga (use análise nodal).
\item Calcule a corrente fornecida por \emph{cada} bateria.
\item Interprete o sinal obtido e explique o risco prático dessa associação.
\end{itensq}
\resposta{(a) $V = \SI{12}{\volt}$; (b) $I_1 = 0$, $I_2 = \SI{6}{\ampere}$;
(c) ver comentário}
\resolucao{\textbf{(a)} Tomando o nó superior com potencial $V$ e o inferior como
terra, a lei dos nós exige que a soma das correntes que \emph{chegam} ao nó seja
igual à que sai pela carga:
\[
\frac{\EMF_1-V}{r_1}+\frac{\EMF_2-V}{r_2}=\frac{V}{R_L}.
\]
Substituindo:
\[
\frac{12-V}{0{,}2}+\frac{12{,}6-V}{0{,}1}=\frac{V}{2}.
\]
Multiplicando por 2: $10(12-V)+20(12{,}6-V)=V$, ou seja,
\[
120-10V+252-20V=V
\;\Longrightarrow\;
372 = 31V
\;\Longrightarrow\;
V = \SI{12}{\volt}.
\]

\textbf{(b)} As correntes de cada fonte:
\[
I_1=\frac{12-12}{0{,}2}=\SI{0}{\ampere},
\qquad
I_2=\frac{12{,}6-12}{0{,}1}=\SI{6}{\ampere},
\qquad
I_L=\frac{12}{2}=\SI{6}{\ampere}.\ \checkmark
\]
\textbf{(c) Interpretação e risco.} A bateria de \emph{maior} fem assume
praticamente toda a carga, enquanto a de menor fem fica ociosa. Se a diferença de
fem fosse maior --- por exemplo, uma bateria descarregada em paralelo com uma
carregada --- a corrente na bateria mais fraca ficaria \textbf{negativa}: ela
seria \emph{carregada} pela outra, com corrente limitada apenas pelas
resistências internas, tipicamente muito pequenas.\\
\textbf{Consequência prática:} correntes de circulação elevadíssimas, aquecimento
e risco de dano ou explosão. Por isso o paralelismo de baterias exige unidades
\emph{idênticas em fem, tecnologia e estado de carga}, ou o uso de diodos de
isolamento e controladores de balanceamento --- prática obrigatória em bancos de
baterias e sistemas fotovoltaicos.}
\end{questao}

% BEGIN BANCO COMPLEMENTAR Fontes reais e Leis de Kirchhoff
% =====================================================================
%  Banco complementar --- Fontes Reais e Leis de Kirchhoff  (REVISADO)
%  Substitui a versao gerada por template (9 clones em N2, 9 em N3 e
%  7 questoes conceituais marcadas como N4).
%  O cap02 ja cobre: duas fontes ideais em oposicao, LKC em no com
%  I1 e I2 entrando, bateria 12 V / 0,5 ohm, bateria com r=0,4 ohm,
%  no com cinco correntes, DOIS circuitos de duas malhas resolvidos
%  por Kirchhoff, quatro fontes ideais com terminais abertos e --- como
%  N4 --- DUAS BATERIAS DE FEM DIFERENTES EM PARALELO. Este ultimo tema
%  motivou a remocao de uma N4 do banco antigo, que o duplicava.
%  O banco de circuitos serie revisado ja traz regulacao de tensao e a
%  determinacao de fem e r por dois pontos de carga.
%  Este banco explora: fonte ABSORVENDO potencia, banco de baterias
%  serie x paralelo, LKC generalizada, contagem de equacoes
%  independentes, propagacao de incerteza na verificacao da LKC e o
%  teorema de Tellegen em forma elementar.
%  Distribuicao mantida: 6 N1 + 9 N2 + 9 N3 + 7 N4 = 31
% =====================================================================

\subsubsection*{Banco complementar --- Fontes reais e leis de Kirchhoff}

\begin{atencao}[Organização do banco]
A fonte real é o modelo $\EMF$ em série com $r$, e dele decorre
$V=\EMF-rI$ quando ela \emph{fornece} --- mas $V=\EMF+rI$ quando ela
\textbf{recebe} corrente, situação de carregamento que aparece com frequência e
costuma ser esquecida. As leis de Kirchhoff, por sua vez, decorrem da
conservação de carga e de energia: elas valem para \emph{qualquer} elemento,
linear ou não. O N3 explora fontes absorvendo potência, bancos de baterias e
contagem de equações; o N4 exige demonstração, ensaio e propagação de
incertezas.
\end{atencao}

\blocofixacao

\begin{questao}[1]{}
Uma fonte de $\EMF=\SI{6}{\volt}$ e $r=\SI{1}{\ohm}$ fornece
$\SI{1}{\ampere}$. Determine a tensão em seus terminais.
\resposta{$V=\SI{5}{\volt}$}
\resolucao{
\[
V=\EMF-rI=6-1(1)=\SI{5}{\volt}.
\]
A diferença de $\SI{1}{\volt}$ é a queda \emph{interna}, que não chega à carga.
Ela cresce com a corrente --- é por isso que a tensão de uma pilha cai quando se
liga um aparelho que exige mais.}
\end{questao}

\begin{questao}[1]{}
Uma fonte de $\EMF=\SI{9}{\volt}$ apresenta $\SI{8.4}{\volt}$ nos terminais
quando fornece $\SI{2}{\ampere}$. Determine sua resistência interna.
\resposta{$r=\SI{0.3}{\ohm}$}
\resolucao{
\[
r=\frac{\EMF-V}{I}=\frac{9-8{,}4}{2}=\SI{0.3}{\ohm}.
\]
\textbf{Método geral:} a queda interna dividida pela corrente. Uma única medição
sob carga, somada ao valor de $\EMF$, determina $r$ --- e $\EMF$ pode ser lida
em vazio com um voltímetro de alta impedância.}
\end{questao}

\begin{questao}[1]{}
Em um nó entram $\SI{2.5}{\ampere}$ e $\SI{1.5}{\ampere}$, e sai uma única
corrente. Determine seu valor.
\resposta{$I=\SI{4}{\ampere}$}
\resolucao{Pela lei dos nós,
\[
\sum I_{entra}=\sum I_{sai}
\;\Longrightarrow\;
2{,}5+1{,}5=\SI{4}{\ampere}.
\]
\textbf{Fundamento:} a LKC expressa a \textbf{conservação da carga}. Um nó não
acumula carga --- tudo o que entra, sai. Isso vale independentemente dos
elementos ligados a ele, e mesmo que sejam não lineares.}
\end{questao}

\begin{questao}[1]{}
Uma malha simples contém uma fonte ideal de $\SI{12}{\volt}$ e resistores de
$\SI{4}{\ohm}$ e $\SI{8}{\ohm}$. Determine a corrente pela lei das malhas.
\resposta{$I=\SI{1}{\ampere}$}
\resolucao{Pela LKT, a soma das quedas iguala a f.e.m.:
\[
12=4I+8I=12I
\;\Longrightarrow\;
I=\SI{1}{\ampere}.
\]
\textbf{Fundamento:} a LKT expressa a \textbf{conservação da energia}. Uma carga
que percorre a malha e retorna ao ponto de partida deve ter ganho e perdido a
mesma energia --- o potencial é uma função de estado.}
\end{questao}

\begin{questao}[1]{}
Sobre a diferença entre fonte ideal e fonte real, é correto afirmar:
\begin{alternativas}[2]
\item a ideal mantém a tensão constante para qualquer corrente
\item a real mantém a tensão constante para qualquer corrente
\item ambas mantêm a corrente constante
\item a ideal tem resistência interna maior
\end{alternativas}
\resposta{(a)}
\resolucao{A fonte \textbf{ideal} tem $r=0$: sua tensão terminal é
$\EMF$ independentemente da corrente, inclusive em curto-circuito --- o que
exigiria corrente infinita, revelando que se trata de uma idealização.\\
A fonte \textbf{real} tem $r>0$, e sua tensão cai como $\EMF-rI$.\\
\textbf{Quando a idealização é aceitável:} sempre que $r\ll R_{carga}$. Uma
fonte de bancada com $r=\SI{10}{\milli\ohm}$ alimentando
$\SI{100}{\ohm}$ é, para todos os efeitos, ideal.}
\end{questao}

\begin{questao}[1]{}
Uma pilha de $\EMF=\SI{1.5}{\volt}$ tem $r=\SI{0.25}{\ohm}$. Determine sua
corrente de curto-circuito.
\resposta{$I_{cc}=\SI{6}{\ampere}$}
\resolucao{
\[
I_{cc}=\frac{\EMF}{r}=\frac{1{,}5}{0{,}25}=\SI{6}{\ampere}.
\]
\textbf{Ressalva importante:} o valor pressupõe $r$ constante --- hipótese
otimista. Em uma pilha real, $r$ cresce com a corrente e com o aquecimento, e a
corrente efetiva é menor. Ainda assim, $\SI{6}{\ampere}$ em uma pilha pequena
significa $\SI{9}{\watt}$ dissipados \emph{dentro} dela: aquecimento rápido,
vazamento de eletrólito e risco de ruptura.}
\end{questao}

\blocoaplicacao

\begin{questao}[2]{}
Uma bateria de $\EMF=\SI{12}{\volt}$ e $r=\SI{0.6}{\ohm}$ alimenta
$R=\SI{5.4}{\ohm}$.
\begin{itensq}
\item Determine a corrente e a tensão terminal.
\item Determine as potências e o rendimento.
\end{itensq}
\resposta{$I=\SI{2}{\ampere}$, $V=\SI{10.8}{\volt}$;
$P_R=\SI{21.6}{\watt}$, $P_r=\SI{2.4}{\watt}$, $\eta=90\%$}
\resolucao{(a)
\[
I=\frac{\EMF}{r+R}=\frac{12}{6}=\SI{2}{\ampere};
\qquad
V=12-0{,}6(2)=\SI{10.8}{\volt}.
\]
(b)
\[
P_{total}=\EMF\,I=\SI{24}{\watt};
\quad
P_R=5{,}4(4)=\SI{21.6}{\watt};
\quad
P_r=0{,}6(4)=\SI{2.4}{\watt};
\]
\[
\eta=\frac{R}{R+r}=\frac{5{,}4}{6}=90\%.
\]
\textbf{Fórmula útil:} o rendimento de uma fonte real é sempre
$R/(R+r)$ --- ele depende só da razão entre as resistências, e não da tensão
nem da corrente.}
\end{questao}

\begin{questao}[2]{}
Um carregador de $\SI{14}{\volt}$ é ligado a uma bateria de
$\SI{12}{\volt}$, com resistência total de $\SI{0.4}{\ohm}$ no circuito.
\begin{itensq}
\item Determine a corrente de carga.
\item Determine a tensão nos terminais da bateria.
\end{itensq}
\resposta{$I=\SI{5}{\ampere}$; $V=\SI{14}{\volt}$}
\resolucao{(a) As duas f.e.m. estão em \textbf{oposição}, e a diferença é que
impulsiona a corrente:
\[
I=\frac{14-12}{0{,}4}=\SI{5}{\ampere}.
\]
(b) A bateria está \textbf{recebendo} corrente, então a queda interna
\emph{soma-se} à f.e.m.:
\[
V=\EMF+rI=12+0{,}4(5)=\SI{14}{\volt}.
\]
\textbf{Inversão do sinal:} quando a fonte fornece, $V=\EMF-rI<\EMF$; quando
recebe, $V=\EMF+rI>\EMF$. É por isso que a tensão de uma bateria de automóvel
sobe de $\SI{12.6}{\volt}$ para cerca de $\SI{14}{\volt}$ com o motor ligado ---
sinal de que o alternador está carregando.}
\end{questao}

\begin{questao}[2]{}
Em um nó, entram $\SI{3}{\ampere}$, $\SI{4}{\ampere}$ e $\SI{2}{\ampere}$;
sai $\SI{6}{\ampere}$ por um ramo e $I_5$ por outro. Determine $I_5$.
\resposta{$I_5=\SI{3}{\ampere}$, saindo}
\resolucao{
\[
3+4+2=6+I_5
\;\Longrightarrow\;
I_5=9-6=\SI{3}{\ampere}.
\]
O sinal positivo confirma o sentido admitido (saindo).\\
\textbf{Convenção alternativa:} adotar todas as correntes como \emph{entrando} e
impor soma nula. As que realmente saem aparecem com sinal negativo. As duas
convenções são equivalentes --- escolha uma e mantenha-a.}
\end{questao}

\begin{questao}[2]{}
Uma malha contém fontes ideais de $\SI{24}{\volt}$ e $\SI{6}{\volt}$ em
oposição, com resistores somando $\SI{6}{\ohm}$. Determine a corrente e
identifique qual fonte absorve energia.
\resposta{$I=\SI{3}{\ampere}$; a de $\SI{6}{\volt}$ absorve}
\resolucao{
\[
I=\frac{24-6}{6}=\SI{3}{\ampere},
\]
no sentido imposto pela fonte maior.\\
\textbf{Balanço:} a de $\SI{24}{\volt}$ fornece $24(3)=\SI{72}{\watt}$; a de
$\SI{6}{\volt}$ \emph{absorve} $6(3)=\SI{18}{\watt}$; os resistores dissipam
$6(9)=\SI{54}{\watt}$. Soma: $18+54=\SI{72}{\watt}$ \checkmark.\\
\textbf{Critério para identificar:} a fonte \emph{absorve} quando a corrente
entra pelo seu terminal positivo. Se ela for recarregável, está sendo
carregada; se não for (uma pilha comum), a energia vira calor e o dispositivo
pode ser danificado.}
\end{questao}

\begin{questao}[2]{}
Duas pilhas de $\SI{1.5}{\volt}$ e $r=\SI{0.2}{\ohm}$ cada são associadas.
Determine a f.e.m. e a resistência interna equivalentes em série e em paralelo.
\resposta{Série: $\SI{3}{\volt}$ e $\SI{0.4}{\ohm}$; paralelo:
$\SI{1.5}{\volt}$ e $\SI{0.1}{\ohm}$}
\resolucao{\textbf{Série:} as f.e.m. se somam e as resistências internas também:
$\EMF=\SI{3}{\volt}$, $r=\SI{0.4}{\ohm}$.\\
\textbf{Paralelo} (pilhas \emph{idênticas}): a f.e.m. permanece
$\SI{1.5}{\volt}$ e as resistências internas ficam em paralelo:
$r=\SI{0.1}{\ohm}$.\\
\textbf{Correntes de curto:} série $3/0{,}4=\SI{7.5}{\ampere}$; paralelo
$1{,}5/0{,}1=\SI{15}{\ampere}$.\\
\textbf{Advertência:} a associação em paralelo só é segura com pilhas
\emph{idênticas e no mesmo estado de carga}. Com f.e.m. diferentes, circula
corrente entre elas mesmo sem carga externa.}
\end{questao}

\begin{questao}[2]{}
Determine o número de equações independentes de LKC e de LKT para um circuito
com $n=5$ nós e $b=8$ ramos.
\resposta{$4$ de LKC e $4$ de LKT, totalizando $8$}
\resolucao{\textbf{LKC:} $n-1=4$ equações. A do quinto nó é combinação linear
das anteriores --- redundante.\\
\textbf{LKT:} $b-n+1=8-5+1=4$ equações, uma por malha independente.\\
\textbf{Total:} $4+4=8=b$ --- exatamente o número de correntes de ramo a
determinar. O sistema é quadrado e bem posto.\\
\textbf{Não é coincidência:} $(n-1)+(b-n+1)=b$ sempre. As leis de Kirchhoff
fornecem \emph{precisamente} as equações necessárias, nem uma a mais nem a
menos.}
\end{questao}

\begin{questao}[2]{}
Uma fonte de $\EMF=\SI{12}{\volt}$ e $r=\SI{0.5}{\ohm}$ alimenta duas cargas
em paralelo, de $\SI{12}{\ohm}$ e $\SI{6}{\ohm}$. Determine a tensão terminal e
a corrente de cada carga.
\resposta{$V=\SI{10.67}{\volt}$; $\SI{0.89}{\ampere}$ e
$\SI{1.78}{\ampere}$}
\resolucao{
\[
R_{par}=\frac{12(6)}{18}=\SI{4}{\ohm};
\qquad
I=\frac{12}{4{,}5}=\SI{2.67}{\ampere};
\qquad
V=12-0{,}5(2{,}67)=\SI{10.67}{\volt}.
\]
\emph{(Conferindo: $V=R_{par}I=4(2{,}67)=\SI{10.67}{\volt}$ \checkmark.)}
\[
I_{12}=\frac{10{,}67}{12}=\SI{0.89}{\ampere};
\qquad
I_{6}=\frac{10{,}67}{6}=\SI{1.78}{\ampere}.
\]
\textbf{Acoplamento pelas cargas:} se a carga de $\SI{6}{\ohm}$ for desligada,
a corrente total cai para $12/12{,}5=\SI{0.96}{\ampere}$ e a tensão sobe para
$\SI{11.52}{\volt}$ --- o que \emph{aumenta} a corrente da carga remanescente.
As cargas se influenciam através de $r$.}
\end{questao}

\begin{questao}[2]{}
Uma bateria em uso apresenta $\EMF=\SI{1.5}{\volt}$ e $r=\SI{0.3}{\ohm}$
quando nova, e $\EMF=\SI{1.3}{\volt}$ com $r=\SI{1.2}{\ohm}$ ao final da vida.
Compare a tensão entregue a uma carga de $\SI{3}{\ohm}$.
\resposta{$\SI{1.36}{\volt}$ contra $\SI{0.93}{\volt}$ --- queda de $32\%$}
\resolucao{\textbf{Nova:} $I=1{,}5/3{,}3=\SI{0.455}{\ampere}$ e
$V=3(0{,}455)=\SI{1.36}{\volt}$.\\
\textbf{Gasta:} $I=1{,}3/4{,}2=\SI{0.310}{\ampere}$ e
$V=3(0{,}310)=\SI{0.93}{\volt}$.\\
\textbf{Observe o que domina:} a f.e.m. caiu apenas $13\%$, mas a resistência
interna \textbf{quadruplicou} --- e é ela a responsável pela maior parte da
perda de desempenho.\\
\textbf{Consequência para diagnóstico:} medir a tensão de uma pilha \emph{em
vazio} é enganoso, pois ela mostra quase $\EMF$ mesmo já esgotada. O teste
correto é sob carga, ou a medição direta de $r$.}
\end{questao}

\begin{questao}[2]{}
Verifique o balanço de potência de um circuito com fonte de
$\EMF=\SI{20}{\volt}$, $r=\SI{1}{\ohm}$ e carga de $\SI{9}{\ohm}$.
\resposta{Fornecido $\SI{40}{\watt}$; dissipado $4+\SI{36}{\watt}$}
\resolucao{
\[
I=\frac{20}{10}=\SI{2}{\ampere};
\qquad
P_{fonte}=\EMF\,I=\SI{40}{\watt};
\]
\[
P_r=1(4)=\SI{4}{\watt};
\qquad
P_R=9(4)=\SI{36}{\watt};
\qquad 4+36=40\ \checkmark
\]
\textbf{Cuidado com a distinção:} a potência \emph{gerada} usa $\EMF$, não a
tensão terminal. Usando $V=\SI{18}{\volt}$ obter-se-ia $\SI{36}{\watt}$ ---
que é a potência \emph{entregue aos terminais}, e não a produzida internamente.
A diferença de $\SI{4}{\watt}$ é exatamente o que $r$ consome.}
\end{questao}

\blocoaprofundamento

\begin{questao}[3]{}
Duas fontes reais alimentam a mesma carga: $\EMF_1=\SI{12}{\volt}$ com
$r_1=\SI{1}{\ohm}$ e $\EMF_2=\SI{6}{\volt}$ com $r_2=\SI{1}{\ohm}$, ambas
ligadas a $R=\SI{4}{\ohm}$.
\begin{itensq}
\item Determine a tensão comum e as três correntes.
\item Identifique o comportamento de cada fonte.
\end{itensq}
\resposta{$V=\SI{8}{\volt}$; $I_1=\SI{4}{\ampere}$,
$I_2=\SI{-2}{\ampere}$, $I_R=\SI{2}{\ampere}$ --- a fonte 2 é
\textbf{carregada}}
\resolucao{(a) LKC no nó comum:
\[
\frac{V-12}{1}+\frac{V-6}{1}+\frac{V}{4}=0
\;\Longrightarrow\;
4V-48+4V-24+V=0
\;\Longrightarrow\;
V=\SI{8}{\volt}.
\]
\[
I_1=\frac{12-8}{1}=\SI{4}{\ampere};
\qquad
I_2=\frac{6-8}{1}=\SI{-2}{\ampere};
\qquad
I_R=\frac{8}{4}=\SI{2}{\ampere}.
\]
\textbf{LKC:} $4+(-2)=2$ \checkmark.\\
(b) \textbf{A fonte 2 tem corrente negativa} --- ela não fornece, \emph{recebe}
$\SI{2}{\ampere}$. A fonte 1, mais forte, alimenta a carga \emph{e} carrega a
fonte 2.\\
\textbf{Balanço completo:}
\begin{center}
\begin{tabular}{lr}
\hline
Fonte 1 fornece & $\SI{48}{\watt}$\\
$r_1$ dissipa & $\SI{16}{\watt}$\\
$R$ dissipa & $\SI{16}{\watt}$\\
Fonte 2 absorve (f.e.m.) & $\SI{12}{\watt}$\\
$r_2$ dissipa & $\SI{4}{\watt}$\\
\hline
\end{tabular}
\end{center}
Total consumido: $16+16+12+4=\SI{48}{\watt}$ \checkmark.\\
\textbf{Se a fonte 2 for recarregável}, os $\SI{12}{\watt}$ são armazenados. Se
não for, viram calor --- e o dispositivo aquece perigosamente.}
\end{questao}

\begin{questao}[3]{}
Um banco é formado por quatro pilhas de $\SI{1.5}{\volt}$ e
$r=\SI{0.2}{\ohm}$.
\begin{itensq}
\item Compare série e paralelo quanto a $\EMF$, $r$ e corrente de
      curto.
\item Determine a potência máxima disponível em cada arranjo.
\item Estabeleça o critério de escolha.
\end{itensq}
\resposta{Série: $\SI{6}{\volt}$/$\SI{0.8}{\ohm}$; paralelo:
$\SI{1.5}{\volt}$/$\SI{0.05}{\ohm}$; $P_{\max}=\SI{11.25}{\watt}$ nos
\textbf{dois}}
\resolucao{(a)
\begin{center}
\begin{tabular}{lccc}
\hline
Arranjo & $\EMF$ & $r$ & $I_{cc}$\\
\hline
Série & $\SI{6}{\volt}$ & $\SI{0.8}{\ohm}$ & $\SI{7.5}{\ampere}$\\
Paralelo & $\SI{1.5}{\volt}$ & $\SI{0.05}{\ohm}$ & $\SI{30}{\ampere}$\\
\hline
\end{tabular}
\end{center}
(b)
\[
\text{Série: } \frac{6^{2}}{4(0{,}8)}=\SI{11.25}{\watt};
\qquad
\text{Paralelo: } \frac{1{,}5^{2}}{4(0{,}05)}=\SI{11.25}{\watt}.
\]
\textbf{Idênticas} --- e não por acaso. Ambos os arranjos usam as mesmas quatro
pilhas, e a energia total disponível é a mesma. A topologia decide \emph{como}
ela é entregue, não \emph{quanto}.\\
(c) \textbf{Critério de escolha.}
\begin{itemize}
\item \textbf{Série} para cargas de \emph{alta} resistência, que precisam de
tensão. A carga ótima é $\SI{0.8}{\ohm}$.
\item \textbf{Paralelo} para cargas de \emph{baixa} resistência, que precisam de
corrente. A carga ótima é $\SI{0.05}{\ohm}$.
\item \textbf{Série-paralelo} (dois ramos de duas) daria $\SI{3}{\volt}$ e
$\SI{0.2}{\ohm}$ --- e, de novo, $\SI{11.25}{\watt}$.
\end{itemize}
\textbf{Regra:} escolha o arranjo cuja resistência interna se aproxime da carga
--- ou, mais realisticamente, cuja \emph{tensão} atenda ao que a carga exige.}
\end{questao}

\begin{questao}[3]{}
Aplique a LKC a uma \textbf{superfície fechada} que envolva dois nós de um
circuito, e explique por que o resultado é válido.
\begin{itensq}
\item Enuncie a forma generalizada.
\item Justifique a partir da conservação de carga.
\item Dê um exemplo de uso.
\end{itensq}
\resposta{A soma das correntes que atravessam qualquer superfície fechada é
nula}
\resolucao{(a) \textbf{LKC generalizada:} para qualquer superfície fechada que
corte o circuito, a soma algébrica das correntes que a atravessam é zero.\\
(b) \textbf{Justificativa:} a carga total dentro da superfície não pode variar
(em regime permanente). Como corrente é fluxo de carga, o que entra deve
igualar o que sai --- exatamente o mesmo argumento que vale para um nó, agora
aplicado a uma região que \emph{contém} vários nós.\\
Formalmente, escrevendo a LKC para cada nó interno e somando, as correntes dos
ramos \emph{internos} aparecem duas vezes com sinais opostos e se cancelam;
sobram apenas as que cruzam a fronteira.\\
(c) \textbf{Exemplo de uso.} Considere um bloco de circuito ligado ao restante
por três fios apenas. Se dois deles conduzem $\SI{5}{\ampere}$ e
$\SI{3}{\ampere}$ para dentro, o terceiro conduz necessariamente
$\SI{8}{\ampere}$ para fora --- \textbf{sem que seja preciso conhecer nada do
interior do bloco}, que pode conter dezenas de componentes.\\
\textbf{Aplicação prática:} é o que permite verificar a corrente de um circuito
inteiro medindo apenas os condutores de entrada, e é também o princípio do
\emph{dispositivo diferencial residual}, que compara fase e neutro --- qualquer
diferença indica corrente escapando por um terceiro caminho.}
\end{questao}

\begin{questao}[3]{}
Uma fonte foi ensaiada sob diferentes cargas, obtendo-se:
\begin{center}
\begin{tabular}{cccc}
\hline
$I$ ($\si{\ampere}$) & $0{,}50$ & $1{,}00$ & $2{,}00$\\
$V$ ($\si{\volt}$) & $11{,}75$ & $11{,}50$ & $11{,}00$\\
\hline
\end{tabular}
\end{center}
\begin{itensq}
\item Determine $\EMF$ e $r$.
\item Verifique a linearidade do modelo.
\item Estime a corrente de curto e discuta a extrapolação.
\end{itensq}
\resposta{$\EMF=\SI{12}{\volt}$, $r=\SI{0.5}{\ohm}$;
$I_{cc}=\SI{24}{\ampere}$}
\resolucao{(a) O modelo $V=\EMF-rI$ é uma reta. De dois pontos quaisquer:
\[
r=-\frac{\Delta V}{\Delta I}=-\frac{11{,}00-11{,}75}{2{,}00-0{,}50}
=\frac{0{,}75}{1{,}50}=\SI{0.5}{\ohm};
\]
\[
\EMF=11{,}75+0{,}5(0{,}50)=\SI{12}{\volt}.
\]
(b) \textbf{Verificação com o terceiro ponto:}
$12-0{,}5(1{,}00)=\SI{11.50}{\volt}$ \checkmark. Os três pontos são
\emph{exatamente} colineares --- o modelo linear descreve bem a fonte nessa
faixa.\\
\textbf{Nota de método:} com apenas dois pontos, sempre passa uma reta e a
linearidade não pode ser testada. O terceiro ponto é o que transforma o ajuste
em \emph{verificação}.\\
(c) $I_{cc}=\EMF/r=12/0{,}5=\SI{24}{\ampere}$.\\
\textbf{Cautela com a extrapolação:} o ensaio cobriu de $0{,}5$ a
$\SI{2}{\ampere}$, e o valor de curto está \emph{doze vezes} além do maior
ponto medido. Fontes reais aquecem, $r$ cresce e proteções atuam --- os
$\SI{24}{\ampere}$ são um \textbf{limite superior}, útil para dimensionar
proteção, não uma previsão.}
\end{questao}

\begin{questao}[3]{}
Um circuito tem $n=6$ nós, $b=9$ ramos e $2$ fontes de corrente ideais.
\begin{itensq}
\item Determine o número de equações de Kirchhoff necessárias.
\item Explique como as fontes de corrente reduzem o trabalho.
\item Compare com a contagem para análise nodal.
\end{itensq}
\resposta{$5$ LKC $+$ $4$ LKT $=9$; as fontes de corrente fixam duas correntes
de ramo, restando $7$ incógnitas}
\resolucao{(a) LKC: $n-1=5$. LKT: $b-n+1=9-6+1=4$. Total: $9$ --- igual ao
número de correntes de ramo.\\
(b) \textbf{Cada fonte de corrente ideal fixa diretamente a corrente de seu
ramo.} Duas delas eliminam duas incógnitas, restando $7$. Em compensação, a
tensão sobre cada fonte de corrente é desconhecida, o que impede escrever a LKT
de malhas que a atravessem --- é preciso escolher malhas que as contornem, ou
tratar essas tensões como novas incógnitas.\\
\textbf{Resultado líquido:} sistema $7\times7$ em vez de $9\times9$.\\
(c) \textbf{Análise nodal:} $n-1=5$ equações, com as fontes de corrente entrando
\emph{diretamente} no vetor independente, sem nenhuma manipulação. Sistema
$5\times5$.\\
\textbf{Conclusão:} a análise nodal é claramente superior aqui. Ela é, na
verdade, uma forma organizada de aplicar a LKC com a LKT já embutida na escolha
dos potenciais --- e por isso produz o sistema menor sempre que há fontes de
corrente.}
\end{questao}

\begin{questao}[3]{}
Uma fonte tem resistência interna que \emph{cresce} com a corrente, segundo
$r=0{,}2+0{,}1\,I$ (com $r$ em $\si{\ohm}$ e $I$ em $\si{\ampere}$), com
$\EMF=\SI{6}{\volt}$.
\begin{itensq}
\item Determine a corrente para uma carga de $\SI{1}{\ohm}$.
\item Compare com a previsão do modelo linear com $r=\SI{0.2}{\ohm}$.
\end{itensq}
\resposta{$I=\SI{3.80}{\ampere}$ contra $\SI{5}{\ampere}$ previstos}
\resolucao{(a) A LKT dá
\[
6=(0{,}2+0{,}1I)I+1\cdot I
\;\Longrightarrow\;
0{,}1I^{2}+1{,}2I-6=0.
\]
\[
I=\frac{-1{,}2+\sqrt{1{,}44+2{,}4}}{0{,}2}
=\frac{-1{,}2+1{,}960}{0{,}2}=\SI{3.80}{\ampere}.
\]
\emph{(Verificação: $r=0{,}2+0{,}380=\SI{0.58}{\ohm}$ e
$6/(0{,}58+1)=\SI{3.80}{\ampere}$ \checkmark.)}\\
(b) \textbf{Modelo linear:} $I=6/1{,}2=\SI{5}{\ampere}$ --- \textbf{$32\%$
acima} do valor real.\\
\textbf{Duas lições.}\\
\emph{(i)} As \textbf{leis de Kirchhoff continuam válidas} --- elas não exigem
linearidade. O que muda é que a equação resultante deixa de ser linear e passa a
exigir solução de segundo grau.\\
\emph{(ii)} O modelo de $r$ constante é uma \emph{aproximação de pequenos
sinais}. Ele descreve bem a fonte perto do ponto onde $r$ foi medido, e falha
progressivamente longe dele. Baterias reais comportam-se assim, e é por isso
que fabricantes especificam $r$ para uma corrente de ensaio definida.}
\end{questao}

\begin{questao}[3]{}
Uma bateria de $\EMF=\SI{12}{\volt}$ e $r=\SI{0.5}{\ohm}$ deve entregar tensão
terminal de pelo menos $\SI{11.4}{\volt}$.
\begin{itensq}
\item Determine a corrente máxima admissível.
\item Determine a menor carga que pode ser conectada.
\item Determine a potência entregue nessa condição.
\end{itensq}
\resposta{$I\le\SI{1.2}{\ampere}$; $R\ge\SI{9.5}{\ohm}$;
$P=\SI{13.68}{\watt}$}
\resolucao{(a) De $V=\EMF-rI\ge11{,}4$:
\[
0{,}5I\le0{,}6
\;\Longrightarrow\;
I\le\SI{1.2}{\ampere}.
\]
(b) $R=V/I=11{,}4/1{,}2=\SI{9.5}{\ohm}$ --- cargas menores exigiriam mais
corrente e derrubariam a tensão.\\
(c) $P=VI=11{,}4(1{,}2)=\SI{13.68}{\watt}$, com rendimento
$9{,}5/10=95\%$.\\
\textbf{Contraste com a máxima transferência:} a potência máxima ocorreria em
$R=\SI{0.5}{\ohm}$, valendo $\SI{72}{\watt}$ --- mas com tensão terminal de
apenas $\SI{6}{\volt}$, muito abaixo do requisito.\\
\textbf{A restrição de tensão é quase sempre a que manda} em eletrônica, porque
circuitos alimentados simplesmente não funcionam abaixo de sua tensão mínima. A
saída de engenharia é reduzir $r$, não aceitar tensão menor.}
\end{questao}

\begin{questao}[3]{}
Um circuito contém uma fonte \textbf{dependente} de tensão de valor
$3\,I_x$ (em volts, com $I_x$ em ampères) em série com $\SI{2}{\ohm}$, além de
uma fonte independente de $\SI{20}{\volt}$ e um resistor de $\SI{8}{\ohm}$, em
malha única. A corrente de controle $I_x$ é a própria corrente da malha.
\begin{itensq}
\item Escreva a LKT e resolva.
\item Explique por que Kirchhoff continua válido.
\end{itensq}
\resposta{$I=\SI{1.54}{\ampere}$ (fonte dependente em oposição)}
\resolucao{(a) Percorrendo a malha, com a fonte dependente \emph{opondo-se}:
\[
20=2I+8I+3I=13I
\;\Longrightarrow\;
I=\frac{20}{13}=\SI{1.54}{\ampere}.
\]
\emph{(Se a fonte dependente auxiliasse, viria $20=10I-3I=7I$ e
$I=\SI{2.86}{\ampere}$ --- o sentido da fonte dependente altera
completamente o resultado, e por isso a convenção de polaridade precisa estar
explícita no esquema.)}\\
(b) \textbf{As leis de Kirchhoff decorrem de conservação de carga e de
energia}, e nenhuma dessas conservações depende de os elementos serem lineares
ou independentes. A LKT continua sendo "a soma das quedas em uma malha é nula";
a LKC continua sendo "o que entra em um nó, sai".\\
\textbf{O que muda:} a fonte dependente introduz uma \textbf{equação de
restrição} adicional, ligando seu valor a uma variável do circuito. Essa
equação é resolvida \emph{junto} com as de Kirchhoff.\\
\textbf{Consequência estrutural:} o termo migra para o lado das incógnitas, o
que pode tornar a matriz assimétrica --- como se observa na análise nodal e na
de malhas. A reciprocidade se perde; as leis, não.}
\end{questao}

\begin{questao}[3]{}
Uma fonte alimenta uma carga através de um cabo de $\SI{0.3}{\ohm}$ (ida e
volta). A fonte tem $\EMF=\SI{24}{\volt}$ e $r=\SI{0.2}{\ohm}$; a carga vale
$\SI{9.5}{\ohm}$.
\begin{itensq}
\item Determine a corrente, a tensão na carga e a tensão nos terminais da
      fonte.
\item Determine onde a energia é perdida.
\end{itensq}
\resposta{$I=\SI{2.4}{\ampere}$; $V_{carga}=\SI{22.8}{\volt}$;
$V_{fonte}=\SI{23.52}{\volt}$}
\resolucao{(a)
\[
I=\frac{24}{0{,}2+0{,}3+9{,}5}=\frac{24}{10}=\SI{2.4}{\ampere};
\]
\[
V_{carga}=9{,}5(2{,}4)=\SI{22.8}{\volt};
\qquad
V_{fonte}=24-0{,}2(2{,}4)=\SI{23.52}{\volt}.
\]
A diferença de $\SI{0.72}{\volt}$ entre os dois é exatamente a queda no cabo.\\
(b) \textbf{Distribuição das perdas:}
\begin{center}
\begin{tabular}{lrr}
\hline
Elemento & Potência & Fração\\
\hline
Carga ($\SI{9.5}{\ohm}$) & $\SI{54.72}{\watt}$ & $95{,}0\%$\\
Cabo ($\SI{0.3}{\ohm}$) & $\SI{1.73}{\watt}$ & $3{,}0\%$\\
Interna ($\SI{0.2}{\ohm}$) & $\SI{1.15}{\watt}$ & $2{,}0\%$\\
\hline
Total & $\SI{57.6}{\watt}$ & $100\%$\\
\hline
\end{tabular}
\end{center}
\textbf{Observação de diagnóstico:} medindo apenas nos terminais da fonte
($\SI{23.52}{\volt}$), a instalação pareceria melhor do que é. A tensão que
importa é a que chega à \emph{carga} --- e ela está $5\%$ abaixo da f.e.m.\\
\textbf{Regra de medição:} sempre meça no ponto de uso, não na fonte. As duas
leituras diferem exatamente pela queda no caminho.}
\end{questao}

\blocodesafio

\begin{questao}[4]{}
\textbf{(Ensaio de caracterização.)} Descreva como obter $\EMF$ e $r$ de uma
fonte por ensaios em vazio e em curto-circuito.
\begin{itensq}
\item Deduza as relações.
\item Analise os riscos do ensaio em curto.
\item Proponha uma alternativa segura e compare.
\end{itensq}
\resposta{$\EMF=V_{oc}$ e $r=V_{oc}/I_{cc}$; o curto é destrutivo em fontes de
baixa impedância}
\resolucao{(a) \textbf{Em vazio} ($I=0$): $V_{oc}=\EMF$ --- a leitura direta,
desde que o voltímetro tenha impedância muito maior que $r$.\\
\textbf{Em curto} ($V=0$): $\EMF=rI_{cc}$, logo
\[
r=\frac{V_{oc}}{I_{cc}}.
\]
(b) \textbf{Riscos do curto.}
\begin{itemize}
\item Uma pilha AA de $\SI{1.5}{\volt}$ com $r=\SI{0.25}{\ohm}$ entrega
$\SI{6}{\ampere}$ e dissipa $\SI{9}{\watt}$ internamente --- aquecimento
rápido e vazamento.
\item Uma bateria de automóvel ($r\approx\SI{5}{\milli\ohm}$) chega a
$\SI{2500}{\ampere}$, com centelhamento, fusão de terminais e risco de
explosão por ignição do hidrogênio liberado.
\item Baterias de lítio podem entrar em fuga térmica e incendiar.
\item Além disso, $r$ \emph{aumenta} durante o curto, de modo que a leitura
subestima a corrente inicial e superestima $r$.
\end{itemize}
(c) \textbf{Alternativa segura: dois pontos de carga.} Meça $(I_1,V_1)$ e
$(I_2,V_2)$ com resistências conhecidas e moderadas, e ajuste a reta:
\[
r=-\frac{V_2-V_1}{I_2-I_1};
\qquad
\EMF=V_1+rI_1.
\]
\textbf{Comparação.} O método de dois pontos é seguro, não perturba a fonte, e
com três ou mais pontos ainda \emph{testa} a linearidade do modelo --- coisa que
o ensaio de curto não faz. O ensaio em curto só se justifica em fontes
projetadas para suportá-lo, com limitação eletrônica de corrente.}
\end{questao}

\begin{questao}[4]{}
\textbf{(Demonstração --- balanço de potência.)} Prove que a soma algébrica das
potências de todos os elementos de um circuito é nula.
\begin{itensq}
\item Estabeleça a convenção de sinais.
\item Demonstre a partir das leis de Kirchhoff.
\item Explique por que o resultado independe da natureza dos elementos.
\end{itensq}
\resposta{$\sum_k v_ki_k=0$ --- consequência direta de LKC e LKT}
\resolucao{(a) \textbf{Convenção do receptor:} para cada ramo $k$, adote a
corrente $i_k$ entrando pelo terminal de maior potencial. A potência
$p_k=v_ki_k$ é então \emph{positiva} se o elemento absorve e \emph{negativa} se
fornece.\\
(b) \textbf{Demonstração.} Escreva cada tensão de ramo como diferença de
potenciais nodais: $v_k=V_a-V_b$, onde $a$ e $b$ são os nós do ramo $k$. Então
\[
\sum_k p_k=\sum_k (V_a-V_b)\,i_k.
\]
Reagrupando por \emph{nó} em vez de por ramo, cada potencial $V_n$ aparece
multiplicado pela soma algébrica das correntes que chegam ao nó $n$:
\[
\sum_k p_k=\sum_n V_n\left(\sum_{k\in n}\pm i_k\right)=\sum_n V_n\cdot 0=0,
\]
pois o parêntese é exatamente a \textbf{LKC} do nó $n$.\\
\emph{(A escrita $v_k=V_a-V_b$ já embute a \textbf{LKT}: ela pressupõe que o
potencial é bem definido em cada nó.)}\\
(c) \textbf{Independência da natureza dos elementos.} A demonstração usou
\emph{apenas} as duas leis de Kirchhoff e a definição $p=vi$. Em nenhum momento
foi preciso saber a relação entre $v_k$ e $i_k$ de cada elemento.\\
Portanto o resultado vale para resistores, fontes, diodos, capacitores,
indutores, elementos ativos --- lineares ou não, com ou sem memória. É o
\textbf{teorema de Tellegen}, e sua generalidade é notável: ele é uma
propriedade da \emph{topologia} do circuito, não dos componentes.\\
\textbf{Uso prático:} verificar o balanço de potência é o teste de consistência
mais completo de qualquer solução --- se ele não fecha, há erro em algum lugar.}
\end{questao}

\begin{questao}[4]{}
\textbf{(Formulação com fonte ideal entre dois nós.)} Um circuito tem dois nós
não referenciais, $A$ e $B$, com uma fonte ideal de tensão $\EMF$ ligada
diretamente entre eles.
\begin{itensq}
\item Explique por que a LKC não pode ser escrita isoladamente em $A$ nem em
      $B$.
\item Formule o sistema corretamente.
\item Determine a corrente na fonte após resolver.
\end{itensq}
\resposta{Usar supernó: uma LKC envolvendo $A$ e $B$, mais a restrição
$V_A-V_B=\EMF$}
\resolucao{(a) A corrente que atravessa uma fonte de tensão \textbf{ideal} é
determinada pelo circuito, não pela fonte --- não existe relação
$i=f(v)$ para ela. Ao escrever a LKC em $A$, essa corrente desconhecida aparece
como uma incógnita a mais, e o mesmo ocorre em $B$. Duas equações, três
incógnitas.\\
(b) \textbf{Formulação por supernó.} Englobe $A$ e $B$ em uma única superfície
fechada. Pela LKC generalizada, a corrente interna da fonte \emph{se cancela}
--- ela entra em um nó e sai do outro:
\[
\sum_{\text{ramos que cruzam a fronteira}} i=0.
\]
Essa é \emph{uma} equação. A segunda vem da própria fonte:
\[
V_A-V_B=\EMF.
\]
Duas equações, duas incógnitas ($V_A$ e $V_B$). \textbf{Sistema fechado.}\\
(c) \textbf{Corrente na fonte:} obtida \emph{depois}, aplicando a LKC a
\emph{um} dos dois nós isoladamente. Com $V_A$ e $V_B$ já conhecidos, todas as
demais correntes daquele nó são calculáveis, e a da fonte é o que falta para
fechar o balanço.\\
\textbf{Observação:} se um dos terminais da fonte for o próprio nó de
referência, o problema desaparece --- o outro potencial fica \emph{conhecido} e
deixa de ser incógnita. \textbf{Escolher o terra em um terminal de fonte é a
simplificação mais barata que existe.}}
\end{questao}

\begin{questao}[4]{}
\textbf{(Projeto de verificação experimental.)} Elabore um ensaio de laboratório
para verificar a LKC em um nó com três ramos, considerando incertezas dos
instrumentos.
\begin{itensq}
\item Descreva o procedimento.
\item Propague as incertezas e estabeleça o critério de aceitação.
\item Interprete um resíduo de $\SI{0.35}{\ampere}$ em correntes de alguns
      ampères.
\end{itensq}
\resposta{Critério: resíduo compatível com $u_c\approx\sqrt{\sum u_k^{2}}$;
$\SI{0.35}{\ampere}$ indica erro real}
\resolucao{(a) \textbf{Procedimento.} Meça as três correntes
\emph{simultaneamente}, com três amperímetros --- ou, se houver apenas um,
faça as três medições sem alterar o circuito entre elas. Registre sentidos e
adote uma convenção única (por exemplo, positivo entrando).\\
\textbf{Cuidado essencial:} inserir o amperímetro \emph{altera} o circuito, pois
acrescenta sua resistência interna ao ramo. Com três inserções sucessivas, o
circuito muda três vezes e a verificação perde sentido. Use amperímetro de baixa
resistência ou, melhor, alicate amperímetro, que não abre o ramo.\\
(b) \textbf{Propagação.} Suponha
$I_1=(5{,}02\pm0{,}05)\,\si{\ampere}$,
$I_2=(3{,}01\pm0{,}03)\,\si{\ampere}$ entrando, e $I_3$ saindo. A previsão é
$I_1+I_2=\SI{8.03}{\ampere}$, com incerteza combinada
\[
u_c=\sqrt{0{,}05^{2}+0{,}03^{2}}=\SI{0.058}{\ampere}.
\]
Incluindo a incerteza da própria medição de $I_3$ (digamos
$\SI{0.08}{\ampere}$), o resíduo esperado tem incerteza
$\sqrt{0{,}058^{2}+0{,}08^{2}}=\SI{0.099}{\ampere}$.\\
\textbf{Critério de aceitação:} $|I_1+I_2-I_3|\le2u\approx
\SI{0.20}{\ampere}$ (dois desvios, cerca de $95\%$ de confiança).\\
(c) \textbf{Resíduo de $\SI{0.35}{\ampere}$} é $3{,}5$ vezes a incerteza
combinada --- \textbf{não é compatível com erro aleatório}. Causas prováveis, em
ordem:
\begin{itemize}
\item \textbf{Erro de sinal ou de sentido} em uma das leituras --- o mais comum.
\item \textbf{Um quarto ramo não contabilizado}, inclusive uma fuga para o terra
ou pela bancada.
\item \textbf{Instrumento descalibrado} ou operando fora de escala.
\item \textbf{Circuito alterado entre medições} (aquecimento, contato
intermitente).
\end{itemize}
\textbf{Conclusão importante:} a LKC \emph{não} pode ser violada --- ela é
conservação de carga. Um resíduo significativo é sempre erro de medição ou de
modelo, nunca da lei. O ensaio, portanto, verifica o \emph{procedimento}, e não
a física.}
\end{questao}

\begin{questao}[4]{}
\textbf{(Baterias em série desequilibradas.)} Quatro células de
$\SI{1.5}{\volt}$ estão em série alimentando uma carga, mas uma delas está
gasta, com $\EMF=\SI{0.9}{\volt}$ e $r=\SI{1.5}{\ohm}$ (as demais mantêm
$r=\SI{0.2}{\ohm}$). A carga vale $\SI{5}{\ohm}$.
\begin{itensq}
\item Determine a corrente e a tensão em cada célula.
\item Identifique o que ocorre com a célula gasta.
\item Discuta a implicação prática.
\end{itensq}
\resposta{$I=\SI{0.76}{\ampere}$; a célula gasta apresenta
$\SI{-0.24}{\volt}$ --- está sendo \textbf{invertida}}
\resolucao{(a)
\[
\EMF_{total}=3(1{,}5)+0{,}9=\SI{5.4}{\volt};
\qquad
r_{total}=3(0{,}2)+1{,}5=\SI{2.1}{\ohm};
\]
\[
I=\frac{5{,}4}{2{,}1+5}=\frac{5{,}4}{7{,}1}=\SI{0.76}{\ampere}.
\]
\textbf{Tensão de cada célula boa:} $1{,}5-0{,}2(0{,}76)=\SI{1.35}{\volt}$.\\
\textbf{Tensão da célula gasta:}
$0{,}9-1{,}5(0{,}76)=0{,}9-1{,}14=\SI{-0.24}{\volt}$.\\
(b) \textbf{A tensão da célula gasta é negativa} --- ou seja, ela está sendo
percorrida por corrente no sentido de \emph{carga}, forçada pelas três células
boas. Ela deixou de ser fonte e passou a ser carga, dissipando
$1{,}5(0{,}76)^{2}=\SI{0.87}{\watt}$ internamente.\\
\textbf{Em células não recarregáveis}, isso provoca a chamada \textbf{inversão
de polaridade}: reações internas indesejadas, geração de gás, aquecimento e
vazamento de eletrólito --- exatamente o que se observa quando uma pilha
"vaza" em um controle remoto esquecido ligado.\\
(c) \textbf{Implicações práticas.}
\begin{itemize}
\item \textbf{Nunca misture pilhas novas e usadas}, nem de marcas ou tecnologias
diferentes, em série. A mais fraca será invertida pelas demais.
\item \textbf{Substitua sempre o conjunto completo.}
\item Em baterias de lítio em série, o problema é grave a ponto de exigir
\textbf{circuito de balanceamento} (BMS), que monitora cada célula
individualmente e impede que qualquer uma seja levada à inversão.
\end{itemize}
\textbf{Diagnóstico:} a célula gasta domina o conjunto por sua \emph{resistência
interna}, não por sua f.e.m. --- $\SI{1.5}{\ohm}$ contra
$\SI{0.6}{\ohm}$ das outras três somadas.}
\end{questao}

\begin{questao}[4]{}
\textbf{(Modelagem --- fonte com limitação de corrente.)} Uma fonte de bancada
tem $\EMF=\SI{20}{\volt}$ e $r=\SI{0.1}{\ohm}$, com limitação eletrônica em
$\SI{2}{\ampere}$.
\begin{itensq}
\item Descreva a característica $V\times I$ completa.
\item Determine o ponto de operação para cargas de $\SI{50}{\ohm}$,
      $\SI{10}{\ohm}$ e $\SI{2}{\ohm}$.
\item Explique por que o modelo de fonte real deixa de valer.
\end{itensq}
\resposta{Duas regiões: fonte de tensão até $\SI{2}{\ampere}$, fonte de corrente
depois}
\resolucao{(a) \textbf{Característica em dois trechos.}
\begin{itemize}
\item \textbf{Região de tensão constante} ($I<\SI{2}{\ampere}$):
$V=20-0{,}1I$ --- reta quase horizontal.
\item \textbf{Região de corrente constante} ($V<\SI{19.8}{\volt}$):
$I=\SI{2}{\ampere}$ --- reta vertical.
\end{itemize}
O "joelho" fica em $(\SI{2}{\ampere};\ \SI{19.8}{\volt})$, com
$R_{joelho}=\SI{9.9}{\ohm}$.\\
(b)
\begin{center}
\begin{tabular}{lccl}
\hline
$R_L$ & $I$ & $V$ & Região\\
\hline
$\SI{50}{\ohm}$ & $\SI{0.399}{\ampere}$ & $\SI{19.96}{\volt}$ & tensão\\
$\SI{10}{\ohm}$ & $\SI{1.98}{\ampere}$ & $\SI{19.80}{\volt}$ & tensão (limite)\\
$\SI{2}{\ohm}$ & $\SI{2.00}{\ampere}$ & $\SI{4.00}{\volt}$ & corrente\\
\hline
\end{tabular}
\end{center}
Note o terceiro caso: o modelo linear preveria
$20/2{,}1=\SI{9.5}{\ampere}$ --- quase cinco vezes o real.\\
(c) \textbf{Por que o modelo falha.} A fonte real é um modelo \emph{linear} de
dois parâmetros, e sua característica é uma \emph{única} reta. A limitação
eletrônica introduz um segundo trecho, tornando a característica
\textbf{quebrada} --- e nenhuma reta descreve as duas regiões.\\
\textbf{Como proceder corretamente:}
\begin{itemize}
\item Determinar em qual região a carga opera, comparando $R_L$ com
$R_{joelho}=\SI{9.9}{\ohm}$.
\item $R_L>R_{joelho}$: use o modelo de \emph{fonte de tensão} $(\EMF,r)$.
\item $R_L<R_{joelho}$: use o modelo de \emph{fonte de corrente} ideal de
$\SI{2}{\ampere}$, e a tensão sai de $V=R_LI$.
\end{itemize}
\textbf{As leis de Kirchhoff seguem válidas nas duas regiões} --- o que muda é
apenas a relação constitutiva da fonte. É o mesmo procedimento usado com
diodos: identificar o trecho e aplicar o modelo correspondente.}
\end{questao}

\begin{questao}[4]{}
\textbf{(Sistema completo por Kirchhoff.)} Um circuito de duas malhas tem
$\EMF_1=\SI{18}{\volt}$ com $r_1=\SI{1}{\ohm}$ no primeiro ramo,
$\EMF_2=\SI{12}{\volt}$ com $r_2=\SI{2}{\ohm}$ no segundo, e
$R=\SI{6}{\ohm}$ no ramo comum.
\begin{itensq}
\item Escreva as equações de Kirchhoff e resolva.
\item Verifique o balanço de potência.
\item Identifique o papel de cada fonte.
\end{itensq}
\resposta{$I_1=\SI{3.6}{\ampere}$, $I_2=\SI{-1.2}{\ampere}$,
$I_R=\SI{2.4}{\ampere}$}
\resolucao{(a) Chamando $V$ a tensão do nó comum:
\[
\frac{V-18}{1}+\frac{V-12}{2}+\frac{V}{6}=0.
\]
Multiplicando por $6$: $6V-108+3V-36+V=0\Rightarrow10V=144
\Rightarrow V=\SI{14.4}{\volt}$.
\[
I_1=\frac{18-14{,}4}{1}=\SI{3.6}{\ampere};
\quad
I_2=\frac{12-14{,}4}{2}=\SI{-1.2}{\ampere};
\quad
I_R=\frac{14{,}4}{6}=\SI{2.4}{\ampere}.
\]
\textbf{LKC:} $3{,}6-1{,}2=2{,}4$ \checkmark.\\
(b) \textbf{Balanço.}
\begin{center}
\begin{tabular}{lr}
\hline
$\EMF_1$ fornece $18(3{,}6)$ & $\SI{64.8}{\watt}$\\
$r_1$ dissipa $1(3{,}6)^{2}$ & $\SI{12.96}{\watt}$\\
$R$ dissipa $6(2{,}4)^{2}$ & $\SI{34.56}{\watt}$\\
$\EMF_2$ absorve $12(1{,}2)$ & $\SI{14.4}{\watt}$\\
$r_2$ dissipa $2(1{,}2)^{2}$ & $\SI{2.88}{\watt}$\\
\hline
Total consumido & $\SI{64.8}{\watt}$\\
\hline
\end{tabular}
\end{center}
Fecha exatamente \checkmark.\\
(c) \textbf{Papéis.} A fonte 1 é a única \emph{geradora}: ela alimenta a carga
$R$ \emph{e} carrega a fonte 2. A fonte 2 comporta-se como \textbf{receptor},
porque sua f.e.m. ($\SI{12}{\volt}$) é menor que a tensão do nó
($\SI{14.4}{\volt}$) --- e corrente sempre entra por onde o potencial externo é
maior.\\
\textbf{Critério geral:} compare a f.e.m. de cada fonte com a tensão do nó a
que está ligada. Maior $\Rightarrow$ fornece; menor $\Rightarrow$ recebe. O
sinal da corrente calculada confirma automaticamente.}
\end{questao}

% END BANCO COMPLEMENTAR

\gabarito[2.4 Fontes reais e Leis de Kirchhoff]
